% Generated mechanically from frozen input p07/ja/pione.tex.
% Do not edit this generated file; edit the bound locale source instead.

% OK, start here.
%

\setcounter{chapter}{57}
\chapter{スキームの基本群}



\phantomsection
\label{pione-section-phantom}


\section{序論}
\label{pione-section-introduction}

\noindent
本章では Grothendieck によるスキームの基本群とその応用を論じる。
基礎的文献は \cite{SGA1} である。
よい入門として \cite{Lenstra} がある。
その他の参考文献として \cite{Murre-lectures} および
\cite{Grothendieck-Murre} がある。










\section{点上エタールなスキーム}
\label{pione-section-schemes-etale-point}

\noindent
本節では、体のスペクトル上でエタールなスキームを記述する。
結果を述べる前に、位相群 $G$ に対する $G$-集合の圏を導入する。

\begin{definition}
\label{pione-definition-G-set-continuous}
位相群 $G$ をとる。
{\it $G$-集合}（{\it 離散 $G$-集合}ともいう）とは、集合 $X$ であって、
左作用 $a : G \times X \to X$ を備え、$a$ が連続となるものをいう。
ただし $X$ には離散位相を、$G \times X$ には積位相を入れる。
{\it $G$-集合の射} $f : X \to Y$ とは、単に任意の $G$-同変写像、
すなわち $X$ から $Y$ への写像のことである。
$G$-集合の圏は {\it $G\textit{-Sets}$} と表される。
\end{definition}

\noindent
$a : G \times X \to X$ が連続であるという条件は、単に任意の
$x \in X$ の安定化群が $G$ の開部分群であることを意味する。
% P07-ERR-1144: frozen-source duplicate group symbol retained; see sidecar.
$G$ が抽象群 $G$（すなわち位相群ではない群）である場合、$G$ に
離散位相を入れれば、これは先の定義と一致する（例えば
Sites, 例 \ref{sites-example-site-on-group} を参照せよ）。

\medskip\noindent
$L/K$ が無限 Galois 拡大ならば、Galois 群
$G = \text{Gal}(L/K)$ には標準的な位相が備わることを思い出そう。
Fields, 節 \ref{fields-section-infinite-galois} を参照せよ。

\begin{lemma}
\label{pione-lemma-sheaves-point}
体 $K$ をとる。$K^{sep}$ を $K$ の分離閉包とする。
副有限群 $G = \text{Gal}(K^{sep}/K)$ を考える。関手
$$
\begin{matrix}
\text{スキーム（}K\text{ 上でエタール）} &
\longrightarrow &
G\textit{-Sets} \\
X/K & \longmapsto &
\Mor_{\Spec(K)}(\Spec(K^{sep}), X)
\end{matrix}
$$
は圏同値である。
\end{lemma}

\begin{proof}
スキーム $X$ が $K$ 上にあり、$K$ 上エタールであるための必要十分条件は、
$X \cong \coprod_{i\in I} \Spec(K_i)$ と書けて、各 $K_i$ が $K$ の
有限分離拡大となることである
（Morphisms, 補題 \ref{morphisms-lemma-etale-over-field}）。
補題の関手は $X$ に $G$-集合
$$
\coprod\nolimits_i \Hom_K(K_i, K^{sep})
$$
を、その自然な左 $G$-作用とともに対応させる。$G$ の位相の定義により、
各元の安定化群は開である。逆に、任意の $G$-集合 $S$ はその軌道の
非交和である。$S = \coprod S_i$ と書く。$s_i \in S_i$ を選び、その
開安定化群を $G_i \subset G$ と表す。Galois 理論
（Fields, 定理 \ref{fields-theorem-inifinite-galois-theory}）により、体
$(K^{sep})^{G_i}$ は $K$ の有限分離拡大である。したがってスキーム
$$
\coprod\nolimits_i \Spec((K^{sep})^{G_i})
$$
は $K$ 上エタールである。これにより補題の関手の逆関手が得られる。
細部の一部は省略する。
\end{proof}

\begin{remark}
\label{pione-remark-covering-surjective}
補題 \ref{pione-lemma-sheaves-point} の対応のもとで、小エタール・サイト
$\Spec(K)_\etale$（体 $K$ のもの）における被覆は、
$G\textit{-Sets}$ における射の
全射族に対応する。
\end{remark}









\section{ガロア圏}
\label{pione-section-galois}

\noindent
この節では、\cite[Expos\'e V, Sections 4, 5, and 6]{SGA1} に収められた
内容の一部を論じる。

\medskip\noindent
$F : \mathcal{C} \to \textit{Sets}$ を関手とする。
本書の規約では、圏の対象は集合をなし、任意の二対象の間の射も集合をなすことを
思い出そう。標準的な単射
\begin{equation}
\label{pione-equation-embedding-product}
\text{Aut}(F)
\longrightarrow
\prod\nolimits_{X \in \Ob(\mathcal{C})} \text{Aut}(F(X))
\end{equation}
がある。集合 $E$ に対して $\text{Aut}(E)$ にはコンパクト開位相を入れる
（Topology, 例 \ref{topology-example-automorphisms-of-a-set} を参照）。
$E$ が有限なら、これはもちろん離散位相であり、この節で関心をもつのは
この場合である\footnote{pro-\'etale 基本群を論じる際には、一般の場合が重要になる。}。
$\text{Aut}(F)$ には、(\ref{pione-equation-embedding-product}) の右辺の積位相から
誘導される位相を入れる。とくに、作用写像
$$
\text{Aut}(F) \times F(X) \longrightarrow F(X)
$$
は、$F(X)$ に離散位相を入れると連続である。実際、作用写像
$\text{Aut}(E) \times E \to E$ は任意の集合 $E$ に対して連続だからである。
$\text{Aut}(F)$ に入れた位相の普遍性は次のとおりである。$G$ を位相群とし、
$G \to \text{Aut}(F)$ を群準同型であって、誘導される作用
$G \times F(X) \to F(X)$ がすべての $X \in \Ob(\mathcal{C})$ に対し、
$F(X)$ に離散位相を入れたとき連続になるものとする。このとき
$G \to \text{Aut}(F)$ は連続である。

\medskip\noindent
次の補題は、有限集合の圏への関手の自己同型群が自動的に副有限群になることを
示す。

\begin{lemma}
\label{pione-lemma-aut-inverse-limit}
$\mathcal{C}$ を圏、$F : \mathcal{C} \to \textit{Sets}$ を関手とする。
写像 (\ref{pione-equation-embedding-product}) により $\text{Aut}(F)$ は
$\prod_{X \in \Ob(\mathcal{C})} \text{Aut}(F(X))$ の閉部分群と同一視される。
とくに、$F(X)$ がすべての $X$ に対して有限なら、$\text{Aut}(F)$ は
副有限群である。
\end{lemma}

\begin{proof}
$\xi = (\gamma_X) \in \prod \text{Aut}(F(X))$ を
$\text{Aut}(F)$ に属さない元とする。このとき、射 $f : X \to X'$ であって
$\mathcal{C}$ に属するものと、元 $x \in F(X)$ であって
$F(f)(\gamma_X(x)) \not = \gamma_{X'}(F(f)(x))$.
を満たすものが存在する。開近傍
$U = \{\gamma \in \text{Aut}(F(X)) \mid \gamma(x) = \gamma_X(x)\}$
（これは $\gamma_X$ の近傍である）と、開近傍
$U' = \{\gamma' \in \text{Aut}(F(X')) \mid \gamma'(F(f)(x)) =
\gamma_{X'}(F(f)(x))\}$.
を考える。このとき
$U \times U' \times \prod_{X'' \not = X, X'} \text{Aut}(F(X''))$
は $\xi$ の開近傍であり、$\text{Aut}(F)$ と交わらない。
最後の主張は、$\prod \text{Aut}(F(X))$ が、各 $F(X)$ が有限なら
副有限空間であることから従う。
\end{proof}

\begin{example}
\label{pione-example-galois-category-G-sets}
$G$ を位相群とする。重要な例は忘却関手
\begin{equation}
\label{pione-equation-forgetful}
\textit{Finite-}G\textit{-Sets} \longrightarrow \textit{Sets}
\end{equation}
である。ここで $\textit{Finite-}G\textit{-Sets}$ は、有限 $G$-集合を対象とする
$G\textit{-Sets}$ の充満部分圏である。$G$-集合の圏 $G\textit{-Sets}$ は
定義 \ref{pione-definition-G-set-continuous} で定義した。
\end{example}

\noindent
$G$ を位相群とする。$G$ の {\it 副有限完備化}とは、副有限位相を備えた
副有限群
$$
G^\wedge =
\lim_{U \subset G\text{ open, normal, finite index}} G/U
$$
のことである。有限指数の開正規部分群の有限個の共通部分も同じ性質をもつので、
この極限は余フィルター型である。副有限完備化の普遍性とは、任意の連続写像
$G \to H$ で、その終域 $H$ が副有限群であるものが、標準的に
$G \to G^\wedge \to H$ と分解することである。

\begin{lemma}
\label{pione-lemma-single-out-profinite}
$G$ を位相群とする。関手 (\ref{pione-equation-forgetful}) の自己同型群に補題
\ref{pione-lemma-aut-inverse-limit} の副有限位相を入れたものは、$G$ の副有限完備化である。
\end{lemma}

\begin{proof}
(\ref{pione-equation-forgetful}) の関手を $F_G$ と書く。任意の射 $X \to Y$ で
$\textit{Finite-}G\textit{-Sets}$ に属するものは $G$ の作用と可換する。
したがって、任意の $g \in G$ は $F_G$ の自己同型を定め、群の標準的な準同型
$G \to \text{Aut}(F_G)$ を得る。任意の有限 $G$-集合 $X$ は、標準的な $G$-作用を
もつ形 $G/H_i$ の $G$-集合の有限非交和である。ここで $H_i \subset G$ は
有限指数の開部分群である。このとき $U_i = \bigcap gH_ig^{-1}$ は開かつ正規で、
有限指数をもつ。さらに $U_i$ は $G/H_i$ に自明に作用するので、
$U = \bigcap U_i$ は $F_G(X)$ に自明に作用する。したがって作用
$G \times F_G(X) \to F_G(X)$ は連続である。$\text{Aut}(F_G)$ の位相の
普遍性により、写像 $G \to \text{Aut}(F_G)$ は連続である。
補題 \ref{pione-lemma-aut-inverse-limit} と副有限完備化の普遍性により、誘導された
連続群準同型
$$
G^\wedge \longrightarrow \text{Aut}(F_G)
$$
が存在する。また、$G/U$ は $G/U$ に忠実に作用するので、この写像は単射である。
その像が稠密なら写像は全射であり、したがって Topology, 補題
\ref{topology-lemma-bijective-map} により同相写像となる。

\medskip\noindent
$\gamma \in \text{Aut}(F_G)$ および $X \in \Ob(\mathcal{C})$ をとる。
$g \in G$ であって、$\gamma$ と $g$ が $F_G(X)$ 上に同じ作用を誘導するものが
存在することを示せば、証明は完了する。前と同様に、$X$ は $G/H_i$ の有限非交和である。
上と同じ $U_i$ および $U$ を用いると、有限 $G$-集合 $Y = G/U$ はすべての
$G/H_i$ 上へすべての $i$ に対して全射となる。したがって、$g \in G$ であって、
$\gamma$ と $g$ が $F_G(G/U) = G/U$ 上に同じ作用を誘導するものを見つければ十分である。
$e \in G$ を単位元とし、$\gamma(eU) = g_0U$ と書ける $g_0 \in G$ をとる。
任意の $g_1 \in G$ に対し、射
$$
R_{g_1} : G/U \longrightarrow G/U,\quad gU \longmapsto gg_1U
$$
は $\textit{Finite-}G\textit{-Sets}$ の射で、$\gamma$ の作用と可換する。ゆえに
$$
\gamma(g_1U) = \gamma(R_{g_1}(eU)) = R_{g_1}(\gamma(eU)) =
R_{g_1}(g_0U) = g_0g_1U
$$
である。したがって $g = g_0$ とすればよい。
\end{proof}

\noindent
完全関手とは、すべての有限極限および有限余極限と可換する関手であった。
とくに、そのような関手は等化子、余等化子、ファイバー積、押し出しなどと可換する。

\begin{lemma}
\label{pione-lemma-second-fundamental-functor}
$G$ を位相群とする。
$F : \textit{Finite-}G\textit{-Sets} \to \textit{Sets}$
を完全関手とし、$F(X)$ はすべての $X$ に対して有限であるとする。
このとき $F$ は関手 (\ref{pione-equation-forgetful}) と同型である。
\end{lemma}

\begin{proof}
$X$ を $\textit{Finite-}G\textit{-Sets}$ の空でない対象とする。図式
$$
\xymatrix{
X \ar[r] \ar[d] & \{*\} \ar[d] \\
\{*\} \ar[r] & \{*\}
}
$$
は余カルテシアンである。したがって $F(X)$ は空でない。
$U \subset G$ を有限指数の開正規部分群とする。等式
$$
G/U \times G/U = \coprod\nolimits_{gU \in G/U} G/U
$$
において、$gU$ に対応する成分は左辺の $(eU, gU)$ の軌道に対応する。したがって
$$
F(G/U) \times F(G/U) = F(G/U \times G/U) = \coprod\nolimits_{gU \in G/U} F(G/U)
$$
$|F(G/U)| = |G/U|$ である。実際、$F(G/U)$ は空でない。ゆえに
$$
% P07-ERR-1145: frozen-source spelling retained; see sidecar.
\lim_{U \subset G\text{ open, normal, finite idex}} F(G/U)
$$
は空でない（Categories, 補題 \ref{categories-lemma-nonempty-limit}）。
この極限の元 $\gamma = (\gamma_U)$ を一つとる。
関手 (\ref{pione-equation-forgetful}) を $F_G$ と書く。$F_G$ は関手
$$
X \longmapsto \colim_U \Mor(G/U, X)
$$
と同一視できる。ここで $f : G/U \to X$ は $f(eU) \in X = F_G(X)$ に対応する
（詳細は省略する）。したがって元 $\gamma$ は、整定義された写像
$$
t : F_G \longrightarrow F
$$
を定める。実際、$x \in X$ が与えられたとき、$U$ と $f : G/U \to X$ を選び、
$eU$ を $x$ に写すものとして、$t_X(x) = F(f)(\gamma_U)$ と置く。
$t$ が全単射 $t_{G/U} : F_G(G/U) \to F(G/U)$ を任意の $U$ に対して
誘導することを示そう。
すると、$t$ が同型であることが直ちに従う（詳細は省略する）。
$|F_G(G/U)| = |F(G/U)|$ であるから、$t_{G/U}$ が全射であることを示せば十分である。
その像は少なくとも元
$t_{G/U}(eU) = F(\text{id}_{G/U})(\gamma_U) = \gamma_U$.
を含む。$g \in G$ に対し、右乗法 $R_g : G/U \to G/U$ と書く。
$F(R_g) : F(G/U) \to F(G/U)$ の不動点集合は $F(\emptyset) = \emptyset$ に等しい。
これは $g \not \in U$ なら $F$ が等化子と可換するためである。
したがって、$g_1, \ldots, g_{|G/U|}$ が $G/U$ の代表元系なら、元
$F(R_{g_i})(\gamma_U)$ は互いに異なり、したがって $F(G/U)$ を尽くす。このとき
$$
t_{G/U}(g_iU) = F(R_{g_i})(\gamma_U)
$$
であり、証明が完了する。
\end{proof}

\begin{example}
\label{pione-example-from-C-F-to-G-sets}
$\mathcal{C}$ を圏、$F : \mathcal{C} \to \textit{Sets}$ を関手とし、
$F(X)$ はすべての $X \in \Ob(\mathcal{C})$ に対して有限であるとする。
補題 \ref{pione-lemma-aut-inverse-limit} により、$G = \text{Aut}(F)$ には標準的に
副有限位相群の構造が備わる。関手
\begin{equation}
\label{pione-equation-remember}
\mathcal{C} \longrightarrow \textit{Finite-}G\textit{-Sets},\quad
X \longmapsto F(X)
\end{equation}
を得る。ここで $F(X)$ には $G$ の誘導作用を入れる。この作用は、
$\text{Aut}(F)$ 上の位相の構成により連続である。
\end{example}

\noindent
ガロア圏を定義する目的は、関手 (\ref{pione-equation-remember}) が同値となるような
対 $(\mathcal{C}, F)$ を特徴づけることにある。ガロア圏を次のように定義する。

\begin{definition}
\label{pione-definition-galois-category}
\begin{reference}
\cite[Expos\'e V, Definition 5.1]{SGA1} の定義とは異なる。
\cite[Definition 7.2.1]{BS} と比較せよ。
\end{reference}
$\mathcal{C}$ を圏、$F : \mathcal{C} \to \textit{Sets}$ を関手とする。
対 $(\mathcal{C}, F)$ が {\it ガロア圏} であるとは、次の条件を満たすことをいう。
\begin{enumerate}
\item $\mathcal{C}$ は有限極限および有限余極限をもつ。
\item
\label{pione-item-connected-components}
$\mathcal{C}$ の各対象は、連結対象の有限余積（空でもよい）である。
\item $F(X)$ はすべての $X \in \Ob(\mathcal{C})$ に対して有限である。
\item $F$ は同型を反映し\footnote{すなわち、射 $f$ が $\mathcal{C}$ に属し、
$F(f)$ が同型なら、$f$ も同型である。}、かつ完全である\footnote{これは $F$ が
有限極限および有限余極限と可換することを意味する。Categories, 節
\ref{categories-section-exact-functor} を参照。}。
\end{enumerate}
ここで、$X \in \Ob(\mathcal{C})$ が連結であるとは、それが始対象でなく、
任意の単射 $Y \to X$ に対して、$Y$ が始対象であるか、または
$Y \to X$ が同型であることをいう。
\end{definition}

\noindent
{\bf 注意：} この定義は、多くの文献にある定義と同一ではない（ただし後に
同値であることが分かる）。実際、\cite[Expos\'e V, Definition 5.1]{SGA1} では、
ガロア圏を、ある副有限群 $G$ に対する $\textit{Finite-}G\textit{-Sets}$ と
同値な圏として定義する。次いで Grothendieck は、上の公理より弱い公理
(G1)--(G6) の列によってガロア圏を特徴づける。本書の定義を選ぶ理由は、
有限極限と有限余極限の存在、および関手 $F$ の完全性を強調するためである。
その代償として、後にこの節の結果を適用する際には、やや多くの作業が必要になる。

\begin{lemma}
\label{pione-lemma-epi-mono}
$(\mathcal{C}, F)$ をガロア圏とし、
$X \to Y \in \text{Arrows}(\mathcal{C})$ とする。このとき次が成り立つ。
\begin{enumerate}
\item $F$ は忠実である。
\item $X \to Y$ が単射であることと、
$\Leftrightarrow F(X) \to F(Y)$ が単射写像であることは同値である。
\item $X \to Y$ が全射であることと、
$\Leftrightarrow F(X) \to F(Y)$ が全射写像であることは同値である。
\item 対象 $A$ が $\mathcal{C}$ の始対象であるための必要十分条件は
$F(A) = \emptyset$ である。
\item 対象 $Z$ が $\mathcal{C}$ の終対象であるための必要十分条件は
$F(Z)$ が一元集合であることである。
\item $X$ と $Y$ が連結なら、$X \to Y$ は全射である。
\item
\label{pione-item-one-element}
$X$ が連結で、$a, b : X \to Y$ が二つの射であるとする。
$a = b$ が成り立つには、$F(a)$ と $F(b)$ が $F(X)$ の一つの元上で一致すれば十分である。
\item $X = \coprod_{i = 1, \ldots, n} X_i$ および
$Y = \coprod_{j = 1, \ldots, m} Y_j$ で、$X_i$, $Y_j$ が連結であるとする。
このとき写像 $\alpha : \{1, \ldots, n\} \to \{1, \ldots, m\}$ が存在し、
$X \to Y$ は射の族 $X_i \to Y_{\alpha(i)}$ から得られる。
\end{enumerate}
\end{lemma}

\begin{proof}
（1）の証明。$a, b : X \to Y$ かつ $F(a) = F(b)$ とする。
$E$ を $a$ と $b$ の等化子とする。このとき $F(E) = F(X)$ であり、
$E = X$ である。これは $F$ が同型を反映するためである。

\medskip\noindent
（2）の証明。$F$ は射 $X \to X \times_Y X$ を写像
$F(X) \to F(X) \times_{F(Y)} F(X)$ に写し、$F$ は同型を反映するからである。

\medskip\noindent
（3）の証明。$F$ は射 $Y \amalg_X Y \to Y$ を写像
$F(Y) \amalg_{F(X)} F(Y) \to F(Y)$ に写し、$F$ は同型を反映するからである。

\medskip\noindent
（4）の証明。始対象 $A$ は存在し、明らかに $F(A) = \emptyset$ である。
一方、対象 $X$ が $F(X) = \emptyset$ を満たすなら、一意な写像 $A \to X$ は
全単射 $F(A) \to F(X)$ を誘導するから、$A \to X$ は同型である。

\medskip\noindent
（5）の証明。終対象 $Z$ は存在し、明らかに $F(Z)$ は一元集合である。
一方、対象 $X$ に対して $F(X)$ が一元集合なら、一意な写像 $X \to Z$ は
全単射 $F(X) \to F(Z)$ を誘導するから、$X \to Z$ は同型である。

\medskip\noindent
（6）の証明。等化子 $E$、すなわち二つの写像 $Y \to Y \amalg_X Y$ の等化子は、
$\mathcal{C}$ の始対象ではない。実際、$X \to Y$ は $E$ を経由して分解し、
$F(X) \not = \emptyset$ である。ゆえに $E = Y$ であり、主張が従う。

\medskip\noindent
（\ref{pione-item-one-element}）の証明。
等化子 $E$、すなわち $a$ と $b$ の等化子には単射 $E \to X$ が備わり、
$F(E) \subset F(X)$ は $F(a)$ と $F(b)$ が一致する元全体の集合である。
$E$ が始対象であるか $E = X$ であるかのいずれかなので、主張が従う。

\medskip\noindent
（8）の証明。各 $i, j$ に対して、$E_{ij} = X_i \times_Y Y_j$ は始対象であるか、
$X_i$ に等しい。$s \in F(X_i)$ を一つ選ぶと、$E_{ij} = X_i$ であるための
必要十分条件は、$s$ が $F(Y_j) \subset F(Y)$ の元に写ることである。
したがって、これは一意な $j = \alpha(i)$ に対して起こる。
\end{proof}

\noindent
上の補題より、連結対象 $X$ がガロア圏 $(\mathcal{C}, F)$ に与えられたとき、
自己同型群 $\text{Aut}(X)$ の位数は高々 $|F(X)|$ である。実際、
$s \in F(X)$ および $g \in \text{Aut}(X)$ に対して、
$F(g)(s) = s$ であることと $g = \text{id}_X$ であることは
（\ref{pione-item-one-element}）により同値である。等号が成り立つとき、$X$ は
{\it ガロア} であるという。同値な言い方をすれば、$X$ がガロアであるとは、
それが連結であり、$\text{Aut}(X)$ が $F(X)$ に推移的に作用することである。

\begin{lemma}
\label{pione-lemma-galois}
$(\mathcal{C}, F)$ をガロア圏とする。任意の連結対象 $X$ で $\mathcal{C}$ に属するものに対して、
ガロア対象 $Y$ と射 $Y \to X$ が存在する。
\end{lemma}

\begin{proof}
以下、補題 \ref{pione-lemma-epi-mono} の結果を断りなく用いる。
$n = |F(X)|$ とする。$X^n$ に $S_n$ の自然な作用を入れたものを考える。
$$
X^n = \coprod\nolimits_{t \in T} Z_t
$$
を連結対象への分解とする。ある $t$ を選び、$F(Z_t)$ が元
$(s_1, \ldots, s_n)$ を含み、そこで $s_i$ は互いに異なるものとする。別の元
$(s'_1, \ldots, s'_n) \in F(Z_t)$ に対しても、$s'_i$ は互いに異なると主張する。
% P07-ERR-1146: frozen-source article error recorded in sidecar.
そうでなく、たとえば $s'_i = s'_j$ なら、$Z_t$ は $X^{n - 1}$ のある連結成分の、
対角射
$$
\Delta_{ij} : X^{n - 1} \longrightarrow X^n
$$
による像である。連結対象の射は全射であり、$F$ を施すと全射写像を誘導するので、
$s_i = s_j$ が従うが、これは成り立たない。

\medskip\noindent
$G \subset S_n$ を、$g(Z_t) = Z_t$ を満たす元からなる部分群とする。
$S_n$ の
$$
F(X)^n = F(X^n) = \coprod\nolimits_{t' \in T} F(Z_{t'})
$$
への作用を見ると、$G = \{g \in S_n \mid g(s_1, \ldots, s_n) \in F(Z_t)\}$ である。
ここで第二の元 $(s'_1, \ldots, s'_n) \in F(Z_t)$ を選ぶ。上で $s'_i$ が互いに
異なることを示した。したがって $g \in S_n$ であって
$g(s_1, \ldots, s_n) = (s'_1, \ldots, s'_n)$ を満たすものが存在する。
すなわち、$G$ の $F(Z_t)$ への作用は推移的であり、証明が完了する。
\end{proof}

\noindent
次が鍵となる補題である。

\begin{lemma}
\label{pione-lemma-tame}
\begin{reference}
\cite[Definition 7.2.4]{BS} と比較せよ。
\end{reference}
$(\mathcal{C}, F)$ をガロア圏とする。$G = \text{Aut}(F)$ は例
\ref{pione-example-from-C-F-to-G-sets} のものとする。任意の連結対象 $X$ で
$\mathcal{C}$ に属するものに対して、$G$ の $F(X)$ への作用は推移的である。
\end{lemma}

\begin{proof}
以下、補題 \ref{pione-lemma-epi-mono} の結果を断りなく用いる。
$I$ を、$\mathcal{C}$ におけるガロア対象の同型類全体の集合とする。
各 $i \in I$ に対して同型類の代表を $X_i$ とする。
$\gamma_i \in F(X_i)$ を各 $i \in I$ に対して選ぶ。$I$ 上の半順序を、
$i \geq i'$ であるための必要十分条件が、射 $f_{ii'} : X_i \to X_{i'}$ の
存在であるように定める。そのような射が与えられたとき、自己同型
$X_{i'} \to X_{i'}$ を後合成することにより
$F(f_{ii'})(\gamma_i) = \gamma_{i'}$ を満たすようにできる。この正規化のもとで、
射 $f_{ii'}$ は一意である。$I$ は有向半順序集合であることに注意する
（Categories, 定義 \ref{categories-definition-directed-set}）。実際、
$i_1, i_2 \in I$ なら、ガロア対象 $Y$ と射
$Y \to X_{i_1} \times X_{i_2}$ が存在する。これは補題 \ref{pione-lemma-galois} を
$X_{i_1} \times X_{i_2}$ の連結成分に適用すればよい。このとき
% P07-ERR-1147: frozen-source index case retained; see sidecar.
$Y \cong X_i$ がある $i \in I$ に対して成り立ち、$i \geq i_1$, $i \geq I_2$ である。

\medskip\noindent
関手 $F$ は関手 $F'$、すなわち $X$ を
$$
F'(X) = \colim_I \Mor_\mathcal{C}(X_i, X)
$$
へ写すものと同型であると主張する。この同型は、次のように定める関手の変換
$t : F' \to F$ による。$f : X_i \to X$ が与えられたとき、
$t_X(f) = F(f)(\gamma_i)$ と置く。（\ref{pione-item-one-element}）を用いると、
$t_X$ は単射である。全射性を示すため、$\gamma \in F(X)$ をとる。
ガロア圏の定義により、直ちに $X$ が連結である場合へ帰着できる。さらに補題
\ref{pione-lemma-galois} により、$X$ はガロアであるとしてよい。この場合、$X$ は
$X_i$ と同型であり、これはある $i$ に対して成り立つ。同型 $X_i \to X$ は、$\gamma_i$ が $\gamma$ に
写るように選べる（ガロア対象の定義による）。したがって $t$ は同型である。

\medskip\noindent
$A_i = \text{Aut}(X_i)$ と置く。$i \geq i'$ に対して標準的な写像
$h_{ii'} : A_i \to A_{i'}$ が存在し、すべての $a \in A_i$ に対して図式
$$
\xymatrix{
X_i \ar[d]_a \ar[r]_{f_{ii'}} & X_{i'} \ar[d]^{h_{ii'}(a)} \\
X_i \ar[r]^{f_{ii'}} & X_{i'}
}
$$
が可換すると主張する。実際、$h_{ii'}(a) = a' : X_{i'} \to X_{i'}$ を、
$F(a')(\gamma_{i'}) = F(f_{ii'} \circ a)(\gamma_i)$ を満たす一意な自己同型とすればよい。
前と同様に、これにより図式は可換し、しかも選択は一意である。したがって
$h_{i'i''} \circ h_{ii'} = h_{ii''}$
が $i \geq i' \geq i''$ のとき成り立つ。
$F(X_i) \to F(X_{i'})$ が全射なので、$A_i \to A_{i'}$ も全射である。
逆極限をとって群
$$
A = \lim_I A_i
$$
を得る。自己同型群は有限なので、これは副有限群である。Categories, 補題
\ref{categories-lemma-nonempty-limit} により、写像 $A \to A_i$ はすべての $i$ に対して
全射である。

\medskip\noindent
$A$ の元は逆系 $X_i$ に作用するので、前合成により $A$ の $F'$ への右作用を得る。
言い換えると、準同型 $A^{opp} \to G$ を得る。$A \to A_i$ は全射だから、
$G$ は $F(X_i)$ にすべての $i$ に対して推移的に作用する。各連結対象は
$X_i$ のいずれかで支配されるので、補題が従う。
\end{proof}

\begin{proposition}
\label{pione-proposition-galois}
\begin{reference}
これは \cite[Expos\'e V]{SGA1} の弱い形である。
証明は \cite[Theorem 7.2.5]{BS} による。
\end{reference}
$(\mathcal{C}, F)$ をガロア圏とする。$G = \text{Aut}(F)$ は例
\ref{pione-example-from-C-F-to-G-sets} のものとする。関手
$F : \mathcal{C} \to \textit{Finite-}G\textit{-Sets}$
% P07-ERR-1148: frozen-source missing copula recorded in sidecar.
(\ref{pione-equation-remember}) は圏同値である。
\end{proposition}

\begin{proof}
以下、補題 \ref{pione-lemma-epi-mono} の結果を断りなく用いる。とくに、関手が忠実であることは
既に分かっている。補題 \ref{pione-lemma-tame} により、任意の連結対象 $X$ に対する
$G$ の $F(X)$ への作用は推移的である。したがって、一般の $X$ に対し、関手 $F$ は
$X$ の連結成分（ガロア圏の公理により存在する）を、$G$ の $F(X)$ における軌道へ写す。
$X$ と $Y$ を対象とし、$s : F(X) \to F(Y)$ を写像とする。このときグラフ
$\Gamma_s \subset F(X) \times F(Y)$、すなわち $s$ のグラフは連結成分の和である。
したがって、連結成分の和 $Z$ で $X \times Y$ に属するものが存在し、単射 $Z \to X \times Y$ を備え、
$F(Z) = \Gamma_s$ を満たす。$F(Z) \to F(X)$ は全単射なので、$Z \to X$ は
同型である。ゆえに $s = F(f)$ であり、ここで $f : X \cong Z \to Y$ は合成である。
したがって $F$ は充満忠実である。

\medskip\noindent
証明を終えるため、$F$ が本質的全射であることを示す。$G/H$ が、有限指数の任意の
開部分群 $H \subset G$ に対して本質像に属することを示せば十分である。
$G$ 上の位相の定義により、有限個の対象 $X_i$ が存在し、
$$
\Ker(G \longrightarrow \prod\nolimits_i \text{Aut}(F(X_i)))
$$
は $H$ に含まれる。$X_i$ はすべての $i$ に対して連結であるとしてよい。
補題 \ref{pione-lemma-galois} により、ガロア対象 $Y$ で $\prod X_i$ の連結成分へ写るものを
選べる。同型 $F(Y) = G/U$ を $G\textit{-sets}$ において選ぶ。ここで
$U \subset G$ はある開部分群である。$Y$ はガロアなので、群
$\text{Aut}(Y) = \text{Aut}_{G\textit{-Sets}}(G/U)$ は $F(Y) = G/U$ に
推移的に作用する。これは $U$ が正規であることを意味する。$F(Y)$ は
$F(X_i)$ 上へ各 $i$ に対して全射となるので、$U \subset H$ である。
$M \subset \text{Aut}(Y)$ を、次に対応する有限部分群とする。
$$
(H/U)^{opp} \subset (G/U)^{opp} = \text{Aut}_{G\textit{-Sets}}(G/U)
= \text{Aut}(Y).
$$
$X = Y/M$ と置く。すなわち、$X$ は射 $m : Y \to Y$、$m \in M$ の
余等化子である。$F$ は完全なので $F(X) = G/H$ となり、証明が完了する。
\end{proof}

\begin{lemma}
\label{pione-lemma-functoriality-galois}
$(\mathcal{C}, F)$ および $(\mathcal{C}', F')$ をガロア圏とする。
$H : \mathcal{C} \to \mathcal{C}'$ を完全関手とする。同型
$t : F' \circ H \to F$ が存在する。$t$ の選択は連続準同型
$h : G' = \text{Aut}(F') \to \text{Aut}(F) = G$ と
$2$-可換図式
$$
\xymatrix{
\mathcal{C} \ar[r]_H \ar[d] & \mathcal{C}' \ar[d] \\
\textit{Finite-}G\textit{-Sets} \ar[r]^h &
\textit{Finite-}G'\textit{-Sets}
}
$$
を定める。写像 $h$ は $t$ の選択に依存するが、$G$ の内部自己同型を除けば
一意である。逆に、連続準同型 $h : G' \to G$ が与えられると、完全関手
$H : \mathcal{C} \to \mathcal{C}'$ と同型 $t$ が存在し、上の構成によって
$h$ が復元される。
\end{lemma}

\begin{proof}
命題 \ref{pione-proposition-galois} と補題 \ref{pione-lemma-single-out-profinite} により、
$\mathcal{C} = \textit{Finite-}G\textit{-Sets}$ で $F$ は忘却関手であるとしてよく、
$\mathcal{C}'$ についても同様である。したがって $t$ の存在は補題
\ref{pione-lemma-second-fundamental-functor} から従う。写像 $h$ は $t$ による構造の移送から得られる。
図式の可換性は明らかである。$h$ が $G$ の元による内部共役を除いて一意であることは、
$t$ の選択が $G$ の元を除いて一意であることから従う。最後の主張は直ちに分かる。
\end{proof}





\section{関手と準同型}
\label{pione-section-translation}

\noindent
$(\mathcal{C}, F)$、$(\mathcal{C}', F')$、$(\mathcal{C}'', F'')$ をガロア圏とする。
$G = \text{Aut}(F)$、$G' = \text{Aut}(F')$、$G'' = \text{Aut}(F'')$ と置く。
$H : \mathcal{C} \to \mathcal{C}'$ および $H' : \mathcal{C}' \to \mathcal{C}''$ を
完全関手とする。補題 \ref{pione-lemma-functoriality-galois} のように、それぞれに対応する
連続準同型を $h : G' \to G$ および $h' : G'' \to G'$ とする。この節では、対応する
$2$-可換図式
\begin{equation}
\label{pione-equation-translation}
\vcenter{
\xymatrix{
\mathcal{C} \ar[r]_H \ar[d] &
\mathcal{C}' \ar[r]_{H'} \ar[d] &
\mathcal{C}'' \ar[d] \\
\textit{Finite-}G\textit{-Sets} \ar[r]^h &
\textit{Finite-}G'\textit{-Sets} \ar[r]^{h'} &
\textit{Finite-}G''\textit{-Sets}
}
}
\end{equation}
を考え、列 $1 \to G'' \to G' \to G \to 1$ の完全性に関する性質と、関手 $H$ および
$H'$ の性質とを関係づける。

\begin{lemma}
\label{pione-lemma-functoriality-galois-surjective}
図式 (\ref{pione-equation-translation}) において、次の条件は同値である。
\begin{enumerate}
\item $h : G' \to G$ は全射である。
\item $H : \mathcal{C} \to \mathcal{C}'$ は充満忠実である。
\item $X \in \Ob(\mathcal{C})$ が連結なら、$H(X)$ も連結である。
\item $X \in \Ob(\mathcal{C})$ が連結で、射 $*' \to H(X)$ が
$\mathcal{C}'$ に存在するなら、射 $* \to X$ が存在する。
\item 任意の対象 $X$ で $\mathcal{C}$ に属するものに対して写像
$\Mor_\mathcal{C}(*, X) \to \Mor_{\mathcal{C}'}(*', H(X))$
は全単射である。
\end{enumerate}
ここで $*$ と $*'$ は、それぞれ $\mathcal{C}$ と $\mathcal{C}'$ の終対象である。
\end{lemma}

\begin{proof}
含意 (5) $\Rightarrow$ (4) および (2) $\Rightarrow$ (5) は明らかである。

\medskip\noindent
（3）を仮定する。$X$ を $\mathcal{C}$ の連結対象とし、$*' \to H(X)$ を射とする。
（3）により $H(X)$ は連結なので、$*' \to H(X)$ は同型である。したがって、
$G'$-集合で $H(X)$ に対応するものはちょうど一つの元をもち、これは
$G$-集合で $X$ に対応するものが一つの元をもつことを意味する。ゆえに $X$ は $\mathcal{C}$ の終対象と
同型であり、とくに写像 $* \to X$ が存在する。以上により (3) $\Rightarrow$ (4) である。

\medskip\noindent
（1）が成り立つなら、関手
$\textit{Finite-}G\textit{-Sets} \to \textit{Finite-}G'\textit{-Sets}$
は充満忠実である。実際、この場合、$G$-集合の写像が $G$ の作用と可換することと、
$G'$ の作用と可換することは同値である。したがって (1) $\Rightarrow$ (2) である。

\medskip\noindent
（1）が成り立つなら、$G$-集合 $X$ における $G$-軌道と $G'$-軌道は一致する。
したがって (1) $\Rightarrow$ (3) である。

\medskip\noindent
証明を終えるには、（4）が（1）を含意することを示せば十分である。（1）が偽、
すなわち $h$ が全射でないとする。このとき、開部分群 $U \subset G$ で
$h(G')$ を含み、$G$ に等しくないものが存在する。有限 $G$-集合 $M = G/U$ への
作用は推移的であるが、$G'$ は不動点をもつ。対象 $X$ で $\mathcal{C}$ に属し、$M$ に
対応するものは（3）に反する。以上により (3) $\Rightarrow$ (1) であり、証明が完了する。
\end{proof}

\begin{lemma}
\label{pione-lemma-composition-trivial}
図式 (\ref{pione-equation-translation}) において、次の条件は同値である。
\begin{enumerate}
\item $h \circ h'$ は自明である。
\item $H' \circ H$ の像は、終対象の有限余積と同型な対象からなる。
\end{enumerate}
\end{lemma}

\begin{proof}
$H$ および $H'$ を、標準的な関手
$\textit{Finite-}G\textit{-Sets} \to \textit{Finite-}G'\textit{-Sets}
\to \textit{Finite-}G''\textit{-Sets}$ で置き換えてよい。これは $h$ および $h'$ が
定める関手である。このとき主張は、
$G''$ のすべての $G$-集合への作用が自明であることと、準同型 $G'' \to G$ が
自明であることが同値だというものであり、これは明らかである。
\end{proof}

\begin{lemma}
\label{pione-lemma-functoriality-galois-ses}
図式 (\ref{pione-equation-translation}) において、次の条件は同値である。
\begin{enumerate}
\item 列 $G'' \xrightarrow{h'} G' \xrightarrow{h} G \to 1$ は、次の意味で完全である。
$h$ は全射、$h \circ h'$ は自明で、$\Ker(h)$ は $\Im(h')$ を含む最小の閉正規部分群である。
\item $H$ は充満忠実であり、対象 $X'$ で $\mathcal{C}'$ に属するものが $H$ の本質像に属するための
必要十分条件は、$H'(X')$ が終対象の有限余積と同型であることである。
% P07-ERR-1149: frozen-source composite order retained; see sidecar.
\item $H$ は充満忠実で、$H \circ H'$ はすべての対象を終対象の有限余積へ写す。
さらに、対象 $X'$ で $\mathcal{C}'$ に属するものが、$H'(X')$ が終対象の有限余積であるという条件を
満たすなら、対象 $X$ で $\mathcal{C}$ に属するものと全射 $H(X) \to X'$ が存在する。
\end{enumerate}
\end{lemma}

\begin{proof}
補題 \ref{pione-lemma-functoriality-galois-surjective} および
\ref{pione-lemma-composition-trivial} により、$H$ は充満忠実、$h$ は全射、
$H' \circ H$ は対象を終対象の非交和へ写し、$h \circ h'$ は自明であるとしてよい。
$N \subset G'$ を、$h'$ の像を含む最小の閉正規部分群とする。明らかに
$N \subset \Ker(h)$ である。関手 $H$ および $H'$ は標準的な関手
$\textit{Finite-}G\textit{-Sets} \to \textit{Finite-}G'\textit{-Sets}
\to \textit{Finite-}G''\textit{-Sets}$ であり、$h$ および $h'$ により定まるとしてよい。

\medskip\noindent
（2）を仮定する。これは、有限 $G'$-集合 $X'$ に $G''$ が自明に作用するなら、
$G'$ の作用が $G$ を経由して分解することを意味する。これを
$X' = G'/U'N$ に適用する。ここで $U'$ は $G'$ の十分小さい開部分群である。
すると、$\Ker(h) \subset U'N$ がすべての $U'$ に対して成り立つ。$N$ は閉だから、
$\Ker(h) \subset N$ が従う。すなわち（1）が成り立つ。

\medskip\noindent
（1）を仮定する。これは $N = \Ker(h)$ を意味する。$X'$ を有限 $G'$-集合とし、
$G''$ が自明に作用するものとする。これは $\Ker(G' \to \text{Aut}(X'))$ が $\Im(h')$ を含む
閉正規部分群であることを意味する。したがって $N = \Ker(h)$ はこの部分群に含まれ、
$G'$ の $X'$ への作用は $G$ を経由して分解する。すなわち（2）が成り立つ。

\medskip\noindent
（3）を仮定する。これは、有限 $G'$-集合 $X'$ に $G''$ が自明に作用するなら、
$G'$-集合の全射 $X \to X'$ で、$X$ が $G$-集合であるものが存在することを意味する。
明らかに、これは $G'$ の $X'$ への作用が $G$ を経由して分解することを意味する。
すなわち（2）が成り立つ。

\medskip\noindent
含意 (2) $\Rightarrow$ (3) は直ちに従う。これで証明が完了する。
\end{proof}

\begin{lemma}
\label{pione-lemma-functoriality-galois-injective}
図式 (\ref{pione-equation-translation}) において、次の条件は同値である。
\begin{enumerate}
\item $h'$ は単射である。
\item 任意の連結対象 $X''$ で $\mathcal{C}''$ に属するものに対し、
対象 $X'$ で $\mathcal{C}'$ に属するものと図式
% P07-ERR-0644: frozen-source functor symbol retained; see sidecar.
$$
X'' \leftarrow Y'' \rightarrow H(X')
$$
が $\mathcal{C}''$ に存在し、$Y'' \to X''$ は全射、$Y'' \to H(X')$ は単射である。
\end{enumerate}
\end{lemma}

\begin{proof}
$H'$ を、有限 $G'$-集合の圏から有限 $G''$-集合の圏への対応する関手で
置き換えてよい。

\medskip\noindent
$h' : G'' \to G'$ が単射であると仮定する。$H'' \subset G''$ を開部分群とする。
$G''$ の位相は $G'$ から誘導される位相なので、開部分群 $H' \subset G'$ で
$(h')^{-1}(H') \subset H''$ を満たすものが存在する。求める図式は
$$
G''/H'' \leftarrow G''/(h')^{-1}(H') \rightarrow G'/H'
$$
である。逆に、関手
$\textit{Finite-}G'\textit{-Sets} \to \textit{Finite-}G''\textit{-Sets}$.
に対して（2）が成り立つと仮定する。$g'' \in \Ker(h')$ とする。任意の開部分群
$H'' \subset G''$ を選ぶ。仮定により有限 $G'$-集合 $X'$ と図式
$$
G''/H'' \leftarrow Y'' \rightarrow X'
$$
が $G''$-集合の圏に存在し、左の矢印は全射、右の矢印は単射である。
$g''$ は $h'$ の核に属するので、$g''$ は $X'$ に自明に作用する。
したがって $g''$ は $Y''$ に自明に作用し、さらに $G''/H''$ にも自明に作用する。
ゆえに $g'' \in H''$ である。これはすべての開部分群について成り立つから、
$g''$ は所望のとおり単位元である。
\end{proof}

\begin{lemma}
\label{pione-lemma-functoriality-galois-normal}
図式 (\ref{pione-equation-translation}) において、次の条件は同値である。
\begin{enumerate}
\item $h'$ の像は正規である。
\item 任意の連結対象 $X'$ で $\mathcal{C}'$ に属し、$\mathcal{C}''$ の終対象から
$H'(X')$ への射が存在するものに対して、$H'(X')$ は終対象の有限余積と同型である。
\end{enumerate}
\end{lemma}

\begin{proof}
これは、連続群準同型 $h' : G'' \to G'$ に関する次の主張に移し換えられる。
$h'$ の像が正規であるための必要十分条件は、任意の開部分群 $U' \subset G'$ で
$h'(G'')$ を含むものが $h'(G'')$ のすべての共役も含むことである。これから結果は
容易に従う。詳細の一部は省略する。
\end{proof}






\section{有限エタール射}
\label{pione-section-finite-etale}

\noindent
この節では、エタール基本群を構成するのに十分な、有限エタール射に関する
基本的な結果を証明する。

\medskip\noindent
$X$ をスキームとする。記号 $\textit{F\'Et}_X$ で、$X$ 上有限かつエタールな
スキームの圏を表す。したがって、
\begin{enumerate}
\item $\textit{F\'Et}_X$ の対象は、有限エタール射 $Y \to X$ であって終域が $X$ のものである。
\item $\textit{F\'Et}_X$ における $Y \to X$ から $Y' \to X$ への射は、
射 $Y \to Y'$ であって図式
$$
\xymatrix{
Y \ar[rr] \ar[rd] & &  Y' \ar[ld] \\
& X
}
$$
を可換にするものである。
\end{enumerate}
$\textit{F\'Et}_X$ の対象を、しばしば $X$ の {\it 有限エタール被覆} と呼ぶ
（$Y$ が空の場合も含む）。スキームの圏上にはスタック
$p : \textit{F\'Et} \to \Sch$ が存在し、その $X$ 上のファイバーは今定義した圏
$\textit{F\'Et}_X$ である。Examples of Stacks, 節
\ref{examples-stacks-section-finite-etale} を参照。

\begin{example}
\label{pione-example-finite-etale-geometric-point}
$k$ を代数閉体とし、$X = \Spec(k)$ とする。この場合、$\textit{F\'Et}_X$ は
有限集合の圏と同値である。より一般に、$k$ が分離代数閉である場合にも成り立つ。
これは、$k$ 上エタールなスキームが、$k$ 上有限分離な体のスペクトルの
非交和だからである。Morphisms, 補題 \ref{morphisms-lemma-etale-over-field} を参照。
\end{example}

\begin{lemma}
\label{pione-lemma-finite-etale-covers-limits-colimits}
$X$ をスキームとする。圏 $\textit{F\'Et}_X$ は有限極限および有限余極限をもち、
任意の射 $X' \to X$ に対して、基底変換関手
$\textit{F\'Et}_X \to \textit{F\'Et}_{X'}$ は完全である。
\end{lemma}

\begin{proof}
有限極限と左完全性。Categories, 補題 \ref{categories-lemma-finite-limits-exist} により、
$\textit{F\'Et}_X$ が終対象とファイバー積をもつことを示せば十分である。これは、
$X$ 上のすべてのスキームの圏が終対象（すなわち $X$）とファイバー積をもつことから
明らかである。また、$X$ 上有限エタールなスキームのファイバー積も
$X$ 上有限エタールである。さらに、基底変換がこれらの操作と可換することは明らかで、
したがって基底変換は左完全である（Categories, 補題
\ref{categories-lemma-characterize-left-exact}）。

\medskip\noindent
有限余極限と右完全性。Categories, 補題 \ref{categories-lemma-colimits-exist} により、
$\textit{F\'Et}_X$ が有限余積と余等化子をもつことを示せば十分である。有限余積は
非交和で与えられる（空余積は空スキームである）。$a, b : Z \to Y$ を
$\textit{F\'Et}_X$ の二つの射とする。$Z \to X$ および $Y \to X$ は有限エタールなので、
$Z = \underline{\Spec}(\mathcal{C})$ および $Y = \underline{\Spec}(\mathcal{B})$ と書ける。
ここで $\mathcal{O}_X$ 上の有限局所自由代数がそれぞれ $\mathcal{C}$ および $\mathcal{B}$ である。
射 $a, b$ は二つの写像 $a^\sharp, b^\sharp : \mathcal{B} \to \mathcal{C}$ を誘導する。
$\mathcal{A} = \text{Eq}(a^\sharp, b^\sharp)$ をその等化子とする。もし
$$
\underline{\Spec}(\mathcal{A}) \longrightarrow X
$$
が有限エタールなら、これが余等化子であることは明らかである（実際、
$\textit{F\'Et}_X$ の任意の対象は $\mathcal{O}_X$-代数の層の相対スペクトルとして書ける）。
この確認は、$X$ をエタール被覆の要素で置き換えた後で行ってよい
（Descent, 補題 \ref{descent-lemma-descending-property-finite} および
\ref{descent-lemma-descending-property-etale}）。したがって \'Etale Morphisms, 補題
\ref{etale-lemma-finite-etale-etale-local} により、
$Y = \coprod_{i = 1, \ldots, n} X$ および $Z = \coprod_{j = 1, \ldots, m} X$
としてよい。このとき
$$
\mathcal{C} = \prod\nolimits_{1 \leq j \leq m} \mathcal{O}_X
\quad\text{かつ}\quad
\mathcal{B} = \prod\nolimits_{1 \leq i \leq n} \mathcal{O}_X
$$
である。さらに開被覆の要素で置き換えることにより、$a, b$ は写像
$a_s, b_s : \{1, \ldots, m\} \to \{1, \ldots, n\}$ に対応するとしてよい。
すなわち、成分 $X$ で $Z$ の添字 $j$ に対応するものは、成分 $X$ で $Y$ のそれぞれ添字
$a_s(j)$、$b_s(j)$ に対応するものへ、射 $a$、$b$ によって写る。
$\{1, \ldots, n\} \to T$ を $a_s, b_s$ の余等化子とする。このとき
$$
\mathcal{A} = \prod\nolimits_{t \in T} \mathcal{O}_X
$$
であり、そのスペクトルは明らかに $X$ 上有限エタールである。これが基底変換と
両立することの確認は省略する。したがって基底変換は右完全関手である。
\end{proof}

\begin{remark}
\label{pione-remark-colimits-commute-forgetful}
$X$ をスキームとする。自然な関手
$F_1 : \textit{F\'Et}_X \to \Sch$ および $F_2 : \textit{F\'Et}_X \to \Sch/X$
を考える。このとき、
\begin{enumerate}
\item 関手 $F_1$ および $F_2$ は有限余極限と可換する。
\item 関手 $F_2$ は有限極限と可換する。
\item 関手 $F_1$ は連結な有限極限、すなわち等化子およびファイバー積と可換する。
\end{enumerate}
極限に関する結果は、補題 \ref{pione-lemma-finite-etale-covers-limits-colimits} の証明における議論と
Categories, 補題 \ref{categories-lemma-connected-limit-over-X} から直ちに従う。
$F_1$ および $F_2$ が有限余積と可換することは明らかである。Categories, 補題
\ref{categories-lemma-characterize-left-exact} の双対により、$F_1$ および $F_2$ が
余等化子と可換することを示せばよい。補題
\ref{pione-lemma-finite-etale-covers-limits-colimits} の証明で、$\textit{F\'Et}_X$ における
余等化子はエタール局所的に次の形になることを見た。
$$
\xymatrix{
\coprod_{j \in J} U \ar@<1ex>[r]^a \ar@<-1ex>[r]_b &
\coprod_{i \in I} U \ar[r] &
\coprod_{t \in \text{Coeq}(a, b)} U
}
$$
これは明らかにスキームの圏における余等化子である。したがって主張は、余等化子で
あることが fpqc 局所的であるという事実から従う。その正確な定式化は
Descent, 補題 \ref{descent-lemma-coequalizer-fpqc-local} にある。
\end{remark}

\begin{lemma}
\label{pione-lemma-internal-hom-finite-etale}
$X$ をスキームとする。$U, V$ が $X$ 上有限エタールであるとき、
スキーム $W$ で $X$ 上有限エタールなものが存在して
$$
\Mor_X(X, W) = \Mor_X(U, V)
$$
を満たし、しかもこの等式は任意の基底変換後にも成り立つ。
\end{lemma}

\begin{proof}
More on Morphisms, 補題
\ref{more-morphisms-lemma-hom-from-finite-locally-free-separated-lqf} により、
スキーム $W$ で $\mathit{Mor}_X(U, V)$ を表現するものが存在する（エタール射は局所準有限で
あることを Morphisms, 補題 \ref{morphisms-lemma-etale-locally-quasi-finite} から用い、
有限射が分離的であることも用いる）。このスキームが任意の基底変換後にも等式を
満たすことは明らかである。証明を終えるには、$W \to X$ が有限エタールであることを
示さなければならない。これは、$X$ をエタール被覆の要素で置き換えた後に示してよい
（Descent, 補題 \ref{descent-lemma-descending-property-finite} および
\ref{descent-lemma-descending-property-separated}）。したがって \'Etale Morphisms, 補題
\ref{etale-lemma-finite-etale-etale-local} により、
$U = \coprod_{i = 1, \ldots, n} X$ および $V = \coprod_{j = 1, \ldots, m} X$
としてよい。この場合、調べれば
$W = \coprod_{\alpha : \{1, \ldots, n\} \to \{1, \ldots, m\}} X$
である（詳細は省略する）。これで証明が完了する。
\end{proof}

\noindent
$X$ をスキームとする。$X$ の {\it 幾何点} とは、射 $\Spec(k) \to X$ であって、
$k$ が代数閉体であるもののことである。このような点は通常、上線を付けた小文字
$\overline{x}$ で表す。$\overline{x}$ は射だけでなくスキーム $\Spec(k)$ を表すためにも
用い、$\kappa(\overline{x})$ で $k$ を表す。$\overline{x}$ が $x$ の {\it 上にある} とは、
$x \in X$ が $\overline{x}$ の像であることをいう。これについては
\'Etale Cohomology, 節 \ref{etale-cohomology-section-stalks} でさらに論じる。
$\overline{x}$ とエタール射 $U \to X$ が与えられたとき、次を考えられる。
$$
|U_{\overline{x}}| : \text{スキームの点からなる台集合 }U_{\overline{x}} = U \times_X \overline{x}
$$
$U_{\overline{x}}$ は $\overline{x}$ 上のスキームとして $\overline{x}$ のコピーの
非交和であるから（Morphisms, 補題 \ref{morphisms-lemma-etale-over-field}）、
この集合は次のようにも記述できる。
$$
|U_{\overline{x}}| =
\left\{
\begin{matrix}
\text{可換} \\
\text{図式}
\end{matrix}
\vcenter{
\xymatrix{
\overline{x} \ar[rd]_{\overline{x}} \ar[r]_{\overline{u}} & U \ar[d] \\
& X
}
}
\right\}
$$
対応 $U \mapsto |U_{\overline{x}}|$ は関手であり、しばしば
$F_{\overline{x}}$ と書く。

\begin{lemma}
\label{pione-lemma-finite-etale-connected-galois-category}
$X$ を連結スキームとし、$\overline{x}$ を幾何点とする。関手
$$
F_{\overline{x}} : \textit{F\'Et}_X \longrightarrow \textit{Sets},\quad
Y \longmapsto |Y_{\overline{x}}|
$$
はガロア圏を定める（定義 \ref{pione-definition-galois-category}）。
\end{lemma}

\begin{proof}
$\textit{F\'Et}_{\overline{x}}$ を有限集合の圏と同一視すると
（例 \ref{pione-example-finite-etale-geometric-point}）、関手 $F_{\overline{x}}$ は射
$\overline{x} \to X$ に対する基底変換関手にほかならない。したがって
$\textit{F\'Et}_X$ は有限極限および有限余極限をもち、補題
\ref{pione-lemma-finite-etale-covers-limits-colimits} により $F_{\overline{x}}$ は完全である。
また、$\textit{F\'Et}_X$ における有限極限は、$X$ 上のスキームの圏における
対応する有限極限と一致することも用いる。注意
\ref{pione-remark-colimits-commute-forgetful} を参照。

\medskip\noindent
$Y' \to Y$ が $\textit{F\'Et}_X$ における単射なら、
$Y' \to Y' \times_Y Y'$ は同型であり、したがって $Y' \to Y$ はスキームの単射である。
ゆえに $Y' \to Y$ は開埋め込みである（\'Etale Morphisms, 定理
\ref{etale-theorem-etale-radicial-open}）。$Y'$ は $X$ 上有限であり、$Y$ は $X$ 上分離的なので、
射 $Y' \to Y$ は有限であり（Morphisms, 補題 \ref{morphisms-lemma-finite-permanence}）、
したがって閉である（Morphisms, 補題 \ref{morphisms-lemma-finite-proper}）。ゆえにこれは
$Y$ の開閉部分スキームの包含である。したがって、$Y$ が圏 $\textit{F\'Et}_X$ の
連結対象であること（定義 \ref{pione-definition-galois-category} の意味で）と、$Y$ がスキームとして
連結であることは同値である。このとき Topology, 補題
\ref{topology-lemma-finite-fibre-connected-components} により、$Y$ はスキームとしても、
定義 \ref{pione-definition-galois-category} の意味でも、その連結成分の有限余積である。

\medskip\noindent
$Y \to Z$ を $\textit{F\'Et}_X$ の射で、全単射
$F_{\overline{x}}(Y) \to F_{\overline{x}}(Z)$ を誘導するものとする。
$Y \to Z$ が同型であることを示さなければならない。上により $Z$ は連結であるとしてよい。
$Y \to Z$ は有限エタール、したがって有限局所自由なので、$Y \to Z$ が次数 $1$ の
有限局所自由射であることを示せば十分である。これは、$Z$ の任意の点で
$\overline{x}$ の上にあるものの近傍で成り立つ。$Z$ は連結で次数は局所定数だから、結論が従う。
\end{proof}



\section{基本群}
\label{pione-section-fundamental-groups}

\noindent
この節では Grothendieck の代数的基本群を定義する。次の定義が意味をもつのは、
補題 \ref{pione-lemma-finite-etale-connected-galois-category} による。

\begin{definition}
\label{pione-definition-fundamental-group}
$X$ を連結スキームとし、$\overline{x}$ を $X$ の幾何点とする。
$X$ の {\it 基本群} で {\it 基点} $\overline{x}$ をもつものとは、群
$$
\pi_1(X, \overline{x}) = \text{Aut}(F_{\overline{x}})
$$
すなわちファイバー関手
$F_{\overline{x}} : \textit{F\'Et}_X \to \textit{Sets}$
の自己同型群に、補題 \ref{pione-lemma-aut-inverse-limit} の標準的な副有限位相を入れたものである。
\end{definition}

\noindent
上の結果と節 \ref{pione-section-galois} の内容を合わせると、次の定理を得る。

\begin{theorem}
\label{pione-theorem-fundamental-group}
$X$ を連結スキームとし、$\overline{x}$ を $X$ の幾何点とする。
\begin{enumerate}
\item ファイバー関手 $F_{\overline{x}}$ は圏同値
$$
\textit{F\'Et}_X \longrightarrow
\textit{Finite-}\pi_1(X, \overline{x})\textit{-Sets}
$$
を定める。
\item 第二の幾何点 $\overline{x}'$ が $X$ に与えられると、同型
$t : F_{\overline{x}} \to F_{\overline{x}'}$ が存在する。これは（1）の圏同値と両立する同型
$\pi_1(X, \overline{x}) \to \pi_1(X, \overline{x}')$ を与える。この同型は、内部共役を除けば
$t$ に依存しない。
\item 連結スキームの射 $f : X \to Y$ が与えられたとき、
$\overline{y} = f \circ \overline{x}$ と書く。標準的な連続準同型
$$
f_* : \pi_1(X, \overline{x}) \to \pi_1(Y, \overline{y})
$$
が存在し、図式
$$
\xymatrix{
\textit{F\'Et}_Y \ar[r]_{\text{基底変換}} \ar[d]_{F_{\overline{y}}} &
\textit{F\'Et}_X \ar[d]^{F_{\overline{x}}} \\
\textit{Finite-}\pi_1(Y, \overline{y})\textit{-Sets} \ar[r]^{f_*} &
\textit{Finite-}\pi_1(X, \overline{x})\textit{-Sets}
}
$$
は可換である。
\end{enumerate}
\end{theorem}

\begin{proof}
（1）は補題 \ref{pione-lemma-finite-etale-connected-galois-category} と命題
\ref{pione-proposition-galois} から従う。（2）は補題 \ref{pione-lemma-functoriality-galois} の特別な場合である。
（3）について、図式
$$
\xymatrix{
\textit{F\'Et}_Y \ar[r] \ar[d]_{F_{\overline{y}}} &
\textit{F\'Et}_X \ar[d]^{F_{\overline{x}}} \\
\textit{Sets} \ar@{=}[r] & \textit{Sets}
}
$$
は可換である（単なる $2$-可換ではなく、実際に可換である）。これは
$\overline{y} = f \circ \overline{x}$ だからである。したがって、暗黙の関手の変換とともに
補題 \ref{pione-lemma-functoriality-galois} を適用すれば（3）を得る。
\end{proof}

\begin{lemma}
\label{pione-lemma-fundamental-group-Galois-group}
$K$ を体とし、$X = \Spec(K)$ と置く。$\overline{K}$ を代数閉包とし、対応する幾何点を
$\overline{x} : \Spec(\overline{K}) \to X$ と書く。$K^{sep} \subset \overline{K}$ を
分離代数閉包とする。
\begin{enumerate}
\item 補題 \ref{pione-lemma-sheaves-point} の関手は圏同値
$$
\textit{F\'Et}_X \longrightarrow
\textit{Finite-}\text{Gal}(K^{sep}/K)\textit{-Sets}.
$$
を誘導し、$F_{\overline{x}}$ および関手
$\textit{Finite-}\text{Gal}(K^{sep}/K)\textit{-Sets} \to \textit{Sets}$ と両立する。
\item これは標準的な同型
$$
\text{Gal}(K^{sep}/K) \longrightarrow \pi_1(X, \overline{x})
$$
を副有限位相群の間に誘導する。
\end{enumerate}
\end{lemma}

\begin{proof}
補題 \ref{pione-lemma-sheaves-point} の関手は関手 $F_{\overline{x}}$ と同じである。実際、
任意の $Y$ で $X$ 上エタールなものに対して
$$
\Mor_X(\Spec(\overline{K}), Y) = \Mor_X(\Spec(K^{sep}), Y)
$$
である。すなわち、補題 \ref{pione-lemma-sheaves-point} の証明で見たように、
$Y = \coprod_{i \in I} \Spec(L_i)$ であり、$L_i/K$ は $K$ 上有限分離である。
したがって任意の $K$-代数準同型 $L_i \to \overline{K}$ は $K^{sep}$ を経由して分解する。
また、$F_{\overline{x}}(Y)$ が有限であること、$I$ が有限であること、および
$Y \to X$ が有限エタールであることは互いに同値である。これで（1）が示された。

\medskip\noindent
（2）は（1）、補題 \ref{pione-lemma-functoriality-galois}、および補題
\ref{pione-lemma-single-out-profinite} の形式的な帰結である（下の注意も参照）。
\end{proof}

\begin{remark}
\label{pione-remark-variance}
補題 \ref{pione-lemma-fundamental-group-Galois-group} の状況で、同型
$\text{Gal}(K^{sep}/K) \to
\pi_1(X, \overline{x}) = \text{Aut}(F_{\overline{x}})$.
をより明示的に構成しよう。等式
$\text{Gal}(K^{sep}/K) = \text{Aut}(\overline{K}/K)$
が成り立つ。実際、$\overline{K}$ は $K^{sep}$ の完全化である。
$F_{\overline{x}}(Y) = \Mor_X(\Spec(\overline{K}), Y)$ なので、写像
$$
\text{Aut}(\overline{K}/K) \times F_{\overline{x}}(Y) \to F_{\overline{x}}(Y),
\quad
(\sigma, \overline{y}) \mapsto
\sigma \cdot \overline{y} = \overline{y} \circ \Spec(\sigma)
$$
を考えられる。これは作用である。実際、
$$
\sigma\tau \cdot \overline{y} =
\overline{y} \circ \Spec(\sigma\tau) =
\overline{y} \circ \Spec(\tau) \circ \Spec(\sigma) =
\sigma \cdot (\tau \cdot \overline{y})
$$
である。この作用は $Y \in \textit{F\'Et}_X$ に関して関手的であり、所望の写像を得る。
\end{remark}





\section{連結スキームのガロア被覆}
\label{pione-section-finite-etale-under-galois}

\noindent
$X$ を、幾何点 $\overline{x}$ をもつ連結スキームとする。
$F_{\overline{x}} : \textit{F\'Et}_X \to \textit{Sets}$ はガロア圏なので
（補題 \ref{pione-lemma-finite-etale-connected-galois-category}）、節 \ref{pione-section-galois} の
内容が適用できる。この節では、その用語と結果の一部をスキームおよび有限エタール射の
設定へ明示的に移す。

\medskip\noindent
有限エタール射 $Y \to X$ が {\it ガロア被覆} であるとは、$Y$ が
$\textit{F\'Et}_X$ のガロア対象を定めることをいう。有限エタール射 $Y \to X$ に対し
$G = \text{Aut}_X(Y)$ と置くと、次の条件は同値である。
\begin{enumerate}
\item $Y$ は $X$ のガロア被覆である。
\item $Y$ は連結であり、$|G|$ は $Y \to X$ の次数に等しい。
\item $Y$ は連結であり、$G$ は $F_{\overline{x}}(Y)$ に推移的に作用する。
\item $Y$ は連結であり、$G$ は $F_{\overline{x}}(Y)$ に単純推移的に作用する。
\end{enumerate}
これは節 \ref{pione-section-galois} の議論から直ちに従う。

\medskip\noindent
任意の有限エタール射 $f : Y \to X$ で $Y$ が連結であるものに対し、有限エタール・
ガロア被覆 $Y' \to X$ で $Y$ を支配するものが存在する（補題 \ref{pione-lemma-galois}）。

\medskip\noindent
$\textit{F\'Et}_X$ のガロア対象は、同値
$$
F_{\overline{x}} : \textit{F\'Et}_X \to
\textit{Finite-}\pi_1(X, \overline{x})\textit{-Sets}
$$
（定理 \ref{pione-theorem-fundamental-group}）を通じて、有限
$\pi_1(X, \overline{x})\textit{-Sets}$ で形 $G = \pi_1(X, \overline{x})/H$ のものに対応する。
ここで $H$ は正規開部分群である。同値な言い方をすれば、$G$ が有限群であり、
$\pi_1(X, \overline{x}) \to G$ が連続全射なら、$G$ を
$\pi_1(X, \overline{x})$-集合と見たものはガロア被覆に対応する。

\medskip\noindent
$Y_i \to X$、$i = 1, 2$ がガロア群 $G_i$ をもつ有限エタール・ガロア被覆なら、
有限エタール・ガロア被覆 $Y \to X$ で、そのガロア群が $G_1 \times G_2$ の
部分群であるものが存在する。実際、対応する連続準同型
$\pi_1(X, \overline{x}) \to G_i$ をとり、$G$ を誘導された連続準同型
$\pi_1(X, \overline{x}) \to G_1 \times G_2$ の像とすればよい。









\section{基本群の位相的不変性}
\label{pione-section-topological-invariance}

\noindent
この節の主結果は、連結スキームの普遍同相写像が基本群の同型を誘導するというものである。
命題 \ref{pione-proposition-universal-homeomorphism} を参照。

\medskip\noindent
二つのスキームの基本群が同じであることを直接証明する代わりに、その有限エタール被覆の
圏が同じであることを証明する場合が多い。両者が連結なら、もちろんこれから基本群が
等しいことが従う。

\begin{lemma}
\label{pione-lemma-what-equivalence-gives}
$f : X \to Y$ を準コンパクト準分離スキームの射とし、基底変換関手
$\textit{F\'Et}_Y \to \textit{F\'Et}_X$ が圏同値であるとする。このとき、
\begin{enumerate}
\item $f$ は同相写像 $\pi_0(X) \to \pi_0(Y)$ を誘導する。
\item $X$、または同値なことに $Y$ が連結なら、
$\pi_1(X, \overline{x}) = \pi_1(Y, \overline{y})$ である。
\end{enumerate}
\end{lemma}

\begin{proof}
$Y = Y_0 \amalg Y_1$ を、空でない開閉部分スキームへの分解とする。
$f(X)$ は両方の $Y_i$ と交わると主張する。そうでなく、たとえば
$f(X) \subset Y_1$ なら、有限エタール射 $V = Y_1 \to Y$ を考えられる。
これは同型ではないが $V \times_Y X \to X$ は同型となり、矛盾である。

\medskip\noindent
$X = X_0 \amalg X_1$ を開閉部分スキームへの分解とする。有限エタール射
$U = X_1 \to X$ を考える。このとき $U = X \times_Y V$ であり、ここで $V \to Y$ は
ある有限エタール射である。射 $V \to Y$ の次数は局所定数だから、開閉部分スキームへの分解
$Y = \coprod_{d \geq 0} Y_d$ で、$V \to Y$ が次数 $d$ を $Y_d$ 上にもつものを得る。
$f^{-1}(Y_d) = \emptyset$ が $d > 1$ に対して成り立つので、上により
$Y_d = \emptyset$ が $d > 1$ に対して成り立つ。したがって
$f^{-1}(Y_i) = X_i$ が $i = 0, 1$ に対して成り立つ。

\medskip\noindent
したがって $f^{-1}$ は、$Y$ の開閉部分集合全体と $X$ の開閉部分集合全体との間の
全単射を誘導する。$X$ と $Y$ はスペクトル空間であることに注意する
（Properties, 補題 \ref{properties-lemma-quasi-compact-quasi-separated-spectral}）。
Topology, 補題 \ref{topology-lemma-connected-component-intersection} により、
スペクトル空間の開閉部分集合の束は連結成分の集合を決定する。ゆえに
$\pi_0(X) \to \pi_0(Y)$ は全単射である。$\pi_0(X)$ と $\pi_0(Y)$ は副有限空間なので
（Topology, 補題 \ref{topology-lemma-pi0-profinite}）、Topology, 補題
\ref{topology-lemma-bijective-map} により $\pi_0(X) \to \pi_0(Y)$ は同相写像である。
これで（1）が示された。（2）は直ちに従う。
\end{proof}

\noindent
次の補題は、Hensel 対の基本群がその閉部分集合の基本群であることを示す。

\begin{lemma}
\label{pione-lemma-gabber}
$(A, I)$ を Hensel 対とする。$X = \Spec(A)$ および $Z = \Spec(A/I)$ と置く。関手
$$
\textit{F\'Et}_X \longrightarrow \textit{F\'Et}_Z,\quad
U \longmapsto U \times_X Z
$$
は圏同値である。
\end{lemma}

\begin{proof}
これは More on Algebra, 補題 \ref{more-algebra-lemma-finite-etale-equivalence} の
移し換えである。
\end{proof}

\noindent
次の補題は、肥厚の基本群がもとの基本群と同じであることを示す。これは、下にある
普遍同相写像に関する強い命題の証明で用いる。

\begin{lemma}
\label{pione-lemma-thickening}
$X \subset X'$ をスキームの肥厚とする。関手
$$
\textit{F\'Et}_{X'} \longrightarrow \textit{F\'Et}_X,\quad
U' \longmapsto U' \times_{X'} X
$$
は圏同値である。
\end{lemma}

\begin{proof}
肥厚に関する議論は More on Morphisms, 節 \ref{more-morphisms-section-thickenings} を参照。
$U' \to X'$ をエタール射とし、$U = U' \times_{X'} X \to X$ が有限エタールであるとする。
このとき $U' \to X'$ も有限エタールである。これはたとえば More on Morphisms, 補題
\ref{more-morphisms-lemma-properties-that-extend-over-thickenings} から従う。
さて、$X \subset X'$ が有限次数の肥厚なら、この注意と \'Etale Morphisms, 定理
\ref{etale-theorem-remarkable-equivalence} を合わせれば補題が従う。以下では一般の肥厚に
ついて補題を証明するが、読者にはこの証明を読み飛ばすことを勧める。

\medskip\noindent
アフィン開被覆を $X' = \bigcup X_i'$ とする。次のように置く。
$X_i = X \times_{X'} X_i'$, $X_{ij}' = X'_i \cap X'_j$,
$X_{ij} = X \times_{X'} X_{ij}'$, $X_{ijk}' = X'_i \cap X'_j \cap X'_k$,
$X_{ijk} = X \times_{X'} X_{ijk}'$.
肥厚 $X_i \subset X'_i$、$X_{ij} \subset X_{ij}'$、および
$X_{ijk} \subset X_{ijk}'$ のそれぞれに対して定理を証明できるとする。このとき、
スキームの相対的貼り合わせにより $X \subset X'$ に対する結果が従う。
Constructions, 節 \ref{constructions-section-relative-glueing} を参照。スキーム
$X_i'$、$X_{ij}'$、$X_{ijk}'$ は、それぞれアフィンスキームの開部分スキームとして
分離的であることに注意する。この議論をもう一度繰り返すと、スキーム
$X'_i$、$X_{ij}'$、$X_{ijk}'$ がアフィンである場合へ帰着する。

\medskip\noindent
アフィンの場合、$X' = \Spec(A')$ および $X = \Spec(A'/I')$ であり、$I'$ は局所冪零イデアルである。
このとき $(A', I')$ は Hensel 対であり（More on Algebra, 補題
\ref{more-algebra-lemma-locally-nilpotent-henselian}）、結果は補題 \ref{pione-lemma-gabber} から従う
（この場合ははるかに容易である）。
\end{proof}

\noindent
次の命題を証明する「正しい」方法は、エタールサイトの不変性から導くことであろう。
\'Etale Cohomology, 定理 \ref{etale-cohomology-theorem-topological-invariance} を参照。

\begin{proposition}
\label{pione-proposition-universal-homeomorphism}
$f : X \to Y$ をスキームの普遍同相写像とする。このとき
$$
\textit{F\'Et}_Y \longrightarrow \textit{F\'Et}_X,\quad
V \longmapsto V \times_Y X
$$
は圏同値である。したがって $X$ と $Y$ が連結なら、$f$ は基本群の同型
$\pi_1(X, \overline{x}) \to \pi_1(Y, \overline{y})$ を誘導する。
\end{proposition}

\begin{proof}
普遍同相写像とは、整かつ普遍的単射かつ全射な射と同じものである。
Morphisms, 補題 \ref{morphisms-lemma-universal-homeomorphism} を参照。とくに対角射
$\Delta : X \to X \times_Y X$ は、Morphisms, 補題
\ref{morphisms-lemma-universally-injective} により肥厚である。したがって補題
\ref{pione-lemma-thickening} により、有限エタール射 $U \to X$ が与えられると、一意な同型
$$
\varphi : U \times_Y X \to X \times_Y U
$$
であって、$X \times_Y X$ 上有限エタールなスキームの同型であり、$\Delta$ による引き戻しが
$\text{id} : U \to U$ を $X$ 上で与えるものが存在する。$X \to X \times_Y X \times_Y X$ も
肥厚なので（これは全単射かつ閉埋め込みである）、$(U, \varphi)$ は $X/Y$ に関する
降下データである。\'Etale Morphisms, 命題 \ref{etale-proposition-effective} により、
$U = X \times_Y V$ である。ここで $V \to Y$ は準コンパクト、分離的、かつエタールである。
$V \to Y$ が有限であることの証明は省略する（ヒント：射 $U \to V$ は全射であり、
$U \to Y$ は整である）。したがって $\textit{F\'Et}_Y \to \textit{F\'Et}_X$ は
本質的全射である。

\medskip\noindent
上と同様に論じると、$V_1 \to Y$ および $V_2 \to Y$ が $\textit{F\'Et}_Y$ に属するものとして
与えられたとき、任意の射 $a : X \times_Y V_1 \to X \times_Y V_2$ は $X$ 上で
標準的な降下データと両立する。したがって \'Etale Morphisms, 補題
\ref{etale-lemma-fully-faithful-cases} により、$a$ は射 $V_1 \to V_2$ に $Y$ 上で降下する。
\end{proof}










\section{固有スキームの有限エタール被覆}
\label{pione-section-finite-etale-over-proper}

\noindent
この節では、ヘンゼル局所環上の連結固有スキームの基本群が、その特殊
ファイバーの基本群と同じであることを示す。また、この結果のヘンゼル対
に対する変種も証明する。

\medskip\noindent
さらに、代数閉体 $k$ 上の連結固有スキームの基本群は、$k$ を代数閉拡大体
で置き換えても変わらないことを示す。

\medskip\noindent
連結の場合に結果を述べて証明する代わりに、一般の場合に結果を証明し、
基本群についての結論は補題
\ref{pione-lemma-what-equivalence-gives} を用いて読者が導くことにする。

\begin{lemma}
\label{pione-lemma-finite-etale-on-proper-over-henselian}
$A$ をヘンゼル局所環とし、$X$ を $A$ 上の固有スキーム、閉ファイバーを
$X_0$ とする。このとき関手
$$
\textit{F\'Et}_X \to \textit{F\'Et}_{X_0},\quad
U \longmapsto U_0 = U \times_X X_0
$$
は圏同値である。
\end{lemma}

\begin{proof}
ここでの証明は、代数化と近似を適用する一例である。いくつかの段階に分けて進める。

\medskip\noindent
\medskip\noindent
$A$ が完備局所ネーター環の場合の本質的全射性。
$X_n = X \times_{\Spec(A)} \Spec(A/\mathfrak m^{n + 1})$ とおく。
\'Etale Morphisms, 定理 \ref{etale-theorem-remarkable-equivalence} により、包含
$$
X_0 \to X_1 \to X_2 \to \ldots
$$
は、$X_0$ 上エタールなスキームの圏と $X_n$ 上エタールなスキームの圏との間に
圏同値を誘導する。
さらに、$U_n \to X_n$ が有限エタール射 $U_0 \to X_0$ に対応するなら、
$U_n \to X_n$ も有限である。例えば More on Morphisms, 補題
\ref{more-morphisms-lemma-thicken-property-morphisms-cartesian}.
この場合、射 $U_0 \to \Spec(A/\mathfrak m)$ は、$X_0$ が $A/\mathfrak m$ 上固有なので固有である。したがって Grothendieck の代数化定理を
（次の形で）適用できる。
Cohomology of Schemes, 補題
\ref{coherent-lemma-algebraize-formal-scheme-finite-over-proper})
これにより、有限射 $U \to X$ が存在し、その $X_0$ への制限は $U_0$ を復元する。
More on Morphisms, 補題
\ref{more-morphisms-lemma-check-smoothness-on-infinitesimal-nbhds}
により、$U \to X$ は $U_0$ のすべての点でエタールである。
しかし $U$ の各点は $U_0$ の点へ特殊化する（$U$ は $A$ 上固有だから）ので、
$U \to X$ はエタールであると結論できる。このようにして関手が本質的全射であることが従う。

\medskip\noindent
\medskip\noindent
$A$ が完備局所ネーター環の場合の忠実充満性。
$U \to X$ と $V \to X$ を有限エタール射とし、$\varphi_0 : U_0 \to V_0$ を
$X_0$ 上の射とする。次の射を考える。
$$
\Gamma_{\varphi_0} : U_0 \longrightarrow U_0 \times_{X_0} V_0
$$
この射は有限エタール射であると同時に閉埋め込みである。
$X = U \times_X V$ に本質的全射性を適用すると、有限エタール射
$W \to U \times_X V$ が得られ、その特殊ファイバーは $\Gamma_{\varphi_0}$ と同型である。
射影 $W \to U$ を考える。これは構成により有限エタールであり、$U_0$ 上では同型である。
\'Etale Morphisms, 補題
\ref{etale-lemma-finite-etale-one-point}
$W \to U$ は $U_0$ の開近傍で同型である。したがってこれは同型であり、合成
$\varphi : U \cong W \to V$ は $\varphi_0$ の求める持ち上げである。

\medskip\noindent
\medskip\noindent
$A$ がヘンゼル局所ネーター G-環の場合の本質的全射性。
$U_0 \to X_0$ を有限エタール射とする。
$A^\wedge$ を極大イデアルに関する $A$ の完備化とし、$X^\wedge$ を
$X$ の $A^\wedge$ への基底変換とする。上の結果により、有限エタール射
$V \to X^\wedge$ が存在し、その特殊ファイバーは $U_0$ である。$A^\wedge = \colim A_i$ と書き、
$A \to A_i$ は有限型とする。
Limits, 補題 \ref{limits-lemma-descend-finite-presentation} により、ある $i$ と有限表示射
$U_i \to X_{A_i}$ が存在し、その $X^\wedge$ への基底変換は $V$ である。$i$ を大きく取れば
$U_i \to X_{A_i}$ は有限かつエタールであると仮定できる
（Limits, 補題 \ref{limits-lemma-descend-finite-finite-presentation} および
\ref{limits-lemma-descend-etale}）。次のように書く。
$$
A_i = A[x_1, \ldots, x_n]/(f_1, \ldots, f_m)
$$
環写像 $A_i \to A^\wedge$ は、解 $(a_1, \ldots, a_n)$ を $A^\wedge$ 内のものとして、
方程式系 $f_j = 0$ に対して読み替えられる。
Smoothing Ring Maps, 定理 \ref{smoothing-theorem-approximation-property} により、
この解を（例えば位数 $11$ まで）解 $(b_1, \ldots, b_n)$ で $A$ 内に近似できる。
逆に戻すと、$A$-代数写像 $A_i \to A$ が得られ、これは元の写像
$A_i \to A^\wedge$ と同じ閉点を与える（$11 > 1$ だからである）。したがって、この環写像による
$U \to X$ は $V \to X_{A_i}$ の基底変換であり、特殊ファイバーが $U_0$ と同型な有限エタール射となる。

\medskip\noindent
\medskip\noindent
$A$ がヘンゼル局所ネーター G-環の場合の忠実充満性。
これは $A$ が完備ネーター環の場合とまったく同じ方法で、本質的全射性から導かれる。

\medskip\noindent
\medskip\noindent
一般の場合。$(A, \mathfrak m)$ をヘンゼル局所環とする。$S = \Spec(A)$ とおき、
閉点を $s \in S$ と記す。Limits, 補題
\ref{limits-lemma-proper-limit-of-proper-finite-presentation-noetherian}
により、$X \to \Spec(A)$ を固有射 $X_i \to S_i$ の余フィルター付き極限として書ける。
$S_i$ は $\mathbf{Z}$ 上有限型である。各 $i$ について、$s_i \in S_i$ を $s$ の像とする。
$S = \lim S_i$ および $A = \mathcal{O}_{S, s}$ なので
$A = \colim \mathcal{O}_{S_i, s_i}$ である。環 $A_i = \mathcal{O}_{S_i, s_i}$ は
ネーター局所 G-環である（More on Algebra, 命題
\ref{more-algebra-proposition-ubiquity-G-ring}).
More on Algebra, 補題 \ref{more-algebra-lemma-henselization-colimit} により
$A = \colim A_i^h$ である。More on Algebra, 補題
\ref{more-algebra-lemma-henselization-G-ring} により環 $A_i^h$ は G-環である。
したがって $A = \colim A_i^h$ および
$$
X = \lim (X_i \times_{S_i} \Spec(A_i^h))
$$
がスキームとして成り立つ。$X$ 上有限エタールなスキームの圏は、
$X_i \times_{S_i} \Spec(A_i^h)$ 上有限エタールなスキームの圏の極限である（Limits, 補題
\ref{limits-lemma-descend-finite-presentation},
\ref{limits-lemma-descend-finite-finite-presentation} および
\ref{limits-lemma-descend-etale})
について同様である。$X_0 = \lim (X_i \times_{S_i} s_i)$ 上有限エタールなスキームにも
同じことが成り立つ。よって、上で扱った
$X / \Spec(A)$ に対する結果を、上で扱った
$(X_i \times_{S_i} \Spec(A_i^h)) / \Spec(A_i^h)$ に対する結果から形式的に導くことができる。
\end{proof}

\begin{lemma}
\label{pione-lemma-finite-etale-on-proper-over-henselian-pair}
$(A, I)$ をヘンゼル対とし、$X$ を $A$ 上の固有スキームとする。
$X_0 = X \times_{\Spec(A)} \Spec(A/I)$ とおく。このとき関手
$$
\textit{F\'Et}_X \to \textit{F\'Et}_{X_0},\quad
U \longmapsto U_0 = U \times_X X_0
$$
は圏同値である。
\end{lemma}

\begin{proof}
この補題の証明は、補題
\ref{pione-lemma-finite-etale-on-proper-over-henselian} の証明とまったく同じである。

\medskip\noindent
\medskip\noindent
$A$ がネーター環かつ $I$-進完備の場合の本質的全射性。
$X_n = X \times_{\Spec(A)} \Spec(A/I^{n + 1})$ とおく。
\'Etale Morphisms, 定理 \ref{etale-theorem-remarkable-equivalence} により、包含
$$
X_0 \to X_1 \to X_2 \to \ldots
$$
は、$X_0$ 上エタールなスキームの圏と $X_n$ 上エタールなスキームの圏との間に
圏同値を誘導する。
さらに、$U_n \to X_n$ が有限エタール射 $U_0 \to X_0$ に対応するなら、
$U_n \to X_n$ も有限である。例えば More on Morphisms, 補題
\ref{more-morphisms-lemma-thicken-property-morphisms-cartesian}.
この場合、射 $U_0 \to \Spec(A/I)$ は、$X_0$ が $A/I$ 上固有なので固有である。
したがって Grothendieck の代数化定理を（次の形で）適用できる。
Cohomology of Schemes, 補題
\ref{coherent-lemma-algebraize-formal-scheme-finite-over-proper})
これにより、有限射 $U \to X$ が存在し、その $X_0$ への制限は $U_0$ を復元する。
More on Morphisms, 補題
\ref{more-morphisms-lemma-check-smoothness-on-infinitesimal-nbhds}
により、$U \to X$ は $U_0$ のすべての点でエタールである。しかし $U$ の各点は
$U_0$ の点へ特殊化する（$U$ は $A$ 上固有だから）ので、$U \to X$ はエタールである。
このようにして関手が本質的全射であることが従う。

\medskip\noindent
\medskip\noindent
$A$ がネーター環かつ $I$-進完備の場合の忠実充満性。
$U \to X$ と $V \to X$ を有限エタール射とし、$\varphi_0 : U_0 \to V_0$ を
$X_0$ 上の射とする。次の射を考える。
$$
\Gamma_{\varphi_0} : U_0 \longrightarrow U_0 \times_{X_0} V_0
$$
この射は有限エタール射であると同時に閉埋め込みである。
$X = U \times_X V$ に本質的全射性を適用すると、有限エタール射
$W \to U \times_X V$ が得られ、その特殊ファイバーは $\Gamma_{\varphi_0}$ と同型である。
射影 $W \to U$ を考える。これは構成により有限エタールであり、$U_0$ 上では同型である。
\'Etale Morphisms, 補題
\ref{etale-lemma-finite-etale-one-point}
$W \to U$ は $U_0$ の開近傍で同型である。したがってこれは同型であり、合成
$\varphi : U \cong W \to V$ は $\varphi_0$ の求める持ち上げである。

\medskip\noindent
\medskip\noindent
$(A, I)$ がヘンゼル対で、$A$ がネーター G-環の場合の本質的全射性。
$U_0 \to X_0$ を有限エタール射とする。$A^\wedge$ を $A$ の $I$ に関する完備化とする。
$A^\wedge$ はネーター環であり $IA^\wedge$-進完備であることに注意する。Algebra, 補題
\ref{algebra-lemma-completion-complete} および
\ref{algebra-lemma-completion-Noetherian-Noetherian} を参照。$X^\wedge$ を $X$ の $A^\wedge$ への
基底変換とする。上の結果により、有限エタール射
$V \to X^\wedge$ が存在し、その特殊ファイバーは $U_0$ である。$A^\wedge = \colim A_i$ と書き、
$A \to A_i$ は有限型とする。Limits, 補題 \ref{limits-lemma-descend-finite-presentation} により、
ある $i$ と有限表示射 $U_i \to X_{A_i}$ が存在し、その $X^\wedge$ への基底変換は $V$ である。
$i$ を大きく取れば $U_i \to X_{A_i}$ は有限かつエタールであると仮定できる
（Limits, 補題 \ref{limits-lemma-descend-finite-finite-presentation} および
\ref{limits-lemma-descend-etale}）。次のように書く。
$$
A_i = A[x_1, \ldots, x_n]/(f_1, \ldots, f_m)
$$
環写像 $A_i \to A^\wedge$ は、解 $(a_1, \ldots, a_n)$ を $A^\wedge$ 内のものとして、
方程式系 $f_j = 0$ に対して読み替えられる。Smoothing Ring Maps, 補題
\ref{smoothing-lemma-henselian-pair} により、この解を（例えば位数 $11$ まで）解
$(b_1, \ldots, b_n)$ で $A$ 内に近似できる。逆に戻すと、$A$-代数写像 $A_i \to A$ が得られ、これは元の写像
$A_i \to A^\wedge$ と同じ閉点を与える（$11 > 1$ だからである）。したがって、この環写像による
$U \to X$ は $V \to X_{A_i}$ の基底変換であり、特殊ファイバーが $U_0$ と同型な有限エタール射となる。

\medskip\noindent
\medskip\noindent
$(A, I$ がヘンゼル対で、$A$ がネーター G-環の場合の忠実充満性。
これは $A$ が完備ネーター環の場合とまったく同じ方法で、本質的全射性から導かれる。

\medskip\noindent
\medskip\noindent
一般の場合。$(A, I)$ をヘンゼル対とする。$S = \Spec(A)$、$S_0 = \Spec(A/I)$ とおく。
Limits, 補題
\ref{limits-lemma-proper-limit-of-proper-finite-presentation-noetherian}
により、$X \to \Spec(A)$ を固有射 $X_i \to S_i$ の余フィルター付き極限として書ける。
ここで $S_i$ は $\mathbf{Z}$ 上アフィン有限型である。$S_i = \Spec(A_i)$ と書き、
$I_i \subset A_i$ を $I$ の逆像とし、これは写像 $A_i \to A$ によるものである。
$S_{i, 0} = \Spec(A_i/I_i)$ とおく。$S = \lim S_i$ なので $A = \colim A_i$ であり、
したがって $I = \colim I_i$ および $A/I = \colim A_i/I_i$ も成り立つ。
環 $A_i$ はネーター G-環である（More on Algebra, 命題
\ref{more-algebra-proposition-ubiquity-G-ring}).
$(A_i^h, I_i^h)$ を $(A_i, I_i)$ のヘンゼル化と記す。More on Algebra, 補題
\ref{more-algebra-lemma-henselization-colimit} により $A = \colim A_i^h$ である。
More on Algebra, 補題 \ref{more-algebra-lemma-henselization-pair-G-ring} により、環 $A_i^h$
は G-環である。したがって $A = \colim A_i^h$ および
$$
X = \lim (X_i \times_{S_i} \Spec(A_i^h))
$$
がスキームとして成り立つ。$X$ 上有限エタールなスキームの圏は、
$X_i \times_{S_i} \Spec(A_i^h)$ 上有限エタールなスキームの圏の極限である（Limits, 補題
\ref{limits-lemma-descend-finite-presentation},
\ref{limits-lemma-descend-finite-finite-presentation} および
\ref{limits-lemma-descend-etale})
について同様である。$X_0 = \lim (X_i \times_{S_i} S_{i, 0})$ 上有限エタールなスキームにも
同じことが成り立つ。よって、上で扱った
よって、$X / \Spec(A)$ に対する結果を、上で扱った
$(X_i \times_{S_i} \Spec(A_i^h)) / \Spec(A_i^h)$ に対する結果から形式的に導くことができる。
\end{proof}

\begin{lemma}
\label{pione-lemma-finite-etale-invariant-over-proper}
$k'/k$ を代数閉体の拡大とする。$X$ を $k$ 上の固有スキームとする。このとき関手
$$
U \longmapsto U_{k'}
$$
は、$X$ 上有限エタールなスキームの圏と $X_{k'}$ 上有限エタールなスキームの圏との
圏同値である。
\end{lemma}

\begin{proof}
関手が本質的全射であることを示す。$U' \to X_{k'}$ を有限エタール射とする。
$k' = \colim A_i$ を有限型 $k$-代数のフィルター付き余極限と書く。
Limits, 補題 \ref{limits-lemma-descend-finite-presentation} により、ある $i$ と有限表示射
$U_i \to X_{A_i}$ が存在し、その $X_{k'}$ への基底変換は $U'$ である。$i$ を大きく取れば
$U_i \to X_{A_i}$ は有限かつエタールであると仮定できる
（Limits, 補題 \ref{limits-lemma-descend-finite-finite-presentation} および
\ref{limits-lemma-descend-etale}）。
$k$ は代数閉なので、$k$-値点 $t$ を $\Spec(A_i)$ 内に取れる。$U = (U_i)_t$ を
$U_i$ のファイバー（$t$ 上）とする。$A_i^h$ を、$(A_i)_{\mathfrak m}$ のヘンゼル化とする。
ここで $\mathfrak m$ は点 $t$ に対応する極大イデアルである。補題
\ref{pione-lemma-finite-etale-on-proper-over-henselian} により、スキームとして
$(U_i)_{A_i^h} = U \times \Spec(A_i^h)$ であり、$X_{A_i^h}$ 上のスキームである。
ここで $A_i^h$ は $A_i$ 上代数的である（例えば Smoothing Ring Maps, 例
\ref{smoothing-example-describe-henselian} の議論を参照）。さらに $k'$ は代数閉なので、
環写像 $A_i^h \to k'$ を、包含 $A_i \subset k'$ の延長として取れる。
したがって $U'$ は $U$ の基底変換と同型である。忠実充満性の証明もまったく同じである。
\end{proof}








\section{局所連結性}
\label{pione-section-unibranch}

\noindent
この節では、$\pi_1(U) \to \pi_1(X)$ が全射となるのは、$U$ がスキーム $X$ の稠密開集合
であるときである。おおよそ、$U \cap B$ が連結となる場合である。ここで $B$ は
$x \in X \setminus U$ の周りの任意の小さな「球」である。

\begin{lemma}
\label{pione-lemma-dense-faithful}
$f : X \to Y$ をスキームの射とする。$f(X)$ が $Y$ で稠密なら、基底変換関手
$\textit{F\'Et}_Y \to \textit{F\'Et}_X$ は忠実である。
\end{lemma}

\begin{proof}
有限エタール被覆の圏は内部 Hom をもつ（補題
\ref{pione-lemma-internal-hom-finite-etale}）ので、次を示せば十分である。$W$ を $Y$ 上有限エタールな
スキームとし、$s : X \to W$ を $X$ 上の射とする。$t : Y \to W$ を、$s = t \circ f$ を
満たす切断とする。このような切断は高々一つである。$t_1, t_2 : Y \to W$ を、
$s = t_1 \circ f = t_2 \circ f$ を満たす二つの切断とする。
を満たす二つの切断とする。$t_1$ と $t_2$ の等化子は $Y$ で閉じており
（Schemes, 補題 \ref{schemes-lemma-where-are-they-equal}）、$f(X)$ は $Y$ で稠密なので、
$t_1$ と $t_2$ は $Y_{red}$ 上で一致する。したがって $t_1$ と $t_2$ は同じ像をもち、
その像は $W$ の開かつ閉な部分スキームで $Y$ へ同型に写る
（\'Etale Morphisms, 命題 \ref{etale-proposition-properties-sections}）。ゆえに両者は等しい。
\end{proof}

\noindent
次の補題における、厳密ヘンゼル化の穿孔スペクトルが連結であるという条件は、例えば
局所環が幾何学的に一分岐であるという仮定から従う。More on Algebra, 補題
\ref{more-algebra-lemma-geometrically-unibranch} を参照。部分的な逆も
Properties, 補題 \ref{properties-lemma-geometrically-unibranch} にある。

\begin{lemma}
\label{pione-lemma-same-etale-extensions}
$(A, \mathfrak m)$ を局所環とする。$X = \Spec(A)$ とおき、
$U = X \setminus \{\mathfrak m\}$ とする。$A$ の厳密ヘンゼル化の穿孔スペクトルが
連結なら、
$$
\textit{F\'Et}_X \longrightarrow \textit{F\'Et}_U,\quad
Y \longmapsto Y \times_X U
$$
は忠実充満関手である。
\end{lemma}

\begin{proof}
まず $A$ が厳密ヘンゼルであると仮定する。この場合、任意の有限エタール被覆 $Y$ は
$X$ の有限個の互いに素なコピーの和、すなわち $X$ のコピーの和と同型である。したがって、射
$U \to U \amalg \ldots \amalg U$ を、$U$ 上の射として、$X \to X \amalg \ldots \amalg X$
へ $X$ 上で一意的に延長することを示せば十分である。
$U$ が連結（特に空でない）ならこれは成り立つ。

\medskip\noindent
一般の場合。有限エタール被覆の圏が内部 Hom をもつ（補題
\ref{pione-lemma-internal-hom-finite-etale}）ので、次を示せば十分である。$Y$ を $X$ 上有限エタールな
スキームとする。$s : U \to Y$ を $X$ 上の射とし、これを $t : X \to Y$ へ $X$ 上で
延長できることを示す。$A^{sh}$ を $A$ の厳密ヘンゼル化とし、
$X^{sh} = \Spec(A^{sh})$、$U^{sh} = U \times_X X^{sh}$、
$Y^{sh} = Y \times_X X^{sh}$ と記す。第一段落と $A$ に関する仮定により、
$s^{sh} : U^{sh} \to Y^{sh}$ を $s$ の基底変換として $t^{sh} : X^{sh} \to Y^{sh}$ に延長できる。
$A' = A^{sh} \otimes_A A^{sh}$ とおく。$t'_1, t'_2$ を $t^{sh}$ の
$X' = \Spec(A')$ への二つの引き戻しとする。これらは、$s'$ の延長であり、
$s$ の $U' = U \times_X X'$ への引き戻しである。
$A \to A'$ は平坦なので、$U' \subset X'$ は位相的に稠密である。これは
$A \to A'$ に対する降下（Algebra, 補題
\ref{algebra-lemma-flat-going-down}）による。
したがって補題 \ref{pione-lemma-dense-faithful} により $t'_1 = t'_2$ である。ゆえに、例えば
Descent, 補題 \ref{descent-lemma-fpqc-universal-effective-epimorphisms} により、
$t^{sh}$ は射 $t : X \to Y$ に降下する。
\end{proof}

\noindent
補題 \ref{pione-lemma-same-etale-extensions} から、環（およびその厳密ヘンゼル化）の
穿孔スペクトルがいつ連結になるかを知ることは興味深い。十分条件を与える
Hartshorne による有名な補題については、Local Cohomology, 補題
\ref{local-cohomology-lemma-depth-2-connected-punctured-spectrum}.

\begin{lemma}
\label{pione-lemma-quasi-compact-dense-open-connected-at-infinity-Noetherian}
$X$ をスキームとし、$U \subset X$ を稠密開集合とする。次を仮定する。
\begin{enumerate}
\item $X$ の台位相空間はネーター的であり、
\item すべての $x \in X \setminus U$ に対し、$\mathcal{O}_{X, x}$ の
厳密ヘンゼル化の穿孔スペクトルが連結である。
\end{enumerate}
このとき $\textit{F\'Et}_X \to \textit{F\'et}_U$ は忠実充満である。
\end{lemma}

\begin{proof}
 $Y_1, Y_2$ を $X$ 上有限エタールとし、
$\varphi : (Y_1)_U \to (Y_2)_U$ を $U$ 上の射とする。$\varphi$ が
射 $Y_1 \to Y_2$ へ $X$ 上で一意的に持ち上がることを示す。
一意性は補題 \ref{pione-lemma-dense-faithful} から従う。

\medskip\noindent
 $x \in X \setminus U$ を $X \setminus U$ の既約成分の一般点とする。
$V = U \times_X \Spec(\mathcal{O}_{X, x})$ とおく。この選択により、$x$ における
$\Spec(\mathcal{O}_{X, x})$ の穿孔スペクトルである。補題
\ref{pione-lemma-same-etale-extensions} により、射
$\varphi_V : (Y_1)_V \to (Y_2)_V$ を一意的に射
$(Y_1)_{\Spec(\mathcal{O}_{X, x})} \to (Y_2)_{\Spec(\mathcal{O}_{X, x})}$ へ延長できる。
Limits, 補題 \ref{limits-lemma-glueing-near-point} により、開集合
$U \subset U'$（$x$ を含む）と、延長
$\varphi' : (Y_1)_{U'} \to (Y_2)_{U'}$ が $\varphi$ の延長として得られる。$X$ の台位相空間は
ネーター的なので、$\varphi$ が定義されている開集合の補集合に関するネーター帰納法で
証明が完了する。
\end{proof}

\begin{lemma}
\label{pione-lemma-retrocompact-dense-open-connected-at-infinity-closed}
 $X$ をスキームとし、$U \subset X$ を稠密開集合とする。次を仮定する。
\begin{enumerate}
\item $U \to X$ は準コンパクトであり、
\item $X \setminus U$ のすべての点は閉点であり、
\item すべての $x \in X \setminus U$ に対し、$\mathcal{O}_{X, x}$ の
厳密ヘンゼル化の穿孔スペクトルが連結である。
\end{enumerate}
このとき $\textit{F\'Et}_X \to \textit{F\'et}_U$ は忠実充満である。
\end{lemma}

\begin{proof}
 $Y_1, Y_2$ を $X$ 上有限エタールとし、
$\varphi : (Y_1)_U \to (Y_2)_U$ を $U$ 上の射とする。$\varphi$ が
射 $Y_1 \to Y_2$ へ $X$ 上で一意的に持ち上がることを示す。
一意性は補題 \ref{pione-lemma-dense-faithful} から従う。

\medskip\noindent
 $x \in X \setminus U$ とし、$V = U \times_X \Spec(\mathcal{O}_{X, x})$ とおく。
$X \setminus U$ のすべての点が閉点なので、$V$ は
$\Spec(\mathcal{O}_{X, x})$ の穿孔スペクトルである。補題
\ref{pione-lemma-same-etale-extensions} により、射
$\varphi_V : (Y_1)_V \to (Y_2)_V$ を一意的に射
$(Y_1)_{\Spec(\mathcal{O}_{X, x})} \to (Y_2)_{\Spec(\mathcal{O}_{X, x})}$ へ延長できる。
Limits, 補題 \ref{limits-lemma-glueing-near-point}（これは $U$ が $X$ で
レトロコンパクトであることを用いる）により、開集合
$U \subset U'_x$（$x$ を含む）と、延長
$\varphi'_x : (Y_1)_{U'_x} \to (Y_2)_{U'_x}$ が $\varphi$ の延長として得られる。
$x, x' \in X \setminus U$ の二点に対して、射 $\varphi'_x$ と $\varphi'_{x'}$ は
$U'_x \cap U'_{x'}$ 上で、$U$ がその開集合で稠密なので一致する
（補題 \ref{pione-lemma-dense-faithful}）。したがって、望む通り
$\varphi$ を $\bigcup U'_x = X$ まで延長できる。
\end{proof}

\begin{lemma}
\label{pione-lemma-quasi-compact-dense-open-connected-at-infinity}
 $X$ をスキームとし、$U \subset X$ を稠密開集合とする。次を仮定する。
\begin{enumerate}
\item $X$ のすべての準コンパクト開集合は有限個の既約成分をもち、
\item すべての $x \in X \setminus U$ に対し、$\mathcal{O}_{X, x}$ の
厳密ヘンゼル化の穿孔スペクトルが連結である。
\end{enumerate}
このとき $\textit{F\'Et}_X \to \textit{F\'et}_U$ は忠実充満である。
\end{lemma}

\begin{proof}
 $Y_1, Y_2$ を $X$ 上有限エタールとし、
$\varphi : (Y_1)_U \to (Y_2)_U$ を $U$ 上の射とする。$\varphi$ が
射 $Y_1 \to Y_2$ へ $X$ 上で一意的に持ち上がることを示す。一意性は補題
\ref{pione-lemma-dense-faithful} から従う。存在性は、$U$ を拡大できることを、$U \not = X$ の
場合に拡大して示し、ツォルンの補題で証明を完了することで得る。

\medskip\noindent
 $x \in X \setminus U$ を $X \setminus U$ の既約成分の一般点とし、
$V = U \times_X \Spec(\mathcal{O}_{X, x})$ とおく。この選択により、$x$ における
$\Spec(\mathcal{O}_{X, x})$ の穿孔スペクトルである。補題
\ref{pione-lemma-same-etale-extensions} により、射
$\varphi_V : (Y_1)_V \to (Y_2)_V$ を（一意的に）射
$(Y_1)_{\Spec(\mathcal{O}_{X, x})} \to (Y_2)_{\Spec(\mathcal{O}_{X, x})}$ へ延長できる。
$W \subset X$ を $x$ のアフィン近傍として選ぶ。$U \cap W$ は $W$ で稠密なので、一般点
$\eta_1, \ldots, \eta_n$ を含む。これらは $W$ の一般点である。アフィン開集合
$W' \subset W \cap U$ を選び、$\eta_1, \ldots, \eta_n$ を含め、
$V' = W' \times_X \Spec(\mathcal{O}_{X, x})$ とおく。
Limits, 補題 \ref{limits-lemma-glueing-near-point} を
$x \in W \supset W'$ に適用すると、開集合
$W' \subset W'' \subset W$（$x \in W''$）と射
$\varphi'' : (Y_1)_{W''} \to (Y_2)_{W''}$ が得られる。これは
$\varphi$ と $W'$ 上で一致する。$W'$ は
$W'' \cap U$ で稠密なので、補題 \ref{pione-lemma-dense-faithful} により
$\varphi$ と $\varphi''$ は $U \cap W'$ 上で一致する。したがって
$\varphi$ と $\varphi''$ を貼り合わせて射を得る。$\varphi'$ は
$U' = U \cup W''$ 上で、$\varphi$ と $U$ 上で一致する。$x \in U'$ なので、$\varphi$ は
真に大きい開集合へ延長された。

\medskip\noindent
組 $\mathcal{S}$ を、組 $(U', \varphi')$ で $U \subset U'$ かつ $\varphi'$ が
$\varphi$ の延長であるもの全体とする。$\mathcal{S}$ に明らかな方法で半順序を入れる。
$(U'_i, \varphi'_i)$ が全順序部分集合なら、最大元 $(U', \varphi')$ をもつ。
$U' = \bigcup U'_i$ とおき、$\varphi' : (Y_1)_{U'} \to (Y_2)_{U'}$ を
$\varphi'_i$ と $U'_i$ 上で一致する射とする。したがってツォルンの補題が適用でき、
$\mathcal{S}$ は極大元をもつ。上の議論により、この極大元は $\varphi$ の延長であり、
$X$ 全体上で定義される。
\end{proof}

\begin{lemma}
\label{pione-lemma-local-exact-sequence}
$(A, \mathfrak m)$ を局所環とする。$X = \Spec(A)$ とおき、
$U = X \setminus \{\mathfrak m\}$ とする。$U^{sh}$ を、厳密ヘンゼル化 $A^{sh}$（$A$ の）の
穿孔スペクトルと記す。$U$ が準コンパクトで $U^{sh}$ が連結であると仮定する。
このとき列
$$
\pi_1(U^{sh}, \overline{u}) \to \pi_1(U, \overline{u}) \to
\pi_1(X, \overline{u}) \to 1
$$
は補題 \ref{pione-lemma-functoriality-galois-ses} (1) の意味で完全である。
\end{lemma}

\begin{proof}
写像 $\pi_1(U) \to \pi_1(X)$ は補題 \ref{pione-lemma-same-etale-extensions} および
\ref{pione-lemma-functoriality-galois-surjective} により全射である。

\medskip\noindent
$X^{sh} = \Spec(A^{sh})$ と書く。$Y \to X$ を有限エタール射とする。
すると $Y^{sh} = Y \times_X X^{sh} \to X^{sh}$ は有限エタール射である。
$A^{sh}$ は厳密ヘンゼルなので、$Y^{sh}$ は $X^{sh}$ の互いに素なコピーの和と同型である。
したがって $Y \times_X U^{sh}$ についても同じことが成り立つ。よって合成
$\pi_1(U^{sh}) \to \pi_1(U) \to \pi_1(X)$ は自明である。補題
\ref{pione-lemma-composition-trivial} を参照。

\medskip\noindent
証明を完了するには、補題 \ref{pione-lemma-functoriality-galois-ses} により次を示せば十分である。
$V \to U$ を有限エタール射とし、$V \times_U U^{sh}$ が $U^{sh}$ の互いに素なコピーの和であるとする。
このとき有限エタール射 $Y \to X$ を見つけ、$V \cong Y \times_X U$ を $U$ 上で満たすことを示す。
仮定により、有限エタール射 $Y^{sh} \to X^{sh}$ と同型
$V \times_U U^{sh} \cong Y^{sh} \times_{X^{sh}} U^{sh}$ が存在する。次の図式を考える。
$$
\xymatrix{
U \ar[d] & U^{sh} \ar[d] \ar[l] &
U^{sh} \times_U U^{sh} \ar[d] \ar@<1ex>[l] \ar@<-1ex>[l] &
U^{sh} \times_U U^{sh} \times_U U^{sh}
\ar[d] \ar@<1ex>[l] \ar[l] \ar@<-1ex>[l] \\
X & X^{sh} \ar[l] &
X^{sh} \times_X X^{sh} \ar@<1ex>[l] \ar@<-1ex>[l] &
X^{sh} \times_X X^{sh} \times_X X^{sh} \ar@<1ex>[l] \ar[l] \ar@<-1ex>[l]
}
$$
仮定により $U \subset X$ は準コンパクトなので、下向きの矢印はすべて準コンパクトな
開埋め込みである。$\xi \in X^{sh} \times_X X^{sh}$ を
$U^{sh} \times_U U^{sh}$ に属さない点とする。このとき $\xi$ は
閉点 $x^{sh}$ の上にある $X^{sh}$ の点である。局所環準同型
$$
A^{sh} = \mathcal{O}_{X^{sh}, x^{sh}} \to
\mathcal{O}_{X^{sh} \times_X X^{sh}, \xi}
$$
これは第一射影 $X^{sh} \times_X X^{sh}$ で定まる。これは局所準同型の
フィルター付き余極限であり、それらは局所化エタール環写像である。
$A^{sh}$ は厳密ヘンゼルなので、これは同型である。補集合のすべての $\xi$ について
成り立つので、これらの点の間に特殊化は存在せず、したがって各 $\xi$ は閉点である
（これを直接示すこともできる）。$\xi$ における局所環は $A^{sh}$ と同型なので、
厳密ヘンゼルであり、連結な穿孔スペクトルをもつ。同様に、点 $\xi$ について、
$X^{sh} \times_X X^{sh} \times_X X^{sh}$ に属し $U^{sh} \times_U U^{sh} \times_U U^{sh}$ に
属さない場合も同様である。
補題 \ref{pione-lemma-retrocompact-dense-open-connected-at-infinity-closed} により、
鉛直矢印に沿う引き戻しは有限エタールスキームの圏の間の忠実充満関手を誘導する。
したがって、$V \times_U U^{sh}$ 上の標準的降下データは、fpqc 被覆
$\{U^{sh} \to U\}$ に関して、$Y^{sh}$ の降下データへ、fpqc 被覆
$\{X^{sh} \to X\}$ に関して移る。
$Y^{sh} \to X^{sh}$ は有限、したがってアフィンなので、この降下データは有効である
（Descent, 補題 \ref{descent-lemma-affine}）。よって、アフィン射 $Y \to X$ と、
降下データと両立する同型 $Y \times_X X^{sh} \to Y^{sh}$ を得る。
降下データの忠実充満性（Descent, 補題
\ref{descent-lemma-refine-coverings-fully-faithful} の場合と同様）により、
同型 $V \to U \times_X Y$ を得る。最後に、$Y \to X$ は
$Y^{sh} \to X^{sh}$ が有限エタールであることから有限エタールである。Descent, 補題
\ref{descent-lemma-descending-property-etale} および
\ref{descent-lemma-descending-property-finite} を参照。
\end{proof}

\noindent
$X$ を既約スキームとし、$\eta \in X$ を一般点とする。標準射
$\eta \to X$ は標準写像を誘導する。
\begin{equation}
\label{pione-equation-inclusion-generic-point}
\text{Gal}(\kappa(\eta)^{sep}/\kappa(\eta)) = \pi_1(\eta, \overline{\eta})
\longrightarrow \pi_1(X, \overline{\eta})
\end{equation}
左辺の同一視は補題 \ref{pione-lemma-fundamental-group-Galois-group} による。

\begin{lemma}
\label{pione-lemma-irreducible-geometrically-unibranch}
$X$ を既約かつ幾何学的一分岐なスキームとする。任意の空でない開集合
$U \subset X$ に対し、標準写像
$$
\pi_1(U, \overline{u}) \longrightarrow \pi_1(X, \overline{u})
$$
は全射である。写像 (\ref{pione-equation-inclusion-generic-point})
$\pi_1(\eta, \overline{\eta}) \to \pi_1(X, \overline{\eta})$ も全射である。
\end{lemma}

\begin{proof}
補題 \ref{pione-lemma-thickening} により $X$ をその簡約へ置き換えてよい。したがって
$X$ は整スキームであると仮定できる。補題 \ref{pione-lemma-functoriality-galois-surjective} により、
この補題の主張は関手 $\textit{F\'Et}_X \to \textit{F\'Et}_U$ と
$\textit{F\'Et}_X \to \textit{F\'Et}_\eta$ が忠実充満であるという主張に読み替えられる。

\medskip\noindent
 $\textit{F\'Et}_X \to \textit{F\'Et}_U$ に対する結果は、補題
\ref{pione-lemma-quasi-compact-dense-open-connected-at-infinity} と、幾何学的一分岐な局所環
$A$ の厳密ヘンゼル化が既約スペクトルをもつという事実から従う。More on Algebra, 補題
\ref{more-algebra-lemma-geometrically-unibranch} を参照。

\medskip\noindent
剰余体 $\kappa(\eta) = \mathcal{O}_{X, \eta}$ は、$\mathcal{O}_X(U)$ のフィルター付き余極限である。
$U \subset X$ は空でないアフィン開集合を走る。したがって $\textit{F\'Et}_\eta$ は、
$\textit{F\'Et}_U$（そのような $U$ に対する）の圏の余極限である。Limits, 補題
\ref{limits-lemma-descend-finite-presentation},
\ref{limits-lemma-descend-finite-finite-presentation}, および
\ref{limits-lemma-descend-etale} を参照。形式的な議論により、
$\textit{F\'Et}_X \to \textit{F\'Et}_\eta$ の忠実充満性は関手
$\textit{F\'Et}_X \to \textit{F\'Et}_U$ の忠実充満性から従う。
\end{proof}

\begin{lemma}
\label{pione-lemma-exact-sequence-finite-nr-closed-pts}
$X$ をスキームとする。$x_1, \ldots, x_n \in X$ を有限個の閉点とし、次を満たすとする。
\begin{enumerate}
\item $U = X \setminus \{x_1, \ldots, x_n\}$ は連結で $X$ のレトロコンパクト開集合であり、
\item 各 $i$ について、穿孔スペクトル $U_i^{sh}$ は、$\mathcal{O}_{X, x_i}$ の
厳密ヘンゼル化の穿孔スペクトルであり、連結である。
\end{enumerate}
このとき写像 $\pi_1(U) \to \pi_1(X)$ は全射であり、その核は、
$\pi_1(U)$ の、$\pi_1(U_i^{sh}) \to \pi_1(U)$ の像（$i = 1, \ldots, n$）を含む
最小の閉正規部分群である。
\end{lemma}

\begin{proof}
全射性は補題 \ref{pione-lemma-retrocompact-dense-open-connected-at-infinity-closed} および
\ref{pione-lemma-functoriality-galois-surjective} から従う。写像の列
を考える。
$$
\pi_1(U)  \to \ldots \to
\pi_1(X \setminus \{x_1, x_2\}) \to \pi_1(X \setminus \{x_1\}) \to \pi_1(X)
$$
群論的議論により、核について $n = 1$ の場合を示せば十分である（詳細は省略）。
$x = x_1$、$U^{sh} = U_1^{sh}$ と書き、$A = \mathcal{O}_{X, x}$ とおき、
$A^{sh}$ を厳密ヘンゼル化とする。次の図式を考える。
$$
\xymatrix{
U \ar[d] &
\Spec(A) \setminus \{\mathfrak m\} \ar[l] \ar[d] &
U^{sh} \ar[d] \ar[l] \\
X & \Spec(A) \ar[l] & \Spec(A^{sh}) \ar[l]
}
$$
補題 \ref{pione-lemma-functoriality-galois-ses} により、$V \to U$ が $U^{sh}$ の自明な被覆へ
引き戻される有限エタール射を、$X$ 上の有限エタールスキームへ延長できることを示す。
補題 \ref{pione-lemma-local-exact-sequence} により、$A$ の穿孔スペクトル上の有限エタールスキーム
について対応する主張を知っている。しかし Limits, 補題
\ref{limits-lemma-glueing-near-closed-point} により、$X$ 上有限表示なスキームは、
$U$ 上有限表示なスキームと $A$ 上有限表示なスキームを $A$ の穿孔スペクトル上で
貼り合わせたものと同じである。これで証明が完了する。
\end{proof}









\section{正規スキームの基本群}
\label{pione-section-normal}

\noindent
整で幾何学的一分岐なスキーム $X$ をとる。前節で、$X$ の基本群は
関数体 $K$ が $X$ のものであるガロア群の商であることを見た。写像は連続なので、
その核はガロア群の正規閉部分群である。したがってこの核はガロア理論
(Fields, 定理 \ref{fields-theorem-inifinite-galois-theory}) により
ガロア拡大 $M/K$ に対応する。本節では、決定すべき $M$ を
$X$ が正規整スキームである場合に求める。

\medskip\noindent
整正規スキーム $X$ をとり、その関数体を $K$ とする。$L/K$ を有限拡大とする。
Morphisms, 節 \ref{morphisms-section-normalization-X-in-Y} で定義した
正規化 $Y \to X$ を、$X$ における射 $\Spec(L) \to X$ に関して考える。
この設定で、{\it $X$ は $L$ において非分岐である} とは、
$Y \to X$ がスキームの非分岐射であることを言う。この条件は補題
\ref{pione-lemma-unramified} で明らかにする。$Y \to X$ の
$\Spec(K)$ 上のスキーム論的ファイバーは $\Spec(L)$ であることに注意する。
したがって体拡大 $L/K$ は、$X$ が $L$ において非分岐なら分離的である。
Morphisms, 補題 \ref{morphisms-lemma-unramified-over-field} を参照。

\begin{lemma}
\label{pione-lemma-unramified-in-L}
上記の状況では次は同値である。
\begin{enumerate}
\item $X$ は $L$ において非分岐である。
\item $Y \to X$ はエタールである。
\item $Y \to X$ は有限エタールである。
\end{enumerate}
\end{lemma}

\begin{proof}
 $Y \to X$ は整射であることに注意する。いずれの場合も、定義により
射 $Y \to X$ は局所有限型である。したがって各場合に
Morphisms, 補題 \ref{morphisms-lemma-finite-integral} により
$Y \to X$ は有限である。特に (2) と (3) は同値である。
エタール射は非分岐なので、(2) は (1) を含意する。

\medskip\noindent
逆に $Y \to X$ が非分岐だと仮定する。正規スキームは幾何学的一分岐である
(Properties, 補題 \ref{properties-lemma-normal-geometrically-unibranch}) から、
More on Morphisms, 補題
\ref{more-morphisms-lemma-global-unramified-dominant-unibranch-is-etale} により
射 $Y \to X$ はエタールである。次の段落では直接証明も与える。

\medskip\noindent
 $x \in X$ とする。次を満たすエタール近傍 $(U, u) \to (X, x)$ を選べる。
$$
Y \times_X U = \coprod V_j \longrightarrow U
$$
は互いに素な閉埋め込みの直和となる。\'Etale Morphisms, 補題
\ref{etale-lemma-finite-unramified-etale-local} を参照。縮小して $U$ を
準コンパクトとしてよい。このとき $U$ は有限個の既約成分をもつ
(Descent, 補題 \ref{descent-lemma-locally-finite-nr-irred-local-fppf})。
$U$ は正規なので (Descent, 補題
\ref{descent-lemma-normal-local-smooth})、$U$ の既約成分は開かつ閉である
(Properties, 補題 \ref{properties-lemma-normal-locally-finite-nr-irreducibles})。
したがって $U$ は既約だと仮定してよい。このとき $U$ は整スキームであり、
その生成点 $\xi$ は $X$ の生成点へ写る。一方、
More on Morphisms, 補題
\ref{more-morphisms-lemma-normalization-smooth-localization} により
$Y \times_X U$ は $U$ の正規化であり、$\Spec(L) \times_X U$ において考えることが
分かる。$\Spec(L) \times_X U$ のすべての点は $\xi$ へ写る。よって
Morphisms, 補題 \ref{morphisms-lemma-normalization-generic} により各 $V_j$ は
$\xi$ へ写る点を含む。よって $V_j \to U$ は、$U = \overline{\{\xi\}}$ なので
同型である。したがって $Y \times_X U \to U$ はエタールである。
Descent, 補題 \ref{descent-lemma-descending-property-etale} により、
$Y \to X$ は $U \to X$ の像（$x$ の開近傍）上でエタールであると結論する。
\end{proof}

\begin{lemma}
\label{pione-lemma-finite-etale-covering-normal-unramified}
$X$ を関数体 $K$ をもつ正規整スキームとする。$Y \to X$ を有限エタール射とする。
$Y$ が連結なら、$Y$ は整な正規スキームであり、$Y$ は
$X$ の正規化（その体は $Y$ の関数体における）である。
\end{lemma}

\begin{proof}
Descent, 補題 \ref{descent-lemma-normal-local-smooth} によりスキーム $Y$ は正規である。
$Y \to X$ は平坦なので、Morphisms, 補題
\ref{morphisms-lemma-generalizations-lift-flat} により $Y$ の各生成点は
$X$ の生成点へ写る。$Y \to X$ は有限なので、$Y$ は有限個の既約成分をもつ。
したがって Properties, 補題
\ref{properties-lemma-normal-locally-finite-nr-irreducibles} により $Y$ は
有限個の整な正規スキームの互いに素な直和である。ゆえに $Y$ が連結なら、
$Y$ は整な正規スキームである。

\medskip\noindent
$L$ を $Y$ の関数体とし、$Y' \to X$ を $X$ の $L$ における正規化とする。
Morphisms, 補題 \ref{morphisms-lemma-characterize-normalization} により
$Y' \to Y \to X$ という分解を得て、$Y' \to Y$ は $Y$ の $L$ における正規化である。
$Y$ は正規なので $Y' = Y$ であることは明らかである（これは
Morphisms, 補題 \ref{morphisms-lemma-finite-birational-over-normal} からも
導ける）。
\end{proof}

\begin{proposition}
\label{pione-proposition-normal}
$X$ を関数体 $K$ をもつ正規整スキームとする。このとき標準写像
(\ref{pione-equation-inclusion-generic-point})
$$
\text{Gal}(K^{sep}/K) = \pi_1(\eta, \overline{\eta})
\longrightarrow \pi_1(X, \overline{\eta})
$$
は商写像
$\text{Gal}(K^{sep}/K) \to \text{Gal}(M/K)$ と同一視される。ここで
$M \subset K^{sep}$ は有限部分拡大 $L$ の合併であり、$X$ は $L$ ごとに
非分岐である。
\end{proposition}

\begin{proof}
$X$ は正規スキームなので幾何学的一分岐である
(Properties, 補題 \ref{properties-lemma-normal-geometrically-unibranch})。
したがって補題 \ref{pione-lemma-irreducible-geometrically-unibranch} を $X$ に適用できる。
ゆえに $\pi_1(\eta, \overline{\eta}) \to \pi_1(X, \overline{\eta})$ は全射であり、
次の可換図式の上側水平矢印は忠実充満である。
$$
\xymatrix{
\textit{F\'Et}_X \ar[r] \ar[d] \ar[rd]_c & \textit{F\'Et}_\eta \ar[d] \\
\textit{Finite-}\pi_1(X, \overline{\eta})\textit{-sets} \ar[r] &
\textit{Finite-}\text{Gal}(K^{sep}/K)\textit{-sets}
}
$$
左側の垂直矢印は定理 \ref{pione-theorem-fundamental-group} の同値であり、
右側の垂直矢印は補題 \ref{pione-lemma-fundamental-group-Galois-group} の同値である。
下側の水平矢印は命題の写像から誘導される。補題
\ref{pione-lemma-unramified-in-L} と \ref{pione-lemma-finite-etale-covering-normal-unramified}
により、$c$ の本質像は次の形の集合と同型な
$\text{Gal}(K^{sep}/K)\textit{-Sets}$ からなることが分かる。
$$
S = \Hom_K(\prod\nolimits_{i = 1, \ldots, n} L_i, K^{sep}) =
\coprod\nolimits_{i = 1, \ldots, n} \Hom_K(L_i, K^{sep})
$$
ここで $L_i/K$ は有限分離拡大で、$X$ は $L_i$ において非分岐である。
したがって、補題の主張にあるように $M \subset K^{sep}$ なら、
$\text{Gal}(K^{sep}/M)$ は、$\text{Gal}(K^{sep}/K)$ の部分群であり、
$c$ の本質像にあるすべての対象に自明に作用する。一方、$c$ の本質像は、
$S$ のうち $\text{Gal}(K^{sep}/K)$ の作用が
$\text{Gal}(K^{sep}/K) \to \pi_1(X, \overline{\eta})$ の全射を経由するものからなる。
したがって
$\text{Gal}(K^{sep}/M)$ は核である。ゆえに
$\text{Gal}(K^{sep}/M)$ は正規部分群、$M/K$ はガロア拡大であり、次の短完全列を得る。
$$
1 \to \text{Gal}(K^{sep}/M) \to
\text{Gal}(K^{sep}/K) \to
\text{Gal}(M/K) \to 1
$$
ガロア理論 (Fields, 定理
\ref{fields-theorem-inifinite-galois-theory} および
補題 \ref{fields-lemma-ses-infinite-galois}) による。証明を終える。
\end{proof}

\begin{lemma}
\label{pione-lemma-local-exact-sequence-normal}
$(A, \mathfrak m)$ を正規局所環とする。$X = \Spec(A)$ とおく。
$A^{sh}$ を $A$ の厳密ヘンゼル化とする。$K$ と $K^{sh}$ をそれぞれ
$A$ と $A^{sh}$ の分数体とする。このとき次の列
$$
\pi_1(\Spec(K^{sh})) \to \pi_1(\Spec(K)) \to \pi_1(X) \to 1
$$
は補題 \ref{pione-lemma-functoriality-galois-ses} (1) の意味で完全である。
\end{lemma}

\begin{proof}
 $A^{sh}$ は正規整域であることに注意する。More on Algebra, 補題
\ref{more-algebra-lemma-henselization-normal} を参照。
$\pi_1(\Spec(K)) \to \pi_1(X)$ は命題
\ref{pione-proposition-normal} により全射である。

\medskip\noindent
$X^{sh} = \Spec(A^{sh})$ と書く。$Y \to X$ を有限エタール射とする。
すると $Y^{sh} = Y \times_X X^{sh} \to X^{sh}$ は有限エタール射である。
$A^{sh}$ は厳密ヘンゼルなので、$Y^{sh}$ は $X^{sh}$ のコピーの
互いに素な直和と同型である。したがって $Y \times_X \Spec(K^{sh})$ も
同様である。ゆえに合成 $\pi_1(\Spec(K^{sh})) \to \pi_1(X)$ は自明である。
補題 \ref{pione-lemma-composition-trivial} を参照。

\medskip\noindent
証明を終えるには、補題 \ref{pione-lemma-functoriality-galois-ses} に従い、次を
示せば十分である。$V \to \Spec(K)$ を有限エタール射で、
$V \times_{\Spec(K)} \Spec(K^{sh})$ が $\Spec(K^{sh})$ のコピーの
互いに素な直和であるとする。このとき
$Y \to X$ を満たし、$V \cong Y \times_X \Spec(K)$（$\Spec(K)$ 上）となる
有限エタール射を見つけられることを示す。$V = \Spec(L)$ と書くと、
$L$ は $K$ の有限分離拡大の有限積である。$B \subset L$ を $A$ の $L$ における
整閉包とする。$A \to B$ がエタールなら $Y = \Spec(B)$ ととれて
証明は完了する。Algebra, 補題
\ref{algebra-lemma-integral-closure-commutes-smooth}（省略した極限論法とともに）
により、$B \otimes_A A^{sh}$ は
$A^{sh}$ の、$L^{sh} = L \otimes_K K^{sh}$ における整閉包である。
仮定は $L^{sh}$ が $K^{sh}$ のコピーの積であるということなので、
$B^{sh}$ は $A^{sh}$ のコピーの積である。したがって
$A^{sh} \to B^{sh}$ はエタールである。$A \to A^{sh}$ は忠実平坦なので、
Descent, 補題 \ref{descent-lemma-descending-property-etale} により
$A \to B$ はエタールとなり、所望の結論を得る。
\end{proof}







\section{群作用と整閉包}
\label{pione-section-group-actions-integral}

\noindent
本節では More on Algebra, 節 \ref{more-algebra-section-group-actions-integral} の議論を続ける。
正規局所環は定義により整域であることを思い出そう。

\begin{lemma}
\label{pione-lemma-get-algebraic-closure}
正規整域 $A$ で、その分数体 $K$ が分離代数閉であるものをとる。
$\mathfrak p \subset A$ を非零素イデアルとする。このとき剰余体
$\kappa(\mathfrak p)$ は代数閉である。
\end{lemma}

\begin{proof}
補題が成り立たないと仮定して矛盾を導く。すると
単首既約多項式 $P(T) \in \kappa(\mathfrak p)[T]$ で次数 $d > 1$ のものが存在する。
分母を払うため、$P$ を $a^d P(a^{-1}T)$ に置き換える
（ここで $a \in A$ は適切に選ぶ）と、$P$ は単首多項式
$Q$（$A[T]$ における）の像であると仮定できる。$Q$ は $K[T]$ で
既約であることに注意する。
実際、$K$ 上の因数分解は Algebra, 補題
\ref{algebra-lemma-polynomials-divide} により $A$ 上の因数分解を導き、
$\mathfrak p$ で剰余をとれば $P$ の因数分解を得てしまう。
$K$ は分離代数閉なので、$Q$ は分離多項式ではない
(Fields, 定義 \ref{fields-definition-separable})。したがって $K$ の標数は
$p > 0$ であり、$Q$ の一次項は消える
(Fields, 定義 \ref{fields-definition-separable})。しかしこのとき
しかし $Q$ を $Q + a T$（$a \in \mathfrak p$ は非零）へ置き換えれば矛盾となる。
\end{proof}

\begin{lemma}
\label{pione-lemma-normal-local-domain-separablly-closed-fraction-field}
分離閉な分数体をもつ正規局所環は厳密ヘンゼルである。
\end{lemma}

\begin{proof}
$(A, \mathfrak m, \kappa)$ を分離閉な分数体 $K$ をもつ正規局所環とする。
$A = K$ なら終わりである。そうでなければ、補題
\ref{pione-lemma-get-algebraic-closure} により剰余体 $\kappa$ は代数閉であり、
$A$ がヘンゼルであることを確認すれば十分である。
$f \in A[T]$ を単首とし、$a_0 \in \kappa$ を
重複度 $1$ の、還元 $\overline{f} \in \kappa[T]$ における根とする。
$f = \prod f_i$ を $K[T]$ における因数分解とする。
Algebra, 補題 \ref{algebra-lemma-polynomials-divide} により
$f_i \in A[T]$ である。したがって $a_0$ は $f_i$ の根であり、ある $i$ が存在する。
$f$ を $f_i$ に置き換えて $f$ が既約だと仮定してよい。
導関数 $f'$ は $A[T]$ で零ではないので、$a_0$ が単根であることから、
$f$ は $K$ が分離代数閉であることにより一次であると結論する。
したがって $A$ はヘンゼルである。Algebra, 定義
\ref{algebra-definition-henselian} を参照。
\end{proof}

\begin{lemma}
\label{pione-lemma-inertia-base-change}
有限群 $G$ が環 $R$ に作用するとする。$R^G \to A$ を環準同型とする。
$\mathfrak q' \subset A \otimes_{R^G} R$ を、素イデアル
$\mathfrak q \subset R$ の上にある素イデアルとする。このとき
$$
I_\mathfrak q = \{\sigma \in G \mid
\sigma(\mathfrak q) = \mathfrak q\text{ かつ }
\sigma \bmod \mathfrak q = \text{id}_{\kappa(\mathfrak q)}\}
$$
は
$$
I_{\mathfrak q'} = \{\sigma \in G \mid
\sigma(\mathfrak q') = \mathfrak q'\text{ かつ }
\sigma \bmod \mathfrak q' = \text{id}_{\kappa(\mathfrak q')}\}
$$
に等しい。
\end{lemma}

\begin{proof}
$\mathfrak q$ は $\mathfrak q'$ の逆像であり、
$\kappa(\mathfrak q) \subset \kappa(\mathfrak q')$ なので、
$I_{\mathfrak q'} \subset I_\mathfrak q$ を得る。
逆に $\sigma \in I_\mathfrak q$ なら、$\sigma$ はファイバー環
$A \otimes_{R^G} \kappa(\mathfrak q)$ 上自明に作用する。
したがって $\sigma$ は $\mathfrak q$ の上にあるすべての素イデアルを固定し、
それらの剰余体上で恒等写像を誘導する。
\end{proof}

\begin{lemma}
\label{pione-lemma-inertia-invariants-etale}
有限群 $G$ が環 $R$ に作用するとする。$\mathfrak q \subset R$ を素イデアルとし、
次のようにおく。
$$
I = \{\sigma \in G \mid \sigma(\mathfrak q) = \mathfrak q
\text{ かつ } \sigma \bmod \mathfrak q = \text{id}_\mathfrak q\}
$$
すると $R^G \to R^I$ は $R^I \cap \mathfrak q$ においてエタールである。
\end{lemma}

\begin{proof}
証明の方針はエタール局所化を用い、$R \to R^I$ が
$R^I \cap \mathfrak p$ において局所同型となる場合へ帰着することである。
$R^G \to A$ をエタール環準同型とする。群 $G$ の
$A \otimes_{R^G} R$ への作用と、素イデアル
$\mathfrak q'$（$A \otimes_{R^G} R$ の素イデアル）が
$\mathfrak q$ の上にある場合について結果が成り立つなら、
元の結果も成り立つと主張する。

\medskip\noindent
これを確認する。$R^G \to A$ は平坦なので
$A = (A \otimes_{R^G} R)^G$ である。More on Algebra,
補題 \ref{more-algebra-lemma-base-change-invariants} を参照。
補題 \ref{pione-lemma-inertia-base-change} により群 $I$ は変わらない。
次に More on Algebra, 補題
\ref{more-algebra-lemma-base-change-invariants} をもう一度適用すると、
$A \otimes_{R^G} R^I = (A \otimes_{R^G} R)^I$ を得る
（$R^I \to A \otimes_{R^G} R^I$ は平坦だからである）。
したがって
$$
\xymatrix{
\Spec((A \otimes_{R^G} R)^I) \ar[d] \ar[r] & \Spec(R^I) \ar[d] \\
\Spec(A) \ar[r] & \Spec(R^G)
}
$$
これはデカルトであり、水平矢印はエタールである。したがって左側の
垂直矢印が、ある開近傍 $W$（$(A \otimes_{R^G} R)^I \cap \mathfrak q'$ を含む）
でエタールなら、右側の垂直矢印は $W$ の $\Spec(R^I)$ における
（開な）像の点でエタールである。Descent, 補題
\ref{descent-lemma-smooth-permanence} を参照。特に射
$\Spec(R^I) \to \Spec(R^G)$ は $R^I \cap \mathfrak q$ において
エタールである。

\medskip\noindent
$\mathfrak p = R^G \cap \mathfrak q$ とおく。
More on Algebra, 補題 \ref{more-algebra-lemma-one-orbit} により、
$\Spec(R) \to \Spec(R^G)$ の $\mathfrak p$ 上のファイバーは有限である。
さらに、これらの点での剰余体拡大は代数的かつ正規で、自己同型群が有限である。
これは More on Algebra, 補題
\ref{more-algebra-lemma-one-orbit-geometric} による。したがって
More on Morphisms, 補題
\ref{more-morphisms-lemma-etale-makes-integral-split} を、整環準同型
$R^G \to R$ と素イデアル $\mathfrak p$ に適用できる。
上の主張と合わせると、$R = A_1 \times \ldots \times A_n$ に帰着する。
各 $A_i$ はただ一つの素イデアル $\mathfrak q_i$ をもち、これは
$\mathfrak p$ の上にあり、剰余体拡大
$\kappa(\mathfrak q_i)/\kappa(\mathfrak p)$ は純非分離である。もちろん
$\mathfrak q$ はこれらの素イデアルの一つであり、$\mathfrak q = \mathfrak q_1$ としてよい。

\medskip\noindent
$G$ が因子 $A_i$ を置換するとは限らない（例えば $A_i$ のスペクトルが連結で、
$R^G$ が局所環ならこれは成り立つ）。これは次のように直せる。
読者自身で考え抜くことを勧めるが、位相の代わりに冪等元を使ってもよい。
積分解は $\Spec(R)$ の開かつ閉な部分集合 $U_i$ による互いに素な
直和分解を与えることを思い出そう。$G$ は有限なので、この被覆を
$\Spec(R) = \coprod_{j \in J} W_j$ による開かつ閉な $W_j$ の有限直和分解に
細分でき、各 $j \in J$ について
$j' \in J$ が存在して $\sigma(W_j) = W_{j'}$ となるようにできる。
各 $W_j$ のうち $\{\mathfrak q_1, \ldots, \mathfrak q_n\}$ と交わらないものの合併は、
$\mathfrak p$ 上のファイバーと交わらない閉集合である。したがって
$\Spec(R^G)$ の、$\mathfrak p$ と交わらない閉集合へ、
$\Spec(R) \to \Spec(R^G)$ が閉写像であることにより写る。
ゆえに $R^G$ を主局所化で置き換えれば（上の主張により許される）、
各 $W_j$ がいずれかの $\mathfrak q_i$ と交わると仮定できる。
そこで $U_i = W_j$ とおく。ただし $\mathfrak q_i \in W_j$ のときである。
対応する積分解 $R = A_1 \times \ldots \times A_n$ は
$G$ が因子 $A_i$ を置換するものとなる。

\medskip\noindent
したがって $R = A_1 \times \ldots \times A_n$ という、$G$-作用と両立する積分解があり、
各 $A_i$ はただ一つの素イデアル $\mathfrak q_i$ をもち、これは
$\mathfrak p$ の上にあり、
体拡大 $\kappa(\mathfrak q_i)/\kappa(\mathfrak p)$ は純非分離であると仮定できる。
$A' = A_2 \times \ldots \times A_n$ と書くと、
$$
R = A_1 \times A'
$$
となる。$\mathfrak q = \mathfrak q_1$ なので、すべての
$\sigma \in I$ は上の積分解を保つ。したがって
$$
R^I = (A_1)^I \times (A')^I
$$
$I = D = \{\sigma \in G \mid \sigma(\mathfrak q) = \mathfrak q\}$ であることに注意する。
これは $\kappa(\mathfrak q)/\kappa(\mathfrak p)$ が純非分離だからである。
$G$ の $\mathfrak p$ 上の素イデアルへの作用は推移的である
(More on Algebra, 補題 \ref{more-algebra-lemma-one-orbit})。したがって
$I$ の $G$ における指数は $n$ であり、
$G = eI \amalg \sigma_2I \amalg \ldots \amalg \sigma_nI$ と書け、
$A_i = \sigma_i(A_1)$（$i = 2, \ldots, n$）である。ゆえに
$$
R^G = (A_1)^I.
$$
写像 $R^G \to R^I$ は $R^I \cap \mathfrak q$ でエタールであり、証明は完了する。
\end{proof}

\noindent
次の補題は More on Algebra, 補題
\ref{more-algebra-lemma-inertial-invariants-unramified} を一般化する。

\begin{lemma}
\label{pione-lemma-inertial-invariants-unramified}
$A$ を分数体 $K$ をもつ正規整域とする。$L/K$ を（無限でもよい）
ガロア拡大とする。$G = \text{Gal}(L/K)$ とおき、$B$ を
$A$ の $L$ における整閉包とする。$\mathfrak q \subset B$ とする。次のようにおく。
$$
I = \{\sigma \in G \mid
\sigma(\mathfrak q) = \mathfrak q \text{ かつ }
\sigma \bmod \mathfrak q = \text{id}_{\kappa(\mathfrak q)}\}
$$
すると $(B^I)_{B^I \cap \mathfrak q}$ はエタールな $A$-代数の
フィルター付き余極限である。
\end{lemma}

\begin{proof}
$L$ を $K$ の有限ガロア拡大のフィルター付き余極限として書ける。
したがって $L/K$ が有限ガロア拡大の場合にこの補題を証明すれば十分である。
Algebra, 補題 \ref{algebra-lemma-colimit-colimit-etale} を参照。
$A = B^G$ である。$A$ は $K = L^G$ において整閉なので、結果は
補題 \ref{pione-lemma-inertia-invariants-etale} から従う。
\end{proof}







\section{分岐理論}
\label{pione-section-ramification}

\noindent
本節では More on Algebra, 節 \ref{more-algebra-section-ramification} の議論を続け、
スキームの基本群についての議論と関連付ける。

\medskip\noindent
$(A, \mathfrak m, \kappa)$ を分数体 $K$ をもつ正規局所環とする。
分離代数閉包 $K^{sep}$ をとる。$A^{sep}$ を $A$ の $K^{sep}$ における整閉包とする。
極大イデアル $\mathfrak m^{sep} \subset A^{sep}$ を一つとる。
$A \subset A^h \subset A^{sh}$ をヘンゼル化および厳密ヘンゼル化とする。
$A^h$ と $A^{sh}$ も正規環であることに注意する
(More on Algebra, 補題 \ref{more-algebra-lemma-henselization-normal})。
それらの分数体をそれぞれ $K^h$ と $K^{sh}$ と書く。
補題 \ref{pione-lemma-normal-local-domain-separablly-closed-fraction-field} により
$(A^{sep})_{\mathfrak m^{sep}}$ は厳密ヘンゼルなので、
$A$-代数準同型 $A^{sh} \to (A^{sep})_{\mathfrak m^{sep}}$ をとることができる。
具体的には、まず $\kappa$-埋め込み\footnote{これは、
$\kappa(\mathfrak m^{sh})$ が $\kappa$ の分離代数閉包であり、
$\kappa(\mathfrak m^{sep})$ が補題 \ref{pione-lemma-get-algebraic-closure} により
$\kappa$ の代数閉包であることから可能である。}
$\kappa(\mathfrak m^{sh}) \to \kappa(\mathfrak m^{sep})$ を選び、
Algebra, 補題 \ref{algebra-lemma-strictly-henselian-functorial} により
$A$-代数準同型へ（ただ一通りに）拡張する。
次の図式を得る。
$$
\xymatrix{
K^{sep} & K^{sh} \ar[l] & K^h \ar[l] & K \ar[l] \\
(A^{sep})_{\mathfrak m^{sep}} \ar[u] &
A^{sh} \ar[u] \ar[l] &
A^h \ar[u] \ar[l] &
A \ar[u] \ar[l]
}
$$
これらの環のスペクトルの基本群を考えることができる。
もちろん $K^{sep}$、$(A^{sep})_{\mathfrak m^{sep}}$、$A^{sh}$ は厳密ヘンゼルなので、
それらについては自明な群を得る。したがって興味深い部分は次である。
\begin{equation}
\label{pione-equation-inertia-diagram-pione}
\vcenter{
\xymatrix{
\pi_1(U^{sh}) \ar[r] \ar[rd]_1 & \pi_1(U^h) \ar[d] \ar[r] & \pi_1(U) \ar[d] \\
& \pi_1(X^h) \ar[r] & \pi_1(X)
}
}
\end{equation}
ここで $X^h$ と $X$ はそれぞれ $A^h$ と $A$ のスペクトルであり、
$U^{sh}$、$U^h$、$U$ はそれぞれ $K^{sh}$、$K^h$、$K$ のスペクトルである。
ラベル $1$ は写像が自明であることを意味する。これは自明群 $\pi_1(X^{sh})$ を経由するためである。
一方、プロ有限群 $G = \text{Gal}(K^{sep}/K)$ は $A^{sep}$ に作用するので、次を定義する。
$$
D = \{\sigma \in G \mid \sigma(\mathfrak m^{sep}) = \mathfrak m^{sep}\}
\supset
I = \{\sigma \in D \mid \sigma \bmod \mathfrak m^{sep} =
\text{id}_{\kappa(\mathfrak m^{sep})}\}
$$
これらの群は、特に $A$ が離散付値環のとき、
それぞれ {\it 分解群} および {\it 慣性群} と呼ばれることがある。

\begin{lemma}
\label{pione-lemma-identify-inertia}
上記の状況では、同型
$\pi_1(U) = \text{Gal}(K^{sep}/K)$ を介して、図式
(\ref{pione-equation-inertia-diagram-pione}) は次の図式に移る。
$$
\xymatrix{
I \ar[r] \ar[rd]_1 & D \ar[d] \ar[r] & \text{Gal}(K^{sep}/K) \ar[d] \\
& \text{Gal}(\kappa(\mathfrak m^{sh})/\kappa) \ar[r] & \text{Gal}(M/K)
}
$$
ここで $K^{sep}/M/K$ は $A$ に関して非分岐な最大部分拡大である。
さらに、垂直矢印は全射であり、左の垂直矢印の核は $I$、
右の垂直矢印の核は $\text{Gal}(K^{sep}/K)$ の $I$ を含む最小閉正規部分群である。
\end{lemma}

\begin{proof}
構成により、群 $D$ は $(A^{sep})_{\mathfrak m^{sep}}$ に $A$ 上で作用する。
剰余体上の写像を指定したときの $A^{sh} \to (A^{sep})_{\mathfrak m^{sep}}$ の一意性
(Algebra, 補題 \ref{algebra-lemma-strictly-henselian-functorial}) から、
$A^{sh} \to (A^{sep})_{\mathfrak m^{sep}}$ の像は
$((A^{sep})_{\mathfrak m^{sep}})^I$ に含まれる。
一方、補題 \ref{pione-lemma-inertial-invariants-unramified} により
$((A^{sep})_{\mathfrak m^{sep}})^I$ は $A$ のエタール拡大のフィルター付き余極限である。
$A^{sh}$ はそのような最大の拡大だから、
$A^{sh} = ((A^{sep})_{\mathfrak m^{sep}})^I$ と結論する。
したがって $K^{sh} = (K^{sep})^I$ である。

\medskip\noindent
$I$ は全射
$D \to \text{Aut}(\kappa(\mathfrak m^{sep})/\kappa)$ の核であることを思い出そう。
More on Algebra, 補題 \ref{more-algebra-lemma-one-orbit-geometric-galois} を参照。
上で見たように、
$\text{Aut}(\kappa(\mathfrak m^{sep})/\kappa) =
\text{Gal}(\kappa(\mathfrak m^{sh})/\kappa)$ であり、これらの体は
$\kappa$ の代数閉包および分離代数閉包である。
一方、Algebra, 補題 \ref{algebra-lemma-henselian-functorial} の一意性により、
任意の $A^{sh}$ の $A$ 上の自己同型は $A^{sh}$ の $A^h$ 上の自己同型である。
さらに $A^{sh}$ は有限エタール拡大 $A^h \subset A'$ の余極限であり、
これらは $1$ 対 $1$ で有限分離拡大 $\kappa'/\kappa$ と対応する。
Algebra, 注意 \ref{algebra-remark-construct-sh-from-h} を参照。したがって
$$
\text{Aut}(A^{sh}/A) = \text{Aut}(A^{sh}/A^h) =
\text{Gal}(\kappa(\mathfrak m^{sh})/\kappa)
$$
$\kappa'/\kappa$ をガロア群 $G$ をもつ有限ガロア拡大とする。
Algebra, 補題 \ref{algebra-lemma-henselian-cat-finite-etale} により
有限エタール拡大 $A^h \subset A'$ をとる。これは $\kappa \subset \kappa'$ に対応する。
分数体と次数を調べれば $(A')^G = A^h$ が従う
（細部は省略する）。余極限をとると
$(A^{sh})^{\text{Gal}(\kappa(\mathfrak m^{sh})/\kappa)} = A^h$ と結論できる。
以上を合わせて $A^h = ((A^{sep})_{\mathfrak m^{sep}})^D$ を得る。
したがって $K^h = (K^{sep})^D$ である。

\medskip\noindent
$U$、$U^h$、$U^{sh}$ はそれぞれ体 $K$、$K^h$、$K^{sh}$ のスペクトルなので、
補題 \ref{pione-lemma-fundamental-group-Galois-group} により二つの図式の上段が対応する。
補題 \ref{pione-lemma-gabber} により
$\pi_1(X^h) = \text{Gal}(\kappa(\mathfrak m^{sh})/\kappa)$ である。
列
$1 \to I \to D \to \text{Gal}(\kappa(\mathfrak m^{sh})/\kappa) \to 1$
が完全であることは上で指摘した。
命題 \ref{pione-proposition-normal} により
$\pi_1(X) = \text{Gal}(M/K)$ である。
最後に、写像
$\text{Gal}(K^{sep}/K) \to \text{Gal}(M/K) = \pi_1(X)$ の核に関する主張は、
補題 \ref{pione-lemma-local-exact-sequence-normal} から従う。これで証明は完了する。
\end{proof}

\noindent
$X$ を関数体 $K$ をもつ正規整スキームとする。
分離代数閉包 $K^{sep}$ を $K$ に対してとる。
$X^{sep} \to X$ を $X$ の $K^{sep}$ における正規化とする。
$G = \text{Gal}(K^{sep}/K)$ は $K^{sep}$ に作用するので、
$G$ の右作用を $X^{sep}$ 上に得る。$y \in X^{sep}$ に対して次を定義する。
$$
D_y = \{\sigma \in G \mid \sigma(y) = y\} \supset
I_y = \{\sigma \in D \mid \sigma \bmod \mathfrak m_y =
\text{id}_{\kappa(y)} \}
$$
上と同様である。一方、$x \in X$ に対し、
$\mathcal{O}_{X, x}^{sh}$ を厳密ヘンゼル化とし、
その分数体を $K_x^{sh}$ とし、$\mathcal{O}_{X, x}^{sh}$ を明示して、
$K$-埋め込み $K_x^{sh} \to K^{sep}$ を選ぶ。

\begin{lemma}
\label{pione-lemma-normal-pione-quotient-inertia}
$X$ を関数体 $K$ をもつ正規整スキームとする。
上記の記号のもとで、次の三つの
$\text{Gal}(K^{sep}/K) = \pi_1(\Spec(K))$ の部分群は等しい。
\begin{enumerate}
\item 全射 $\text{Gal}(K^{sep}/K) \longrightarrow \pi_1(X)$ の核。
\item $I_y$ を含む最小の閉正規部分群（すべての $y \in X^{sep}$ について）。
\item $\text{Gal}(K^{sep}/K_x^{sh})$ を含む最小の閉正規部分群
(すべての $x \in X$ について)。
\end{enumerate}
\end{lemma}

\begin{proof}
補題 \ref{pione-lemma-identify-inertia} により (2) と (3) は同値である。
この補題は、$I_y$ が $\text{Gal}(K^{sep}/K_x^{sh})$ と共役であることを、
$y$ が $x$ の上にある場合に述べている。
補題 \ref{pione-lemma-local-exact-sequence-normal} により
$\text{Gal}(K^{sep}/K_x^{sh})$ は
$\pi_1(\Spec(\mathcal{O}_{X, x}))$ へ自明に写る。
したがって (2) と (3) の部分群
$N \subset G = \text{Gal}(K^{sep}/K)$ は
$G \longrightarrow \pi_1(X)$ の核に含まれる。

\medskip\noindent
他方の包含を示す。$N$ は正規なので、$N \subset U \subset G$ かつ $U$ が開正規であるとき、
商写像 $G \to G/U$ が $\pi_1(X)$ を経由することを示せば十分である。
言い換えると、$L/K$ を $U$ に対応するガロア拡大とすれば、
$X$ が $L$ で非分岐であることを示せばよい
(節 \ref{pione-section-normal}、特に命題 \ref{pione-proposition-normal})。
$X$ がアフィンの場合に示せば十分である（証明の残りで代数の結果を参照するためである）。
$Y \to X$ を $X$ の $L$ における正規化とする。
$L \subset K^{sep}$ の包含は射 $\pi : X^{sep} \to Y$ を誘導する。
$y \in X^{sep}$ に対し、$\pi(y)$ の慣性群は
$\text{Gal}(L/K)$ における $I_y$ の $\text{Gal}(L/K)$ における像である。これは More on Algebra, 補題
\ref{more-algebra-lemma-one-orbit-geometric-galois-compare} から従う。
$N \subset U$ なので、これらすべての慣性群は自明である。
補題 \ref{pione-lemma-inertia-invariants-etale} を適用すれば
$Y \to X$ はエタールであると結論する。
（別法として、補題 \ref{pione-lemma-local-exact-sequence-normal} により、
$Y$ の $\Spec(\mathcal{O}_{X, x})$ への引き戻しが、すべての $x \in X$ について
エタールであることを示し、少し作業を加えて結論してもよい。）
\end{proof}

\begin{example}
\label{pione-example-bigger-codim}
$X$ を関数体 $K$ をもつ正規整ネーター・スキームとする。
分岐軌跡の純粋性（下記参照）により、$X$ が正則ならば、
補題 \ref{pione-lemma-normal-pione-quotient-inertia} では
慣性群 $I = \pi_1(\Spec(K_x^{sh}))$ を、点 $x$（余次元 $1$ の $X$ における点）について
考えれば十分である。
一般にはこれでは十分でない。実際、
$Y = \mathbf{A}_k^n = \Spec(k[t_1, \ldots, t_n])$ とし、
$k$ は標数 $2$ でない体とする。
$G = \{\pm 1\}$ を位数 $2$ の群とし、座標の乗法により $Y$ に作用させ、
$$
X = \Spec(k[t_it_j, i, j \in \{1, \ldots, n\}])
$$
埋め込み $k[t_it_j] \subset k[t_1, \ldots, t_n]$ は、
次数 $2$ の射 $Y \to X$ を定める。この射は、極大イデアル
$\mathfrak m = (t_it_j)$ の上を除いて至る所で非分岐である。
このイデアルは余次元 $n$ の点であり、$X$ に属する。
\end{example}

\begin{lemma}
\label{pione-lemma-unramified}
$X$ を関数体 $K$ をもつ正規整スキームとする。
$L/K$ を有限拡大とし、$Y \to X$ を $X$ の $L$ における正規化とする。
次は同値である。
\begin{enumerate}
\item $X$ が節 \ref{pione-section-normal} の意味で $L$ において非分岐である。
\item $Y \to X$ がスキームの非分岐射である。
\item $Y \to X$ がスキームのエタール射である。
\item $Y \to X$ が有限エタール射である。
\item $x \in X$ に対する射
$Y \times_X \Spec(\mathcal{O}_{X, x}) \to \Spec(\mathcal{O}_{X, x})$ が非分岐である。
\item (5) と同じ。ただし $\mathcal{O}_{X, x}^h$ を用いる。
\item (5) と同じ。ただし $\mathcal{O}_{X, x}^{sh}$ を用いる。
\item $x \in X$ に対するスキーム論的ファイバー $Y_x$ が
$x$ 上エタールで次数 $\geq [L : K]$ である。
\end{enumerate}
さらに $L/K$ がガロア群 $G$ をもつガロア拡大なら、これらは次とも同値である。
\begin{enumerate}
\item[(9)] $y \in Y$ に対し、群
$I_y = \{g \in G \mid g(y) = y\text{ かつ }
g \bmod \mathfrak m_y = \text{id}_{\kappa(y)}\}$ が自明である。
\end{enumerate}
\end{lemma}

\begin{proof}
 (1) と (2) の同値性は (1) の定義である。
(2)、(3)、(4) の同値性は補題 \ref{pione-lemma-unramified-in-L} による。
(4) $\Rightarrow$ (5)、(5) $\Rightarrow$ (6)、(6) $\Rightarrow$ (7) の証明は容易である。

\medskip\noindent
 (7) を仮定する。$\mathcal{O}_{X, x}^{sh}$ は正規局所整域であることに注意する
(More on Algebra, 補題 \ref{more-algebra-lemma-henselization-normal})。
$L^{sh} = L \otimes_K K_x^{sh}$ とおく。ここで $K_x^{sh}$ は
$\mathcal{O}_{X, x}^{sh}$ の分数体である。このとき
$L^{sh} = \prod_{i = 1, \ldots, n} L_i$ であり、$L_i/K_x^{sh}$ は有限分離拡大である。
Algebra, 補題 \ref{algebra-lemma-integral-closure-commutes-smooth}
(および省略する極限の議論) により、
$Y \times_X \Spec(\mathcal{O}_{X, x}^{sh})$ は
$\Spec(\mathcal{O}_{X, x}^{sh})$ の $L^{sh}$ における整閉包である。
したがって補題 \ref{pione-lemma-unramified-in-L} を
$L_i$ を $L^{sh}$ の因子として適用すると、
$Y \times_X \Spec(\mathcal{O}_{X, x}^{sh}) \to \Spec(\mathcal{O}_{X, x}^{sh})$
は有限エタールである。一般点を調べれば次数は $[L : K]$ に等しいので、(8) が成り立つ。

\medskip\noindent
 (8) を仮定する。$x \in X$ とし、スキーム論的ファイバー $Y_x$ が
$x$ 上エタールで次数 $\geq [L : K]$ であるとする。これは $Y$ が
$\geq [L : K]$ 個の幾何点を $x$ の上にもつことを意味する。
$Y \to X$ が $x$ の近傍上有限エタールであることを示す。これで (1) が従う。
証明のため $X = \Spec(R)$、点 $x$ が素イデアル $\mathfrak p \subset R$ に対応し、
$Y = \Spec(S)$ と仮定してよい。More on Morphisms,
補題 \ref{more-morphisms-lemma-etale-makes-integral-split} を適用すると、
エタール近傍 $(U, u) \to (X, x)$ で
$Y \times_X U = V_1 \amalg \ldots \amalg V_m$ となり、各 $V_i$ は
ただ一つの点 $v_i$ をもち、その点は $u$ の上にあり、$\kappa(v_i)/\kappa(u)$ は純非分離となるものが得られる。
$U$ を必要なら縮小し、$U$ を一般点 $\xi$ をもつ正規整スキームと仮定してよい
(Descent, 補題 \ref{descent-lemma-locally-finite-nr-irred-local-fppf}、
\ref{descent-lemma-normal-local-smooth}、Properties, 補題
\ref{properties-lemma-normal-locally-finite-nr-irreducibles} を用いる)。
幾何点についての上の注意から $m \geq [L : K]$ である。
一方、More on Morphisms, 補題 \ref{more-morphisms-lemma-normalization-smooth-localization} により、
$\coprod V_i \to U$ は $U$ の $\Spec(L) \times_X U$ における正規化である。
$K \subset \kappa(\xi)$ は有限分離なので、
$\Spec(L) \times_X U = \Spec(\prod_{i = 1, \ldots, n} L_i)$ と書け、
$L_i/\kappa(\xi)$ は有限で $[L : K] = \sum [L_i : \kappa(\xi)]$ である。
 $V_j$ は各 $j$ について空でなく、$m \geq [L : K]$ なので、
$m = n$ および $[L_i : \kappa(\xi)] = 1$（すべての $i$）と結論する。
したがって $V_j \to U$ は同型、特にエタールであり、$Y \times_X U \to U$ もエタールである。
Descent, 補題 \ref{descent-lemma-descending-property-etale} により、
$Y \to X$ は $U \to X$ の像（$x$ の開近傍）上エタールである。

\medskip\noindent
 $L/K$ がガロア拡大で (9) が成り立つと仮定する。このとき補題
\ref{pione-lemma-inertial-invariants-unramified} により $Y \to X$ はエタールである。
(1) が (9) を含意する証明は省略する。
\end{proof}

\noindent
離散付値環の無限ガロア拡大の場合には、もう少し詳しく述べることができる。
そのために次の記法を導入する。整数の部分集合 $S \subset \mathbf{N}$ が
{\it 乗法的有向} であるとは、$1 \in S$ であり、$n, m \in S$ に対して
$k \in S$ で $n | k$ および $m | k$ を満たすものが存在することをいう。
$S$ 上の半順序を $n \geq_S m$ を $m | n$ と同値とする規則で定める。
体 $\kappa$ に対して、有限群の逆系 $\{\mu_n(\kappa)\}_{n \in S}$ と遷移写像を得る。
$$
\mu_n(\kappa) \longrightarrow \mu_m(\kappa),\quad
\zeta \longmapsto \zeta^{n/m}
$$
これは $n \geq_S m$ のときである。したがってプロ有限群
$$
\lim_{n \in S} \mu_n(\kappa)
$$
この極限は余フィルター付きである（$S$ が有向だからである）。
この構成は $\kappa$ に関して関手的であり、特に
$\text{Aut}(\kappa)$ はこのプロ有限群に作用する。
例えば $S = \{1, n\}$ ならば、これは $\mu_n(\kappa)$ を与える。
$S = \{1, \ell, \ell^2, \ell^3, \ldots\}$ で、$\ell$ が $\kappa$ の標数と異なる素数なら、
これは $\lim_n \mu_{\ell^n}(\kappa)$ を与える。
これはしばしば $\ell$ 進 Tate 加群と呼ばれる $\kappa$ の乗法群である
(More on Algebra, 例 \ref{more-algebra-example-spectral-sequence-principal} を参照)。

\begin{lemma}
\label{pione-lemma-structure-decomposition}
離散付値環 $A$（分数体 $K$）とする。
$L/K$ を（無限でもよい）ガロア拡大とする。
$B$ を $A$ の $L$ における整閉包とし、極大イデアル $\mathfrak m$ をとる。
これは $B$ の極大イデアルである。
$G = \text{Gal}(L/K)$、
$D = \{\sigma \in G \mid \sigma(\mathfrak m) = \mathfrak m\}$、
$I = \{\sigma \in D \mid \sigma \bmod \mathfrak m =
\text{id}_{\kappa(\mathfrak m)}\}$ とおく。
分解群 $D$ は次の標準的な完全列に入る。
$$
1 \to I \to D \to \text{Aut}(\kappa(\mathfrak m)/\kappa_A) \to 1
$$
慣性群 $I$ は次の標準的な完全列に入る。
$$
1 \to P \to I \to I_t \to 1
$$
ここで次が成り立つ。
\begin{enumerate}
\item $P$ は $D$ の正規部分群である。
\item $P$ はプロ-$p$-群であり、$\kappa_A$ の標数が $p > 1$ のときそうである。
$P = \{1\}$ であるのは $\kappa_A$ の標数が 0 のときである。
\item 乗法的に有向な $S \subset \mathbf{N}$ が存在し、
$\kappa(\mathfrak m)$ は原始 $n$ 乗根を各 $n \in S$ について含む
(ただし $S$ の元は $p$ と互いに素である)。
\item 次の標準的な全射が存在する。
$$
\theta_{can} : I \to \lim_{n \in S} \mu_n(\kappa(\mathfrak m))
$$
その核は $P$ であり、
$\theta_{can}(\tau \sigma \tau^{-1}) = \tau(\theta_{can}(\sigma))$
を $\tau \in D$、$\sigma \in I$ に対して満たし、
$I_t \to \lim_{n \in S} \mu_n(\kappa(\mathfrak m))$ という同型を誘導する。
\end{enumerate}
\end{lemma}

\begin{proof}
これは主に More on Algebra, 節 \ref{more-algebra-section-ramification} で証明した
有限ガロア拡大に関する結果の言い換えである。
$D \to \text{Aut}(\kappa(\mathfrak m)/\kappa)$ の全射性は
More on Algebra, 補題 \ref{more-algebra-lemma-one-orbit-geometric-galois} による。
これで最初の完全列を得る。

\medskip\noindent
第二の短完全列を構成するため、有限ガロア部分拡大の集合を $\Lambda$ とする。
すなわち $\lambda \in \Lambda$ は $L/L_\lambda/K$ に対応する。
$G_\lambda = \text{Gal}(L_\lambda/K)$ とおく。
$G_\lambda$ は全射な遷移写像をもつ有限群の逆系であり、
$G = \lim_{\lambda \in \Lambda} G_\lambda$ であることを思い出そう。
Fields, 補題 \ref{fields-lemma-infinite-galois-limit} を参照。
$B_\lambda$ を $A$ の $L_\lambda$ における整閉包とする。
次に $\mathfrak m_\lambda = \mathfrak m \cap B_\lambda$ とおき、
$P_\lambda, I_\lambda, D_\lambda$ を $\mathfrak m_\lambda$ の
野生慣性群、慣性群、分解群と書く。More on Algebra, 補題
\ref{more-algebra-lemma-galois-inertia} を参照。
$\lambda \geq \lambda'$ のとき、制限により次の可換図式を得る。
$$
\xymatrix{
P_\lambda \ar[d] \ar[r] &
I_\lambda \ar[d] \ar[r] &
D_\lambda \ar[d] \ar[r] &
G_\lambda \ar[d] \\
P_{\lambda'} \ar[r] &
I_{\lambda'} \ar[r] &
D_{\lambda'} \ar[r] &
G_{\lambda'}
}
$$
垂直写像は全射である。More on Algebra, 補題
\ref{more-algebra-lemma-compare-inertia} を参照。

\medskip\noindent
定義から直ちに、上の同型のもとで
$I = \lim I_\lambda$ および $D = \lim D_\lambda$、さらに $G = \lim G_\lambda$
であることが従う。
$L = \colim L_\lambda$ なので、$B = \colim B_\lambda$ および
$\kappa(\mathfrak m) = \colim \kappa(\mathfrak m_\lambda)$ である。
$D_\lambda$ の系の遷移写像は
$D_\lambda \to \text{Aut}(\kappa(\mathfrak m_\lambda)/\kappa)$ と両立する
(More on Algebra, 補題 \ref{more-algebra-lemma-compare-inertia})。
したがって写像 $D \to \text{Aut}(\kappa(\mathfrak m)/\kappa)$ は、
写像 $D_\lambda \to \text{Aut}(\kappa(\mathfrak m_\lambda)/\kappa)$ の極限である。

\medskip\noindent
標準的な写像
$$
\theta_{\lambda, can} :
I_\lambda
\longrightarrow
\mu_{n_\lambda}(\kappa(\mathfrak m_\lambda))
$$
が存在する。ここで $n_\lambda = |I_\lambda|/|P_\lambda|$ であり、
$\mu_{n_\lambda}(\kappa(\mathfrak m_\lambda))$ は位数 $n_\lambda$ をもち、
$\theta_{\lambda, can}(\tau \sigma \tau^{-1}) =
\tau(\theta_{\lambda, can}(\sigma))$
を $\tau \in D_\lambda$ および $\sigma \in I_\lambda$ に対して満たす。
さらに次の可換図式を得る。
$$
\xymatrix{
I_\lambda \ar[r]_-{\theta_{\lambda, can}} \ar[d] &
\mu_{n_\lambda}(\kappa(\mathfrak m_\lambda))
\ar[d]^{(-)^{n_\lambda/n_{\lambda'}}} \\
I_{\lambda'} \ar[r]^-{\theta_{\lambda', can}} &
\mu_{n_{\lambda'}}(\kappa(\mathfrak m_{\lambda'}))
}
$$
More on Algebra, 注意 \ref{more-algebra-remark-canonical-inertia-character} を参照。

\medskip\noindent
$S \subset \mathbf{N}$ を整数 $n_\lambda$ たちの集合とする。
$\Lambda$ は有向なので、$S$ は乗法的に有向である。
上に示した可換図式により、写像 $\theta_{\lambda, can}$ の極限をとって
$$
\theta_{can} : I \to \lim_{n \in S} \mu_n(\kappa(\mathfrak m)).
$$
を得る。この写像は連続である（細部は省略する）。
$I_\lambda$ の系の遷移写像は全射であり、$\Lambda$ は有向なので、
射影 $I \to I_\lambda$ は全射である。各 $\lambda$ について次の図式は
$$
\xymatrix{
I \ar[d] \ar[r]_-{\theta_{can}} &
\lim_{n \in S} \mu_n(\kappa(\mathfrak m)) \ar[d] \\
I_{\lambda} \ar[r]^-{\theta_{\lambda, can}} &
\mu_{n_\lambda}(\kappa(\mathfrak m_\lambda))
}
$$
可換である。したがって $\theta_{can}$ の像は有限群
$\mu_{n_\lambda}(\kappa(\mathfrak m)) =
\mu_{n_\lambda}(\kappa(\mathfrak m_\lambda))$（位数 $n_\lambda$）へ全射する
（上記参照）。ゆえに $\theta_{can}$ の像は稠密である。
一方、$\theta_{can}$ は連続で、その始域はプロ有限群である。
したがって位相的議論により $\theta_{can}$ は全射である。

\medskip\noindent
$\theta_{can}(\tau \sigma \tau^{-1}) = \tau(\theta_{can}(\sigma))$
という性質が $\tau \in D$、$\sigma \in I$ に対して成り立つことは、
写像 $\theta_{\lambda, can}$ の対応する性質と、写像
$D \to \text{Aut}(\kappa(\mathfrak m))$ と写像
$D_\lambda \to \text{Aut}(\kappa(\mathfrak m_\lambda))$ の両立性から従う。
$P = \Ker(\theta_{can})$ とおくと、これは $P$ が $D$ の正規部分群であることを意味する。
$I_t = I/P$ とおく。同型
$I_t \to \lim_{n \in S} \mu_n(\kappa(\mathfrak m))$ は
$\theta_{can}$ の全射性から得られる。

\medskip\noindent
証明を終えるため、$P = \lim P_\lambda$ を示す。これにより $P$ はプロ-$p$-群である。
野生でない慣性群 $I_{\lambda, t} = I_\lambda/P_\lambda$ の位数は $n_\lambda$ である。
遷移写像 $P_\lambda \to P_{\lambda'}$ は全射であり、$\Lambda$ は有向なので、
次の短完全列を得る。
$$
1 \to \lim P_\lambda \to I \to \lim I_{\lambda, t} \to 1
$$
（詳細は省略する）。各 $\lambda$ について写像 $\theta_{\lambda, can}$
が同型 $I_{\lambda, t} \cong \mu_{n_\lambda}(\kappa(\mathfrak m))$
を誘導するので、求める結果が従う。
\end{proof}

\begin{lemma}
\label{pione-lemma-structure-decomposition-separable-closure}
離散付値環 $A$（分数体 $K$）とする。
分離閉包 $K^{sep}$ を $K$ に対してとる。
$A^{sep}$ を $A$ の $K^{sep}$ における整閉包とする。
極大イデアル $\mathfrak m^{sep}$ を $A^{sep}$ にとる。
$\mathfrak m = \mathfrak m^{sep} \cap A$ とおき、
$\kappa = A/\mathfrak m$、さらに
$\overline{\kappa} = A^{sep}/\mathfrak m^{sep}$ とおく。
すると $\overline{\kappa}$ は $\kappa$ の代数閉包である。
$G = \text{Gal}(K^{sep}/K)$、
$D = \{\sigma \in G \mid \sigma(\mathfrak m^{sep}) = \mathfrak m^{sep}\}$、
および
$I = \{\sigma \in D \mid \sigma \bmod \mathfrak m^{sep} =
\text{id}_{\kappa(\mathfrak m^{sep})}\}$ とおく。
分解群 $D$ は次の標準的な完全列に入る。
$$
1 \to I \to D \to \text{Gal}(\kappa^{sep}/\kappa) \to 1
$$
$\kappa^{sep} \subset \overline{\kappa}$ は $\kappa$ の分離閉包である。
慣性群 $I$ は次の標準的な完全列に入る。
$$
1 \to P \to I \to I_t \to 1
$$
ここで次が成り立つ。
\begin{enumerate}
\item $P$ は $D$ の正規部分群である。
\item $P$ はプロ-$p$-群であり、$\kappa_A$ の標数が $p > 1$ のときそうである。
$P = \{1\}$ であるのは $\kappa_A$ の標数が 0 のときである。
\item 次の標準的な全射が存在する。
$$
\theta_{can} : I \to \lim_{n\text{ prime to }p} \mu_n(\kappa^{sep})
$$
その核は $P$ であり、
$\theta_{can}(\tau \sigma \tau^{-1}) = \tau(\theta_{can}(\sigma))$
を $\tau \in D$、$\sigma \in I$ に対して満たし、さらに同型
$I_t \to \lim_{n\text{ prime to }p} \mu_n(\kappa^{sep})$.
\end{enumerate}
\end{lemma}

\begin{proof}
$\overline{\kappa}$ は補題 \ref{pione-lemma-get-algebraic-closure} により
$\kappa$ の代数閉包である。
このほかの主張の大部分は、補題 \ref{pione-lemma-structure-decomposition} の
対応する部分から直ちに従う。例えば
$\text{Aut}(\overline{\kappa}/\kappa) = \text{Gal}(\kappa^{sep}/\kappa)$
なので、最初の完全列を得る。
残る証明すべき主張は、補題 \ref{pione-lemma-structure-decomposition} の $S$ に対して
極限 $\lim_{n \in S} \mu_n(\overline{\kappa})$ が
$\lim_{n\text{ prime to }p} \mu_n(\kappa^{sep})$ に等しいことである。
これを示すには、$n$ が $p$ と互いに素であり、$S$ の元を割ることを示せば十分である。
一様化元 $\pi \in A$ をとり、$L$ を多項式 $X^n - \pi$ の分解体とする。
この多項式は分離的なので、$L$ は $K$ の有限ガロア拡大である。
$L \to K^{sep}$ という埋め込みを一つ選ぶ。
$B$ を $A$ の $L$ における整閉包とすると、
$A \to B_{\mathfrak m^{sep} \cap B}$ の分岐指数は
$n$ で割り切れる（$\pi$ の $n$ 乗根が $B$ にあるためである。実際には
分岐指数は $n$ に等しいが、ここではそれを必要としない）。
補題 \ref{pione-lemma-structure-decomposition} の証明における
$S$ の構成から、$n$ は $S$ の元を割ることが従う。
\end{proof}








\section{幾何学的および算術的基本群}
\label{pione-section-galois-action}

\noindent
本節では、スキーム $X$ と体 $k$ の基本群を、
$X_{\overline{k}}$ の基本群と比較したとき何が起こるかを調べる。
ここで $\overline{k}$ は $k$ の代数閉包である。

\begin{lemma}
\label{pione-lemma-limit}
有向集合 $I$ とする。$X_i$ を準コンパクトかつ準分離なスキームの
$I$ 上の、遷移射がアフィンである逆系とする。
Limits, 節 \ref{limits-section-limits} のように $X = \lim X_i$ とおく。
すると圏の同値
$$
\colim \textit{F\'Et}_{X_i} = \textit{F\'Et}_X
$$
 $X_i$ が十分大きい $i$ について連結であり、幾何点 $\overline{x}$ が
$X$ のものであるなら、
$$
\pi_1(X, \overline{x}) = \lim \pi_1(X_i, \overline{x})
$$
\end{lemma}

\begin{proof}
圏の同値は Limits, 補題
\ref{limits-lemma-descend-finite-presentation},
\ref{limits-lemma-descend-finite-finite-presentation}, および
\ref{limits-lemma-descend-etale}.
第二の主張は圏についての主張から形式的に従う。
\end{proof}

\begin{lemma}
\label{pione-lemma-perfection}
体 $k$ とその完全化 $k^{perf}$ をとる。連結スキーム $X$ を $k$ 上にとる。
すると $X_{k^{perf}}$ は連結であり、
$\pi_1(X_{k^{perf}}) \to \pi_1(X)$ は同型である。
\end{lemma}

\begin{proof}
基本群の位相的不変性の特別な場合である。
命題 \ref{pione-proposition-universal-homeomorphism} を参照。
$\Spec(k^{perf}) \to \Spec(k)$ が普遍同相であることを見るには、
Algebra, 補題 \ref{algebra-lemma-radicial-integral-bijective} を用いればよい。
\end{proof}

\begin{lemma}
\label{pione-lemma-ses-field}
体 $k$ とその代数閉包 $\overline{k}$ をとる。
準コンパクトかつ準分離なスキーム $X$ を体 $k$ 上にとる。
基底変換 $X_{\overline{k}}$ が連結なら、短完全列
$$
1 \to \pi_1(X_{\overline{k}}) \to \pi_1(X) \to \pi_1(\Spec(k)) \to 1
$$
をもつプロ有限位相群の完全列である。
\end{lemma}

\begin{proof}
連結な $\textit{F\'Et}_{\Spec(k)}$ の対象は、
$\Spec(k') \to \Spec(k)$ の形であり、ここで $k'/k$ は有限分離拡大である。
すると $X_{\Spec{k'}}$ は連結である。なぜなら射
$X_{\overline{k}} \to X_{\Spec(k')}$ は全射であり、
仮定により $X_{\overline{k}}$ は連結だからである。したがって
$\pi_1(X) \to \pi_1(\Spec(k))$ は
補題 \ref{pione-lemma-functoriality-galois-surjective} により全射である。

\medskip\noindent
先に進む前に、$k$ が完全体であると仮定してよいことに注意する。
実際、補題 \ref{pione-lemma-perfection} により
$\pi_1(X_{k^{perf}}) = \pi_1(X)$ および
$\pi_1(\Spec(k^{perf})) = \pi_1(\Spec(k))$ である。

\medskip\noindent
関手の合成
$\textit{F\'Et}_{\Spec(k)} \to \textit{F\'Et}_X \to
\textit{F\'Et}_{X_{\overline{k}}}$ は、対象を
$X_{\Spec(\overline{k})}$ のコピーの互いに素な合併へ送ることは明らかである。
したがって合成
$\pi_1(X_{\overline{k}}) \to \pi_1(X) \to \pi_1(\Spec(k))$
は、補題 \ref{pione-lemma-composition-trivial} により自明な準同型である。

\medskip\noindent
有限エタール射 $U \to X$ で $U$ が連結なものをとる。
$U \times_X X_{\overline{k}} = U_{\overline{k}}$ に注意する。
$U_{\overline{k}} \to X_{\overline{k}}$ が断面
$s : X_{\overline{k}} \to U_{\overline{k}}$ をもつと仮定する。
すると $s(X_{\overline{k}})$ は $U_{\overline{k}}$ の開連結成分である。
$\sigma \in \text{Gal}(\overline{k}/k)$ に対し、その基底変換を $s^\sigma$ と書く。
これは $s$ の $\Spec(\sigma)$ による基底変換である。
$U_{\overline{k}} \to X_{\overline{k}}$ は有限エタールなので、
断面は有限個しかない。したがって
$$
\overline{T} = \bigcup s^\sigma(X_{\overline{k}})
$$
これは有限合併であり、$\overline{T}$ は
$\text{Gal}(\overline{k}/k)$-安定な開閉部分集合である。
Varieties, 補題 \ref{varieties-lemma-closed-fixed-by-Galois} により、
$\overline{T}$ は閉部分集合 $T \subset U$ の逆像である。
$U_{\overline{k}} \to U$ は開写像である
(Morphisms, 補題 \ref{morphisms-lemma-scheme-over-field-universally-open})
から、$T$ も開である。$U$ は連結なので $T = U$ である。
したがって $U_{\overline{k}}$ は $X_{\overline{k}}$ のコピーの
(有限な)互いに素な合併である。補題
\ref{pione-lemma-functoriality-galois-normal} により、
$\pi_1(X_{\overline{k}}) \to \pi_1(X)$ の像は正規である。

\medskip\noindent
有限エタール被覆 $V \to X_{\overline{k}}$ をとる。
$\overline{k}$ は $k$ の有限分離拡大の合併であることを思い出そう。
補題 \ref{pione-lemma-limit} により、有限分離拡大 $k'/k$ と有限エタール射
$U \to X_{k'}$ であって
$V = X_{\overline{k}} \times_{X_{k'}} U =
U \times_{\Spec(k')} \Spec(\overline{k})$
となるものを得る。
すると合成 $U \to X_{k'} \to X$ は有限エタールであり、
$U \times_{\Spec(k)} \Spec(\overline{k})$ は
$V = U \times_{\Spec(k')} \Spec(\overline{k})$ を開閉部分スキームとして含む。
（$\Spec(\overline{k})$ は
$\Spec(k') \times_{\Spec(k)} \Spec(\overline{k})$ の開閉部分スキームであり、
これは積写像 $k' \otimes_k \overline{k} \to \overline{k}$ による。）
補題 \ref{pione-lemma-functoriality-galois-injective} により、
$\pi_1(X_{\overline{k}}) \to \pi_1(X)$ は単射である。

\medskip\noindent
最後に、任意の有限エタール射 $U \to X$ で、$U_{\overline{k}}$ が
$X_{\overline{k}}$ のコピーの互いに素な合併であるものに対して、
有限エタール射 $V \to \Spec(k)$ と全射
$V \times_{\Spec(k)} X \to U$ が存在することを示す。
補題 \ref{pione-lemma-functoriality-galois-ses} を参照。
上と同様に補題 \ref{pione-lemma-limit} を用いると、有限分離拡大 $k'/k$ で
同型
$U_{k'} \cong \coprod_{i = 1, \ldots, n} X_{k'}$
が存在するものを得る。
したがって $V = \coprod_{i = 1, \ldots, n} \Spec(k')$ とおけばよい。
\end{proof}






\section{ホモトピー完全列}
\label{pione-section-homotopy-exact-sequence}

\noindent
本節では次の結果を議論する。
$f : X \to S$ を有限表示の平坦固有射で、その幾何学的ファイバーが
連結かつ被約であるものとする。
$S$ は連結であり、幾何点 $\overline{s}$ を $S$ からとると仮定する。
このとき基本群の完全列
$$
\pi_1(X_{\overline{s}}) \to \pi_1(X) \to \pi_1(S) \to 1
$$
を得る。命題 \ref{pione-proposition-first-homotopy-sequence} を参照。

\begin{lemma}
\label{pione-lemma-stein-factorization-etale}
\begin{reference}
\cite[Expose X, Proposition 1.2, p. 262]{SGA1}.
\end{reference}
スキームの固有射 $f : X \to S$ をとる。
$X \to S' \to S$ を $f$ の Stein 因子分解とする。More on Morphisms, 定理
\ref{more-morphisms-theorem-stein-factorization-general} を参照。
$f$ が有限表示、平坦で、幾何学的に被約なファイバーをもつなら、
$S' \to S$ は有限エタールである。
\end{lemma}

\begin{proof}
これは Derived Categories of Schemes,
補題 \ref{perfect-lemma-proper-flat-geom-red} と、
More on Morphisms, 定理
\ref{more-morphisms-theorem-stein-factorization-general} に含まれる情報から従う。
\end{proof}

\begin{proposition}
\label{pione-proposition-first-homotopy-sequence}
$f : X \to S$ を有限表示の平坦固有射で、その幾何学的ファイバーが
連結かつ被約であるものとする。$S$ は連結であり、幾何点
$\overline{s}$ を $S$ からとると仮定する。このとき基本群の完全列
$$
\pi_1(X_{\overline{s}}) \to \pi_1(X) \to \pi_1(S) \to 1
$$
を得る。
\end{proposition}

\begin{proof}
有限エタール射 $Y \to X$ をとる。Stein 因子分解
$$
\xymatrix{
Y \ar[d] \ar[r] & X \ar[d] \\
T \ar[r] & S
}
$$
これは $Y \to S$ の Stein 因子分解である。
補題 \ref{pione-lemma-stein-factorization-etale} により射 $T \to S$ は有限エタールである。
このように関手 $\textit{F\'Et}_X \to \textit{F\'Et}_S$ を得る。
任意の有限エタール射 $U \to S$ に対し、射
 $Y \to U \times_S X$ は $X$ 上の射であり、射
 $Y \to U$ は $S$ 上の射と同じであり、
このような射は Stein 因子分解を一意に経由する。すなわち一意な射
$T \to U$ に対応する
(Stein 因子分解を More on Morphisms, 補題
\ref{more-morphisms-lemma-stein-universally-closed} における相対正規化として
構成し、さらに Morphisms, 補題
\ref{morphisms-lemma-characterize-normalization} により因子分解することによる)。
したがって関手
$\textit{F\'Et}_X \to \textit{F\'Et}_S$ と
$\textit{F\'Et}_S \to \textit{F\'Et}_X$ は随伴である。
$U \times_S X \to S$ の Stein 因子分解は $U$ であることにも注意する。
これは $U \times_S X \to U$ のファイバーが幾何学的に連結だからである。

\medskip\noindent
上の議論と Categories, 補題
\ref{categories-lemma-adjoint-fully-faithful} により、
$\textit{F\'Et}_S \to \textit{F\'Et}_X$
は充満忠実である。すなわち
$\pi_1(X) \to \pi_1(S)$ は全射である
(補題 \ref{pione-lemma-functoriality-galois-surjective})。

\medskip\noindent
合成
$\textit{F\'Et}_S \to \textit{F\'Et}_X \to \textit{F\'Et}_{X_{\overline{s}}}$
が任意の $U$ を $X_{\overline{s}}$ のコピーの互いに素な合併へ送ることは明らかである。
したがって $\pi_1(X_{\overline{s}}) \to \pi_1(X) \to \pi_1(S)$ は
補題 \ref{pione-lemma-composition-trivial} により自明である。

\medskip\noindent
有限エタール射 $Y \to X$ で $Y$ が連結であり、
 $Y \times_X X_{\overline{s}}$ が連結成分 $Z$ を含み、
 $X_{\overline{s}}$ と同型であるものをとる。
上記の Stein 因子分解 $T$ を考える。
$\overline{t} \in T_{\overline{s}}$ をファイバー $Z$ に対応する点とする。
$T$ は連結スキームの像なので連結であり、上の全射性により
$T \times_S X$ も連結である。
次の因子分解を考える。
$$
\pi : Y \longrightarrow T \times_S X
$$
任意の閉点 $\overline{x} \in X_{\overline{s}}$ をとる。
$\kappa(\overline{t}) = \kappa(\overline{s}) = \kappa(\overline{x})$
は代数閉体であることに注意する。
すると $\pi$ の $(\overline{t}, \overline{x})$ 上のファイバーはただ一つの点、
すなわち $\overline{z} \in Z$ からなる。この点は
$\overline{x} \in X_{\overline{s}}$ に同型 $Z \to X_{\overline{s}}$ によって対応する。
したがって有限エタール射 $\pi$ は次数 $1$ を
$(\overline{t}, \overline{x})$ の近傍でともなう。
$T \times_S X$ は連結なので、全体で次数 $1$ であり、
$Y \cong T \times_S X$ を得る。よって
$Y \times_X X_{\overline{s}}$ は完全分解する。
以上を合わせると、補題
\ref{pione-lemma-functoriality-galois-ses} および
\ref{pione-lemma-functoriality-galois-normal} の双方を適用でき、
証明が完了する。
\end{proof}




\section{特殊化写像}
\label{pione-section-specialization-map}

\noindent
本節では特殊化写像を構成する。
$f : X \to S$ を幾何学的に連結なファイバーをもつスキームの固有射とする。
$s' \leadsto s$ を $S$ の点の特殊化とする。
$\overline{s}$ および $\overline{s}'$ をそれぞれ $s$ および $s'$
上にある幾何点とする。このとき特殊化写像
$$
sp : \pi_1(X_{\overline{s}'}) \longrightarrow \pi_1(X_{\overline{s}})
$$
この写像の構成は次のとおりである。$A$ を
$\mathcal{O}_{S, s}$ の厳密ヘンゼル化で、$\kappa(s) \subset \kappa(s)^{sep} \subset \kappa(\overline{s})$
に関するものとする。Algebra, 定義
\ref{algebra-definition-henselization} を参照。
$s' \leadsto s$ なので、点 $s'$ に対応する
$\Spec(\mathcal{O}_{S, s})$ の点を考える。したがって $\Spec(A)$ の上の点 $s'$ で、その剰余体が $\kappa(s')$ の分離代数的拡大となるものが
少なくとも一つ（場合によっては複数）存在する。
$\kappa(\overline{s}')$ は代数閉なので、可換図式を与える射
$\varphi : \overline{s}' \to \Spec(A)$ を選べる。
$$
\xymatrix{
\overline{s}' \ar[r]_-\varphi \ar[rd] &
\Spec(A) \ar[d] &
\overline{s} \ar[l] \ar[ld] \\
& S
}
$$
特殊化写像は次の合成である。
$$
\pi_1(X_{\overline{s}'}) \longrightarrow
\pi_1(X_A) =
\pi_1(X_{\kappa(s)^{sep}}) =
\pi_1(X_{\overline{s}})
$$
ここで第一の等号は補題
\ref{pione-lemma-finite-etale-on-proper-over-henselian} によるものであり、
第二の等号は補題 \ref{pione-lemma-perfection} および
\ref{pione-lemma-finite-etale-invariant-over-proper} から従う。
構成から、特殊化写像は次の可換図式に入る。
$$
\xymatrix{
\pi_1(X_{\overline{s}'}) \ar[rr]_{sp} \ar[rd] & &
\pi_1(X_{\overline{s}}) \ar[ld] \\
& \pi_1(X)
}
$$
ここで $X$ は連結と仮定する。特殊化写像は上の $\varphi : \overline{s}' \to \Spec(A)$ の選択に依存する。
この依存性を示すときは $sp_\varphi$ と書く。

\begin{lemma}
\label{pione-lemma-specialization-map-base-change}
次の可換図式を考える。
$$
\xymatrix{
Y \ar[d]_g \ar[r] & X \ar[d]^f \\
T \ar[r] & S
}
$$
ここで $f$ と $g$ は幾何学的に連結なファイバーをもつ固有射である。
$t' \leadsto t$ を $T$ の点の特殊化とし、上で構成した特殊化写像
$sp : \pi_1(Y_{\overline{t}'}) \to \pi_1(Y_{\overline{t}})$ を考える。
このとき次の可換図式が存在する。
$$
\xymatrix{
\pi_1(Y_{\overline{t}'}) \ar[r]_{sp} \ar[d] & \pi_1(Y_{\overline{t}}) \ar[d] \\
\pi_1(X_{\overline{s}'}) \ar[r]^{sp} & \pi_1(X_{\overline{s}})
}
$$
これは特殊化写像の図式であり、$\overline{s}$ と $\overline{s}'$
はそれぞれ $\overline{t}$ と $\overline{t}'$ の像である。
\end{lemma}

\begin{proof}
 $B$ を $\mathcal{O}_{T, t}$ の厳密ヘンゼル化で、$\kappa(t) \subset \kappa(t)^{sep} \subset \kappa(\overline{t})$
に関するものとする。
特殊化写像の構成と同様に、$\psi : \overline{t}' \to \Spec(B)$ を選ぶ。
これは $\overline{t}' \to T$ の持ち上げである。
$s$ および $s'$ を、$t$ および $t'$ の $S$ における像とする。
$A$ を $\mathcal{O}_{S, s}$ の厳密ヘンゼル化で、
$\kappa(s) \subset \kappa(s)^{sep} \subset \kappa(\overline{s})$ に関するものとする。
$\kappa(\overline{s}) = \kappa(\overline{t})$ であるから、
厳密ヘンゼル化の函手性（Algebra, 補題
\ref{algebra-lemma-strictly-henselian-functorial}）により、次の可換図式に入る
環準同型 $A \to B$ を得る。
$$
\xymatrix{
\overline{t}' \ar[r]_-\psi \ar[d] & \Spec(B) \ar[d] \ar[r] & T \ar[d] \\
\overline{s}' \ar[r]^-\varphi & \Spec(A) \ar[r] & S
}
$$
ここで射 $\varphi : \overline{s}' \to \Spec(A)$ は単に合成
$\overline{t}' \to \Spec(B) \to \Spec(A)$ とする。
基底変換を適用すると次の可換図式を得る。
$$
\xymatrix{
Y_{\overline{t}'} \ar[r] \ar[d] & Y_B \ar[d] \\
X_{\overline{s}'} \ar[r] & X_A
}
$$
特殊化写像の構成から、この図式の可換性は補題の図式の可換性を意味する。
\end{proof}

\begin{lemma}
\label{pione-lemma-specialization-map-composition}
幾何学的に連結なファイバーをもつ固有射 $f : X \to S$ をとる。
$s'' \leadsto s' \leadsto s$ を $S$ の点の特殊化とする。
特殊化写像の合成
$\pi_1(X_{\overline{s}''}) \to \pi_1(X_{\overline{s}'}) \to
\pi_1(X_{\overline{s}})$ は特殊化写像
$\pi_1(X_{\overline{s}''}) \to \pi_1(X_{\overline{s}})$ である。
\end{lemma}

\begin{proof}
 $\mathcal{O}_{S, s} \to A$ を $\kappa(s) \to \kappa(\overline{s})$ を用いて構成した厳密ヘンゼル化とする。
第一の特殊化写像の構成に用いた写像を $A \to \kappa(\overline{s}')$ とする。
$\mathcal{O}_{S, s'} \to A'$ を $\kappa(s') \subset \kappa(\overline{s}')$ を用いて構成した厳密ヘンゼル化と書く。
厳密ヘンゼル化の函手性により、写像
$A \to A'$ が存在し、$A' \to \kappa(\overline{s}')$ との合成は与えられた写像である。
(Algebra, 補題 \ref{algebra-lemma-map-into-henselian-colimit}).
次に、第二の特殊化写像の構成に用いた写像を
$A' \to \kappa(\overline{s}'')$ とする。このとき第一および第二の特殊化写像の
合成は特殊化写像
$\pi_1(X_{\overline{s}''}) \to \pi_1(X_{\overline{s}})$
$A \to A' \to \kappa(\overline{s}'')$ を用いて構成したものであることは明らかである。
\end{proof}

\noindent
幾何学的に連結なファイバーをもつ固有射 $X \to S$ をとる。
$R$ を分数体が代数閉な厳密ヘンゼル付値環とし、射 $\Spec(R) \to S$
をとる。$\eta, s \in \Spec(R)$ を一般点および閉点とする。
このとき特殊化写像
$$
sp_R : \pi_1(X_\eta) \to \pi_1(X_s)
$$
を基底変換 $X_R/\Spec(R)$ に対して考えられる。$\eta$ と $s$ はともに
剰余体が代数閉なので、これは意味をもつ。

\begin{lemma}
\label{pione-lemma-specialization-map-valuation-ring}
幾何学的に連結なファイバーをもつ固有射 $f : X \to S$ をとる。
$s' \leadsto s$ を $S$ の点の特殊化とし、
$sp : \pi_1(X_{\overline{s}'}) \to \pi_1(X_{\overline{s}})$
を特殊化写像とする。このとき、$R$ が分数体が代数閉な $S$ 上の厳密ヘンゼル付値環で、
$sp$ が上で定義した $sp_R$ と同型となるものが存在する。
\end{lemma}

\begin{proof}
$\mathcal{O}_{S, s} \to A$ を $\kappa(s) \to \kappa(\overline{s})$ を用いて構成した厳密ヘンゼル化とする。構成に用いた写像
$A \to \kappa(\overline{s}')$ により $sp$ を構成した。$R \subset \kappa(\overline{s}')$ を、
分数体が $\kappa(\overline{s}')$ で $A$ の像を支配する付値環とする。
Algebra, 補題 \ref{algebra-lemma-dominate} を参照。
例えば補題 \ref{pione-lemma-normal-local-domain-separablly-closed-fraction-field}
および Algebra, 補題 \ref{algebra-lemma-valuation-ring-normal} により、
$R$ は厳密ヘンゼルであることに注意する。したがって補題が従う。
\end{proof}

\noindent
幾何学的に連結なファイバーをもつ固有射 $X \to S$ をとる。
$R$ を厳密ヘンゼル離散付値環とし、射 $\Spec(R) \to S$ をとる。
$\eta, s \in \Spec(R)$ を一般点および閉点とする。このとき特殊化写像
$$
sp_R : \pi_1(X_{\overline{\eta}}) \to \pi_1(X_s)
$$
を基底変換 $X_R/\Spec(R)$ に対して考えられる。$s$ の剰余体は代数閉なので、
これは意味をもつ。

\begin{lemma}
\label{pione-lemma-specialization-map-discrete-valuation-ring}
幾何学的に連結なファイバーをもつ固有射 $f : X \to S$ をとる。
$s' \leadsto s$ を $S$ の点の特殊化とし、
$sp : \pi_1(X_{\overline{s}'}) \to \pi_1(X_{\overline{s}})$
を特殊化写像とする。$S$ がネーターならば、$R$ が $S$ 上の厳密ヘンゼル離散付値環で、
$sp$ が上で定義した $sp_R$ と同型となるものが存在する。
\end{lemma}

\begin{proof}
$\mathcal{O}_{S, s} \to A$ を $\kappa(s) \to \kappa(\overline{s})$ を用いて構成した厳密ヘンゼル化とする。構成に用いた写像
$A \to \kappa(\overline{s}')$ により $sp$ を構成した。$R \subset \kappa(\overline{s}')$ を
$A$ の像を支配する離散付値環とする。Algebra, 補題
\ref{algebra-lemma-exists-dvr} を参照。体の図式
$$
\xymatrix{
\kappa(\overline{s}) \ar[r] & k \\
A/\mathfrak m_A \ar[r] \ar[u] & R/\mathfrak m_R \ar[u]
}
$$
を $k$ が代数閉となるように選ぶ。$R^{sh}$ を $R$ から、$R \to k$ を用いて構成した
厳密ヘンゼル化とする。$A \to R^{sh}$ なる写像を得る。Algebra, 補題
\ref{algebra-lemma-strictly-henselian-functorial}.
More on Algebra, 補題 \ref{more-algebra-lemma-henselization-dvr} により
環 $R^{sh}$ は離散付値環である。その一般点および閉点を $\eta, o$ と書く。
$\Spec(R^{sh})$ のスキーム図式
$$
\xymatrix{
\overline{\eta} \ar[d] \ar[r] & \Spec(R^{sh}) \ar[d] &
\Spec(k) \ar[d] \ar[l] \\
\overline{s}' \ar[r] & \Spec(A) & \overline{s} \ar[l]
}
$$
は可換なので、次の可換図式を得る。
$$
\xymatrix{
\pi_1(X_{\overline{\eta}}) \ar[d] \ar[r]_{sp_{R^{sh}}} & \pi_1(X_o) \ar[d] \\
\pi_1(X_{\overline{s}'}) \ar[r]^{sp} & X_{\overline{s}}
}
$$
この写像の構成により、これは特殊化写像の図式である。
縦の矢印は同型であるから（補題
\ref{pione-lemma-finite-etale-invariant-over-proper}）、補題が証明される。
\end{proof}








\section{閉部分スキームへの制限}
\label{pione-section-lefschetz}

\noindent
本節では制限関手についていくつかの結果を証明する。
$$
\textit{F\'Et}_X \longrightarrow \textit{F\'Et}_Y,\quad
U \longmapsto V = U \times_X Y
$$
ここで $X$ はスキーム、$Y$ は閉部分スキームである。基本群の位相的不変性を用いると、
この関手の研究を有限局所自由加群上の完備化関手と関係づけられる。

\medskip\noindent
以下の補題では、Cohomology of Schemes, 節
\ref{coherent-section-coherent-formal} で定義した連接形式加群の概念を用いる。
ネーター・スキームと準連接イデアル層 $\mathcal{I} \subset \mathcal{O}_X$ が与えられたとき、
対象 $(\mathcal{F}_n)$ を $\textit{Coh}(X, \mathcal{I})$ の対象とし、
それが{\it 有限局所自由}であるとは、各 $\mathcal{F}_n$ が有限局所自由な
$\mathcal{O}_X/\mathcal{I}^n$-加群であることとする。

\begin{lemma}
\label{pione-lemma-restriction-fully-faithful}
 $X$ をネーター・スキーム、$Y \subset X$ をイデアル層
$\mathcal{I} \subset \mathcal{O}_X$ をもつ閉部分スキームとする。
完備化関手
$$
\textit{Coh}(\mathcal{O}_X)
\longrightarrow 
\textit{Coh}(X, \mathcal{I}),\quad
\mathcal{F} \longmapsto \mathcal{F}^\wedge
$$
が有限局所自由対象の充満部分圏上で忠実充満である（上記参照）と仮定する。
このとき制限関手 $\textit{F\'Et}_X \to \textit{F\'Et}_Y$ は忠実充満である。
\end{lemma}

\begin{proof}
有限エタール被覆の圏が内部 Hom をもつので（補題
\ref{pione-lemma-internal-hom-finite-etale}）、次を示せば十分である。
$U$ を $X$ 上有限エタールなものとし、射 $t : Y \to U$ を
$X$ 上にとると、一意な切断 $s : X \to U$ が存在して
$t = s|_Y$ を満たす。図式は次のとおり。
$$
\xymatrix{
& U \ar[d]^f \\
Y \ar[r] \ar[ru] & X \ar@{..>}@/^1em/[u]
}
$$
点線の矢印 $s$ を見つけることは、次の $\mathcal{O}_X$-代数写像を見つけることと同値である。
$$
s^\sharp : f_*\mathcal{O}_U \longrightarrow \mathcal{O}_X
$$
これは $Y$ のイデアル層で剰余をとると、与えられた代数写像
$t^\sharp : f_*\mathcal{O}_U \to \mathcal{O}_Y$ になる。
補題 \ref{pione-lemma-thickening} により、$t$ を写像の整合的な系
$t_n : Y_n \to U$ へ一意に持ち上げられ、したがって次の写像を得る。
$$
\lim t_n^\sharp : f_*\mathcal{O}_U \longrightarrow \lim \mathcal{O}_{Y_n}
$$
これは $X$ 上の代数層の写像である。
$f_*\mathcal{O}_U$ は有限局所自由な $\mathcal{O}_X$-加群なので、
一意な $\mathcal{O}_X$-加群写像
$\sigma : f_*\mathcal{O}_U \to \mathcal{O}_X$ を得る。その完備化は
$\lim t_n^\sharp$ である。
$\sigma$ が代数準同型であることを見るには、次の図式が可換であることを確かめればよい。
$$
\xymatrix{
f_*\mathcal{O}_U \otimes_{\mathcal{O}_X} f_*\mathcal{O}_U
\ar[r] \ar[d]_{\sigma \otimes \sigma} &
f_*\mathcal{O}_U \ar[d]^\sigma \\
\mathcal{O}_X \otimes_{\mathcal{O}_X} \mathcal{O}_X \ar[r] &
\mathcal{O}_X
}
$$
が可換である。任意の $n$ について、$Y_n$ へ制限した後、すなわち完備化関手を適用した後には
この図式が可換であることが分かる。したがって完備化関手の忠実性から
代数準同型であると結論する。
\end{proof}

\begin{lemma}
\label{pione-lemma-restriction-equivalence}
 $X$ をネーター・スキーム、$Y \subset X$ をイデアル層
$\mathcal{I} \subset \mathcal{O}_X$ をもつ閉部分スキームとする。
完備化関手
$$
\textit{Coh}(\mathcal{O}_X)
\longrightarrow 
\textit{Coh}(X, \mathcal{I}),\quad
\mathcal{F} \longmapsto \mathcal{F}^\wedge
$$
が有限局所自由対象の充満部分圏の間の同値である（上記参照）と仮定する。
このとき制限関手 $\textit{F\'Et}_X \to \textit{F\'Et}_Y$ は同値である。
\end{lemma}

\begin{proof}
制限関手が忠実充満であることは補題
\ref{pione-lemma-restriction-fully-faithful} による。

\medskip\noindent
有限エタール射 $U_1 \to Y$ をとる。証明を終えるため、$U_1$ が制限関手の本質像に入ることを示す。

\medskip\noindent
 $n \geq 1$ に対し、$Y_n$ を $n$ 次の $Y$ の無限小近傍とする。
補題 \ref{pione-lemma-thickening} により、一意な有限エタール射
$\pi_n : U_n \to Y_n$ が存在し、その $Y = Y_1$ への基底変換は
$U_1 \to Y_1$ を復元する。層
$\mathcal{F}_n = \pi_{n, *}\mathcal{O}_{U_n}$ を考える。
$\mathcal{F}_n$ を $\mathcal{O}_X$-加群として $X$ 上でみなすと、それは局所的に
$(\mathcal{O}_X/f^{n + 1}\mathcal{O}_X)^{\oplus r}$ と同型である。
この $(\mathcal{F}_n)$ は $\textit{Coh}(X, \mathcal{I})$ の有限局所自由対象である。
仮定により、有限局所自由な $\mathcal{O}_X$-加群 $\mathcal{F}$ と整合的な同型の系
$$
\mathcal{F}/\mathcal{I}^n\mathcal{F} \to \mathcal{F}_n
$$
を得る。これは $\mathcal{O}_X$-加群の同型である。

\medskip\noindent
 $\mathcal{F}$ 上の代数構造を構成するため、乗法写像
$\mathcal{F}_n \otimes_{\mathcal{O}_X} \mathcal{F}_n \to \mathcal{F}_n$ を考える。
これは $\mathcal{F}_n = \pi_{n, *}\mathcal{O}_{U_n}$ が代数層であることから得られる。
これらは次の写像を定める。
$$
(\mathcal{F}\otimes_{\mathcal{O}_X} \mathcal{F})^\wedge
\longrightarrow
\mathcal{F}^\wedge
$$
これは圏 $\textit{Coh}(X, \mathcal{I})$ の写像である。したがって仮定により、
$\mu : \mathcal{F}\otimes_{\mathcal{O}_X} \mathcal{F} \to \mathcal{F}$ を得て、
その $Y_n$ への制限が上の乗法写像を与えるようにできる。
補題の記述にある関手の忠実性から、$\mu$ は
$\mathcal{O}_X$-代数構造を $\mathcal{F}$ 上に定め、$\mathcal{F}_n$ の
与えられた代数構造と両立する。
次のように置く。
$$
U = \underline{\Spec}_X((\mathcal{F}, \mu))
$$
すると、有限局所自由スキーム $\pi : U \to X$ を得、その $Y$ への制限は
$U_1$ と同型である。
$\pi$ の判別式は、次の切断の零点集合である。
$$
\det(Q_\pi) :
\mathcal{O}_X
\longrightarrow
\wedge^{top}(\pi_*\mathcal{O}_U)^{\otimes -2}
$$
これは Discriminants, 節 \ref{discriminant-section-discriminant} で構成したものである。
その $Y_n$ への制限はすべての $n$ について同型であるから、Discriminants, 補題
\ref{discriminant-lemma-discriminant} により、それ自身も同型である。したがって $\pi$ は
Discriminants, 補題 \ref{discriminant-lemma-discriminant} によりエタールである。
\end{proof}

\begin{lemma}
\label{pione-lemma-restriction-fully-faithful-general}
 $X$ をネーター・スキーム、$Y \subset X$ をイデアル層
$\mathcal{I} \subset \mathcal{O}_X$ をもつ閉部分スキームとする。
$\mathcal{V}$ を開部分スキーム $V \subset X$ で $Y$ を含むものの集合とし、
逆包含で順序づける。完備化関手
$$
\colim_\mathcal{V} \textit{Coh}(\mathcal{O}_V)
\longrightarrow
\textit{Coh}(X, \mathcal{I}),
\quad
\mathcal{F} \longmapsto \mathcal{F}^\wedge
$$
が有限局所自由対象の充満部分圏上で忠実充満である（上記参照）と仮定する。
このとき制限関手
$\colim_\mathcal{V} \textit{F\'Et}_V \to \textit{F\'Et}_Y$
は忠実充満である。
\end{lemma}

\begin{proof}
 $\mathcal{V}$ は有向集合なので、余極限は Categories, 節
\ref{categories-section-directed-colimits} のとおりである。
残りの議論は補題 \ref{pione-lemma-restriction-fully-faithful} の証明と
ほとんど同じであるため、読者には読み飛ばすことを勧める。

\medskip\noindent
有限エタール被覆の圏が内部 Hom をもつので（補題
\ref{pione-lemma-internal-hom-finite-etale}）、次を示せば十分である。
$U$ を $V \in \mathcal{V}$ 上有限エタールなものとし、射 $t : Y \to U$ を
$V$ 上にとると、ある $V' \geq V$ と射 $s : V' \to U$ が存在して、
$V$ 上で $t = s|_Y$ となる。図式は次のとおり。
$$
\xymatrix{
& & U \ar[d]^f \\
Y \ar[r] \ar[rru] & V' \ar@{..>}[ru] \ar[r] & V
}
$$
点線の矢印 $s$ を見つけることは、次の
$\mathcal{O}_{V'}$-代数写像を見つけることと同値である。
$$
s^\sharp : f_*\mathcal{O}_U|_{V'} \longrightarrow \mathcal{O}_{V'}
$$
これは $Y$ のイデアル層で剰余をとると、与えられた代数写像
$t^\sharp : f_*\mathcal{O}_U \to \mathcal{O}_Y$ になる。
補題 \ref{pione-lemma-thickening} により、$t$ を写像の整合的な系
$t_n : Y_n \to U$ へ一意に持ち上げ、したがって次の写像を得る。
$$
\lim t_n^\sharp : f_*\mathcal{O}_U \longrightarrow \lim \mathcal{O}_{Y_n}
$$
これは $V$ 上の代数層の写像である。
$f_*\mathcal{O}_U$ は有限局所自由な $\mathcal{O}_V$-加群である。
したがって、ある $V' \geq V$ と写像
$\sigma : f_*\mathcal{O}_U|_{V'} \to \mathcal{O}_{V'}$
を得、その完備化は $\lim t_n^\sharp$ である。
$\sigma$ が代数準同型であることを見るには、次の図式
$$
\xymatrix{
(f_*\mathcal{O}_U \otimes_{\mathcal{O}_V} f_*\mathcal{O}_U)|_{V'}
\ar[r] \ar[d]_{\sigma \otimes \sigma} &
f_*\mathcal{O}_U|_{V'} \ar[d]^\sigma \\
\mathcal{O}_{V'} \otimes_{\mathcal{O}_{V'}} \mathcal{O}_{V'} \ar[r] &
\mathcal{O}_{V'}
}
$$
が可換であることを確かめればよい。任意の $n$ について $Y_n$ へ制限した後、
すなわち完備化関手を適用した後には、この図式が可換であることが分かる。
したがって完備化関手の忠実性から、ある $V'' \geq V'$ が存在して
$\sigma|_{V''}$ が代数準同型となることを得る。
\end{proof}

\begin{lemma}
\label{pione-lemma-restriction-equivalence-general}
 $X$ をネーター・スキーム、$Y \subset X$ をイデアル層
$\mathcal{I} \subset \mathcal{O}_X$ をもつ閉部分スキームとする。
$\mathcal{V}$ を開部分スキーム $V \subset X$ で $Y$ を含むものの集合とし、
逆包含で順序づける。完備化関手
$$
\colim_\mathcal{V} \textit{Coh}(\mathcal{O}_V)
\longrightarrow
\textit{Coh}(X, \mathcal{I}),
\quad
\mathcal{F} \longmapsto \mathcal{F}^\wedge
$$
が有限局所自由対象の充満部分圏の間の同値を定める（上記の説明参照）と仮定する。
このとき制限関手
$$
\colim_\mathcal{V} \textit{F\'Et}_V \to \textit{F\'Et}_Y
$$
は同値である。
\end{lemma}

\begin{proof}
 $\mathcal{V}$ は有向集合なので、余極限は Categories, 節
\ref{categories-section-directed-colimits} のとおりである。
残りの議論は補題 \ref{pione-lemma-restriction-equivalence} の証明と
ほとんど同じであるため、読者には読み飛ばすことを勧める。

\medskip\noindent
制限関手が忠実充満であることは補題
\ref{pione-lemma-restriction-fully-faithful-general} による。

\medskip\noindent
有限エタール射 $U_1 \to Y$ をとる。証明を終えるため、
$U_1$ が制限関手の本質像に入ることを示す。

\medskip\noindent
 $n \geq 1$ に対し、$Y_n$ を $n$ 次の $Y$ の無限小近傍とする。
補題 \ref{pione-lemma-thickening} により、一意な有限エタール射
$\pi_n : U_n \to Y_n$ が存在し、その $Y = Y_1$ への基底変換は
$U_1 \to Y_1$ を復元する。層
$\mathcal{F}_n = \pi_{n, *}\mathcal{O}_{U_n}$ を考える。
$\mathcal{F}_n$ を $\mathcal{O}_X$-加群として $X$ 上でみなすと、それは局所的に
$(\mathcal{O}_X/f^{n + 1}\mathcal{O}_X)^{\oplus r}$.
この $(\mathcal{F}_n)$ は $\textit{Coh}(X, \mathcal{I})$ の有限局所自由対象である。
仮定により、ある $V \in \mathcal{V}$、有限局所自由な
$\mathcal{O}_V$-加群 $\mathcal{F}$、および整合的な同型の系
$$
\mathcal{F}/\mathcal{I}^n\mathcal{F} \to \mathcal{F}_n
$$
が存在する。これは $\mathcal{O}_V$-加群の同型である。

\medskip\noindent
 $\mathcal{F}$ 上の代数構造を構成するため、乗法写像
$\mathcal{F}_n \otimes_{\mathcal{O}_V} \mathcal{F}_n \to \mathcal{F}_n$ を考える。
これは $\mathcal{F}_n = \pi_{n, *}\mathcal{O}_{U_n}$ が代数層であることから得られる。
これらは次の写像を定める。
$$
(\mathcal{F}\otimes_{\mathcal{O}_V} \mathcal{F})^\wedge
\longrightarrow
\mathcal{F}^\wedge
$$
これは圏 $\textit{Coh}(X, \mathcal{I})$ の写像である。したがって仮定により、
$V$ を縮小した後、次の写像が存在すると仮定できる。
$\mu : \mathcal{F}\otimes_{\mathcal{O}_V} \mathcal{F} \to \mathcal{F}$
その $Y_n$ への制限は上の乗法写像を与える。
さらに必要なら開部分を縮小して、$\mu$ が可換な
$\mathcal{O}_V$-代数構造を $\mathcal{F}$ 上に定め、$\mathcal{F}_n$ の与えられた
代数構造と両立すると仮定できる。次のように置く。
$$
U = \underline{\Spec}_V((\mathcal{F}, \mu))
$$
すると $V$ 上の有限局所自由スキームを得、その $Y$ への制限は
$U_1$ と同型である。したがって $U \to V$ は $Y$ の上にあるすべての点で
エタールである。More on Morphisms, 補題
\ref{more-morphisms-lemma-check-smoothness-on-infinitesimal-nbhds}.
により、さらに $V$ を一度縮小すれば $U \to V$ は有限エタールであると
仮定できる。これで証明が終わる。
\end{proof}

\begin{lemma}
\label{pione-lemma-restriction-faithful}
 $X$ をスキーム、$Y \subset X$ を閉部分スキームとする。
$X$ のすべての連結成分が $Y$ と交わるならば、
制限関手 $\textit{F\'Et}_X \to \textit{F\'Et}_Y$ は忠実である。
\end{lemma}

\begin{proof}
 $a, b : U \to U'$ を $X$ 上有限エタールなスキームの射とし、
$Y$ への制限が同じであるとする。$U$ の連結成分の像は $X$ の連結成分である。
これは $U \to X$ を $X$ の連結成分へ制限して Topology, 補題
\ref{topology-lemma-finite-fibre-connected-components} を適用すれば分かる。
したがって仮定により、$U$ のすべての連結成分の像は $Y$ と交わる。
よって \'Etale Morphisms, 命題
\ref{etale-proposition-equality} により、$a = b$ である。これは $U$ の各連結成分への
制限後に得られる。$a$ と $b$ の等化子は $U$ の開部分スキームである
（$U'$ の対角射が $X$ 上で開だから）ので、結論を得る。
\end{proof}

\begin{lemma}
\label{pione-lemma-restriction-fully-faithful-special}
 $X$ をネーター・スキーム、$Y \subset X$ を閉部分スキームとする。
$Y_n \subset X$ を $n$ 次の $Y$ の $X$ における無限小近傍とする。
次のいずれかが成り立つと仮定する。
\begin{enumerate}
\item $X$ は準アフィンであり、
$\Gamma(X, \mathcal{O}_X) \to \lim \Gamma(Y_n, \mathcal{O}_{Y_n})$
が同型である、または
\item $X$ は豊富な可逆加群 $\mathcal{L}$ をもち、
$\Gamma(X, \mathcal{L}^{\otimes m}) \to
\lim \Gamma(Y_n, \mathcal{L}^{\otimes m}|_{Y_n})$
がすべての $m \gg 0$ について同型である、または
\item 任意の有限局所自由な $\mathcal{O}_X$-加群 $\mathcal{E}$ について写像
$\Gamma(X, \mathcal{E}) \to \lim \Gamma(Y_n, \mathcal{E}|_{Y_n})$
が同型である。
\end{enumerate}
このとき制限関手 $\textit{F\'Et}_X \to \textit{F\'Et}_Y$ は忠実充満である。
\end{lemma}

\begin{proof}
この補題は、補題 \ref{pione-lemma-restriction-fully-faithful} と
Algebraic and Formal Geometry, 補題
\ref{algebraization-lemma-completion-fully-faithful} から形式的に従う。
\end{proof}

\begin{lemma}
\label{pione-lemma-restriction-fully-faithful-general-special}
 $X$ をネーター・スキーム、$Y \subset X$ を閉部分スキームとする。
$Y_n \subset X$ を $n$ 次の $Y$ の $X$ における無限小近傍とする。
$\mathcal{V}$ を開部分スキーム $V \subset X$ で $Y$ を含むものの集合とし、
逆包含で順序づける。次のいずれかが成り立つと仮定する。
\begin{enumerate}
\item $X$ は準アフィンであり、
$$
\colim_\mathcal{V} \Gamma(V, \mathcal{O}_V)
\longrightarrow
\lim \Gamma(Y_n, \mathcal{O}_{Y_n})
$$
が同型である、または
\item $X$ は豊富な可逆加群 $\mathcal{L}$ をもち、
$$
\colim_\mathcal{V} \Gamma(V, \mathcal{L}^{\otimes m})
\longrightarrow
\lim \Gamma(Y_n, \mathcal{L}^{\otimes m}|_{Y_n})
$$
がすべての $m \gg 0$ について同型である、または
\item 任意の $V \in \mathcal{V}$ と任意の有限局所自由な
$\mathcal{O}_V$-加群 $\mathcal{E}$ について写像
$$
\colim_{V' \geq V} \Gamma(V', \mathcal{E}|_{V'})
\longrightarrow
\lim \Gamma(Y_n, \mathcal{E}|_{Y_n})
$$
が同型である。
\end{enumerate}
このとき関手
$$
\colim_\mathcal{V} \textit{F\'Et}_V \to \textit{F\'Et}_Y
$$
は忠実充満である。
\end{lemma}

\begin{proof}
この補題は、補題 \ref{pione-lemma-restriction-fully-faithful-general} と
Algebraic and Formal Geometry, 補題
\ref{algebraization-lemma-completion-fully-faithful-general} から形式的に従う。
\end{proof}






\section{押し出しと基本群}
\label{pione-section-pushouts}

\noindent
主結果は次のとおりである。

\begin{lemma}
\label{pione-lemma-pushout-along-closed-immersion-and-integral}
More on Morphisms, 状況
\ref{more-morphisms-situation-pushout-along-closed-immersion-and-integral}
において、例えば $Z \to Y$ と $Z \to X$ がスキームの閉埋め込みであれば、
圏の同値
$$
\textit{F\'Et}_{Y \amalg_Z X}
\longrightarrow
\textit{F\'Et}_Y
\times_{\textit{F\'Et}_Z}
\textit{F\'Et}_X
$$
が存在する。
\end{lemma}

\begin{proof}
More on Morphisms, 命題
\ref{more-morphisms-proposition-pushout-along-closed-immersion-and-integral}.
により押し出しは存在する。この関手は、押し出し上有限エタールなスキーム $U$ を
基底変換 $Y' = U \times_{Y \amalg_Z X} Y$ と
$X' = U \times_{Y \amalg_Z X} X$、および自然な同型
$Y' \times_Y Z \to X' \times_X Z$（$Z$ 上）へ送ることで与えられる。
この関手が同値であることを示すため、More on Morphisms, 補題
\ref{more-morphisms-lemma-pushout-functor-equivalence-flat}
を用いて準逆関手を構成する。残るのは次を示すことだけである。
分離的、準有限かつエタールな射 $U \to Y \amalg_Z X$ が与えられ、
$X' \to X$ と $Y' \to Y$ が有限ならば、$U \to Y \amalg_Z X$ は有限である。
これは対応する代数的事実（More on Algebra, 補題
\ref{more-algebra-lemma-finite-module-over-fibre-product})
から従う。あるいは
$$
X' \amalg Y' \to U
$$
が全射であり、$X'$ と $Y'$ が $Y \amalg_Z X$ 上固有であることからも分かる
（ここでは More on Morphisms, 命題
\ref{more-morphisms-proposition-pushout-along-closed-immersion-and-integral})
における押し出しの記述を用いる）。そこで Morphisms, 補題
\ref{morphisms-lemma-scheme-theoretic-image-is-proper} を適用すると、
$U$ は $Y \amalg_Z X$ 上固有であると分かる。
準有限かつ固有な射は有限なので（More on Morphisms, 補題
\ref{more-morphisms-lemma-characterize-finite}）、結論を得る。
\end{proof}





\section{穿孔スペクトルの有限エタール被覆 I}
\label{pione-section-pi1-punctured-spec}

\noindent
まず Lefschetz 型の結果をいくつか証明する。

\begin{situation}
\label{pione-situation-local-lefschetz}
 $(A, \mathfrak m)$ をネーター局所環、$f \in \mathfrak m$ とする。
$X = \Spec(A)$ および $X_0 = \Spec(A/fA)$ と置き、
$U = X \setminus \{\mathfrak m\}$ および
$U_0 = X_0 \setminus \{\mathfrak m\}$ を、それぞれ
$A$ および $A/fA$ の穿孔スペクトルとする。
\end{situation}

\noindent
スキーム $X$ に対し、$X$ 上有限エタールなスキームの圏を
$\textit{F\'Et}_X$ と表したことを思い出そう。節
\ref{pione-section-finite-etale} を参照されたい。
状況 \ref{pione-situation-local-lefschetz} では、基底変換関手
$$
\xymatrix{
\textit{F\'Et}_X \ar[d] \ar[r] & \textit{F\'Et}_U \ar[d] \\
\textit{F\'Et}_{X_0} \ar[r] & \textit{F\'Et}_{U_0}
}
$$
を調べる。多くの場合、右側の垂直矢印は忠実である。

\begin{lemma}
\label{pione-lemma-faithful}
状況 \ref{pione-situation-local-lefschetz} において、次のいずれかが
成り立つと仮定する。
\begin{enumerate}
\item $\dim(A/\mathfrak p) \geq 2$ が、すべての極小素イデアル
$\mathfrak p \subset A$ で $f \not \in \mathfrak p$ を満たすものについて成り立つ、または
\item $U$ のすべての連結成分が $U_0$ と交わる。
\end{enumerate}
このとき
$$
\textit{F\'Et}_U \longrightarrow \textit{F\'Et}_{U_0},\quad
V \longmapsto V_0 = V \times_U U_0
$$
は忠実関手である。
\end{lemma}

\begin{proof}
場合 (2) は補題 \ref{pione-lemma-restriction-faithful} から直ちに従う。
仮定 (1) のもとでは $U$ のすべての既約成分が $U_0$ と交わる。
Algebra, 補題 \ref{algebra-lemma-one-equation} を参照されたい。
したがって (1) は (2) から従う。
\end{proof}

\noindent
より興味深い結果を証明する前に、いくつかの補題が必要である。

\begin{lemma}
\label{pione-lemma-fill-in-missing}
状況 \ref{pione-situation-local-lefschetz} において、$V \to U$ を有限射とする。
$A^\wedge$ を $\mathfrak m$-進の $A$ の完備化とし、
$X' = \Spec(A^\wedge)$ と置く。また $U'$ と $V'$ をそれぞれ
$U$ と $V$ の $X'$ への基底変換とする。有限射 $Y' \to X'$ が
$V' = Y' \times_{X'} U'$ を満たすならば、有限射 $Y \to X$ で
$V = Y \times_X U$ および $Y' = Y \times_X X'$ を満たすものが存在する。
\end{lemma}

\begin{proof}
これは More on Algebra, 命題
\ref{more-algebra-proposition-equivalence} の直接的な応用である。
実際、$f_1, \ldots, f_t$ を $\mathfrak m$ の生成元として選ぶ。
各 $i$ に対して $V \times_U D(f_i) = \Spec(B_i)$ と書く。
$1 \leq i, j \leq n$ に対し、同型
$\alpha_{ij} : (B_i)_{f_j} \to (B_j)_{f_i}$ を $A_{f_if_j}$-代数の間に得る。
両方のスペクトルが
$V \times_U D(f_if_j)$ を表すからである。
$Y' = \Spec(B')$ と書く。$V \times_U U' = Y \times_{X'} U'$ なので、
同型 $\alpha_i : B'_{f_i} \to B_i \otimes_A A^\wedge$ を得る。
直接的な議論により、$(B', B_i, \alpha_i, \alpha_{ij})$ は
$\text{Glue}(A \to A^\wedge, f_1, \ldots, f_t)$ の対象である。
More on Algebra, 注意 \ref{more-algebra-remark-glueing-data} を参照されたい。
上で引用した命題を適用し（代数構造を得るため More on Algebra, 注意
\ref{more-algebra-remark-formal-glueing-algebras} を用いる）、$A$-代数 $B$ で
$\text{Can}(B)$ が $(B', B_i, \alpha_i, \alpha_{ij})$ と同型になるものを得る。
$Y = \Spec(B)$ と置くと、$Y \to X$ は射であり、求める整合的な同型
$V \cong Y \times_X U$ および $Y' = Y \times_X X'$ を伴うことが分かる。
\end{proof}

\begin{lemma}
\label{pione-lemma-fully-faithful-henselian-completion}
状況 \ref{pione-situation-local-lefschetz} において、$A$ はヘンゼル的、
またはより一般に $(A, (f))$ はヘンゼル対であると仮定する。
$A^\wedge$ を $\mathfrak m$-進の $A$ の完備化とし、
$X' = \Spec(A^\wedge)$ と置く。また $U'$ と $U'_0$ をそれぞれ
$U$ と $U_0$ の $X'$ への基底変換とする。
$\textit{F\'Et}_{U'} \to \textit{F\'Et}_{U'_0}$ が忠実充満ならば、
$\textit{F\'Et}_U \to \textit{F\'Et}_{U_0}$ も忠実充満である。
\end{lemma}

\begin{proof}
 $\textit{F\'Et}_{U'} \longrightarrow \textit{F\'Et}_{U'_0}$ が
忠実充満であると仮定する。$X' \to X$ は忠実平坦なので、関手
$V \to V_0 = V \times_U U_0$ が忠実であることは直ちに分かる。
有限エタール被覆の圏が内部 Hom をもつので（補題
\ref{pione-lemma-internal-hom-finite-etale}）、$V$ が $U$ 上有限エタールならば
次を示せば十分である。
$$
\Mor_U(U, V) = \Mor_{U_0}(U_0, V_0)
$$
そこで射 $s_0 : U_0 \to V_0$ が $U_0$ 上に与えられたと仮定し、
射 $s : U \to V$ を $U$ 上に構成する。

\medskip\noindent
仮定により、射 $s' : U' \to V'$ で、その $V'_0$ への制限が
$s'_0$、すなわち $s_0$ の基底変換となるものが存在する。
$V' \to U'$ は有限エタールなので、$V' = s'(U') \amalg W'$ となり、
$W' \to U'$ は有限エタールである。有限射 $Z' \to X'$ で
$W' = Z' \times_{X'} U'$ を満たすものを選ぶ。これは More on Morphisms, 補題
\ref{more-morphisms-lemma-quasi-finite-separated-pass-through-finite}
に述べられた形の Zariski の主定理により可能である
（細部を一つ省略した）。このとき
$$
V' = s'(U') \amalg W' \longrightarrow X' \amalg Z' = Y'
$$
は開埋め込みであり、$V' = Y' \times_{X'} U'$ を満たす。
補題 \ref{pione-lemma-fill-in-missing} により、有限な $Y \to X$ で
$V = Y \times_X U$ および $Y' = Y \times_X X'$ を満たすものを得る。
$Y = \Spec(B)$ と書けば $Y' = \Spec(B \otimes_A A^\wedge)$ である。
このとき $B \otimes_A A^\wedge$ は冪等元 $e'$ をもつ。これは開かつ閉な
部分スキーム $X'$ に対応し、ここで $Y' = X' \amalg Z'$ である。

\medskip\noindent
まず $A$ がヘンゼル的な場合を考える（こちらが少し容易である）。
像 $\overline{e}$、すなわち $e'$ の
$B \otimes_A \kappa(\mathfrak m) = B/\mathfrak mB$ における像は、
冪等元 $e$ を $B$ の中に持ち上げられる。これは $A$ がヘンゼル的だからである
（Algebra, 補題 \ref{algebra-lemma-characterize-henselian} により
$B$ は局所環の積である）。冪等元の持ち上げの一意性から、$e$ は
$e'$ へ写る（ここでは $B \otimes_A A^\wedge$ が局所環の積であることを用いる）。
開かつ閉な部分スキーム $Y_1 \subset Y$ を $e$ に対応するものとする。
このとき $Y_1 \times_X X' = s'(X')$ であり、忠実平坦降下により
$Y_1 \to X$ が同型であることが従い、求める切断を与える。

\medskip\noindent
次に $(A, (f))$ がヘンゼル対の場合を考える。ここでは $s'$ が
$s'_0$ の持ち上げであることを用いる。具体的には、
$Y_{0, 1} \subset Y_0 = Y \times_X X_0$ を
$s_0(U_0) \subset V_0 = Y_0 \times_{X_0} U_0$ の閉包とする。
$X' \to X$ は平坦なので、基底変換 $Y'_{0, 1} \subset Y'_0$ は
$s'_0(U'_0)$ の閉包であり、これは $X'_0 \subset Y'_0$ に等しい
（Morphisms, 補題
\ref{morphisms-lemma-flat-base-change-scheme-theoretic-image} を参照）。
$Y'_0 \to Y_0$ は商写像なので（Morphisms, 補題
\ref{morphisms-lemma-fpqc-quotient-topology}）、$Y_{0, 1}$ は
$Y_0$ の開かつ閉な部分スキームである。対応する冪等元を
$e_0 \in B/fB$ とする。More on Algebra, 補題
\ref{more-algebra-lemma-characterize-henselian-pair}
により、$e_0$ を冪等元 $e \in B$ へ持ち上げられる。
あとは前と同様に結論を得る。
\end{proof}

\noindent
状況 \ref{pione-situation-local-lefschetz} において、制限関手
$\textit{F\'Et}_U \longrightarrow \textit{F\'Et}_{U_0}$
の忠実充満性は、かなり緩やかな仮定のもとで成り立つ。
特に、その仮定から $U$ が連結スキームであるとは限らないが、結論により
$U$ と $U_0$ の連結成分の個数は等しいことが保証される。

\begin{lemma}
\label{pione-lemma-fully-faithful-simple}
状況 \ref{pione-situation-local-lefschetz} において、次を仮定する。
\begin{enumerate}
\item[(a)] $A$ は双対化複体をもつ。
\item[(b)] 対 $(A, (f))$ はヘンゼル的である。
\item[(c)] 次のいずれかが成り立つ。
\begin{enumerate}
\item[(i)] $A_f$ は $(S_2)$ を満たし、$X$ の既約成分で
$X_0$ に含まれないものはすべて次元 $\geq 3$ である。
\item[(ii)] 任意の素イデアル $\mathfrak p \subset A$ で
$f \not \in \mathfrak p$ を満たすものについて
$\text{depth}(A_\mathfrak p) + \dim(A/\mathfrak p) > 2$.
\end{enumerate}
\end{enumerate}
このとき制限関手
$\textit{F\'Et}_U \longrightarrow \textit{F\'Et}_{U_0}$
は忠実充満である。
\end{lemma}

\begin{proof}
 $A'$ を $\mathfrak m$-進の $A$ の完備化とする。仮定が $A'$ に対しても
成り立つことを示す。条件 (a) と (b) については明らかである。条件 (c)(ii) は
Local Cohomology, 補題 \ref{local-cohomology-lemma-change-completion}
により保存される。次に (c)(i) が成り立つと仮定する。$A$ は普遍カテナリーなので
（Dualizing Complexes, 補題 \ref{dualizing-lemma-universally-catenary}）、
$\Spec(A')$ の既約成分で $V(f)$ に含まれないものはすべて次元
$\geq 3$ である。More on Algebra, 命題
\ref{more-algebra-proposition-ratliff} を参照されたい。
$A \to A'$ は Gorenstein ファイバーをもつ平坦射なので、
$A_f$ が $(S_2)$ を満たすことから $A'_f$ も $(S_2)$ を満たすことが従う。
ここで用いた参照は Dualizing Complexes, 節
\ref{dualizing-section-formal-fibres}、More on Algebra, 節
\ref{more-algebra-section-properties-formal-fibres}、および Algebra, 補題
\ref{algebra-lemma-Sk-goes-up} である。したがって補題
\ref{pione-lemma-fully-faithful-henselian-completion} により、$A$ は
ネーター完備局所環であると仮定してよい。

\medskip\noindent
他の仮定に加えて $A$ が完備局所環であると仮定する。
補題 \ref{pione-lemma-restriction-fully-faithful} により、結果は
Algebraic and Formal Geometry, 補題
\ref{algebraization-lemma-fully-faithful-simple-two} から従う。
\end{proof}

\begin{lemma}
\label{pione-lemma-fully-faithful-minimal}
\begin{reference}
\cite[Corollary 1.11]{Bhatt-local}
\end{reference}
状況 \ref{pione-situation-local-lefschetz} において、次を仮定する。
\begin{enumerate}
\item $H^1_\mathfrak m(A)$ と $H^2_\mathfrak m(A)$ は
$f$ のある冪で零化される。
\item $A$ はヘンゼル的、またはより一般に $(A, (f))$ はヘンゼル対である。
\end{enumerate}
このとき制限関手
$\textit{F\'Et}_U \longrightarrow \textit{F\'Et}_{U_0}$
は忠実充満である。
\end{lemma}

\begin{proof}
補題 \ref{pione-lemma-fully-faithful-henselian-completion} により、
$A$ はネーター完備局所環であると仮定してよい
（仮定は保たれる。Dualizing Complexes, 補題
\ref{dualizing-lemma-torsion-change-rings} を用いる）。
補題 \ref{pione-lemma-restriction-fully-faithful} により、結果は
Algebraic and Formal Geometry, 補題
\ref{algebraization-lemma-fully-faithful-alternative} から従う。
\end{proof}

\begin{lemma}
\label{pione-lemma-fully-faithful}
状況 \ref{pione-situation-local-lefschetz} において、$A$ の深さが $\geq 3$ であり、
$A$ はヘンゼル的、またはより一般に $(A, (f))$ はヘンゼル対であると仮定する。
このとき制限関手
$\textit{F\'Et}_U \to \textit{F\'Et}_{U_0}$
は忠実充満である。
\end{lemma}

\begin{proof}
深さの仮定により
$H^1_\mathfrak m(A) = H^2_\mathfrak m(A) = 0$ となる。
Dualizing Complexes, 補題 \ref{dualizing-lemma-depth} を参照されたい。
したがって補題 \ref{pione-lemma-fully-faithful-minimal} が適用できる。
\end{proof}



























\section{局所的な場合の純性 I}
\label{pione-section-local-purity}

\noindent
$(A, \mathfrak m)$ をネーター局所環とする。$X = \Spec(A)$ と置き、
$U = X \setminus \{\mathfrak m\}$ を穿孔スペクトルとする。
{\it $(A, \mathfrak m)$ に対して純性が成り立つ}とは、制限関手
$$
\textit{F\'Et}_X \longrightarrow \textit{F\'Et}_U
$$
が本質的全射であることをいう。本節では、この問題について、$X$ から超曲面
$X_0$（$X$ 内）へ移ると何が変わるかを明らかにする。言い換えれば、穿孔スペクトルの
基本群に対する一種の局所 Lefschetz 性質を調べる。これらの結果は、局所純性に関する
主結果、すなわち補題 \ref{pione-lemma-local-purity}、命題
\ref{pione-proposition-purity-complete-intersection}、および命題
\ref{pione-proposition-purity-smooth-over-depth2} の証明を次元に関する帰納法で進める際に
有用となる。

\begin{lemma}
\label{pione-lemma-sections-over-punctured-spec}
$(A, \mathfrak m)$ をネーター局所環とする。$X = \Spec(A)$ と置き、
$U = X \setminus \{\mathfrak m\}$ とする。$\pi : Y \to X$ を有限射とし、
$\text{depth}(\mathcal{O}_{Y, y}) \geq 2$ がすべての閉点
$y \in Y$ について成り立つと仮定する。
このとき $Y$ は $B = \mathcal{O}_Y(\pi^{-1}(U))$ のスペクトルである。
\end{lemma}

\begin{proof}
$V = \pi^{-1}(U)$ と置き、$\pi' : V \to U$ で $\pi$ の制限を表す。
次の $\mathcal{O}_X$-加群写像を考える。
$$
\pi_*\mathcal{O}_Y \longrightarrow j_*\pi'_*\mathcal{O}_V
$$
ここで $j : U \to X$ は包含射である。この写像には Divisors, 補題
\ref{divisors-lemma-depth-2-hartog} が適用できると主張する。
そうであれば $B = \Gamma(Y, \mathcal{O}_Y)$ となり、補題が従う。
$x \in X$ を閉点とする。次を示せば十分である。
$\text{depth}((\pi_*\mathcal{O}_Y)_x) \geq 2$.
$y_1, \ldots, y_n \in Y$ を $x$ へ写る点とする。
Algebra, 補題 \ref{algebra-lemma-depth-goes-down-finite} により、
$\text{depth}(\mathcal{O}_{Y, y_i}) \geq 2$ が $i = 1, \ldots, n$ について
成り立つことを示せば十分である。
これは補題の仮定そのものなので、証明が終わる。
\end{proof}

\begin{lemma}
\label{pione-lemma-reformulate-purity}
$(A, \mathfrak m)$ をネーター局所環とする。$X = \Spec(A)$ と置き、
$U = X \setminus \{\mathfrak m\}$ とする。$V$ は $U$ 上有限エタールとする。
$A$ の深さが $\geq 2$ であると仮定する。次は同値である。
\begin{enumerate}
\item $V = Y \times_X U$ となる有限エタール射 $Y \to X$ が存在する。
\item $B = \Gamma(V, \mathcal{O}_V)$ は $A$ 上有限エタールである。
\end{enumerate}
\end{lemma}

\begin{proof}
与えられた有限エタール射を $\pi : V \to U$ と表す。
(1) の $Y$ が存在すると仮定する。$x \in X$ を $\mathfrak m$ に対応する点とする。
$y \in Y$ を $x$ へ写る点とする。次を主張する。
$\text{depth}(\mathcal{O}_{Y, y}) \geq 2$.
これは $Y \to X$ がエタールであり、したがって
$A = \mathcal{O}_{X, x}$ と $\mathcal{O}_{Y, y}$ の深さが等しいことから従う
（Algebra, 補題 \ref{algebra-lemma-apply-grothendieck}）。
よって補題 \ref{pione-lemma-sections-over-punctured-spec} が適用でき、
$Y = \Spec(B)$ となる。

\medskip\noindent
含意 (2) $\Rightarrow$ (1) はより容易なので、細部を省略する。
\end{proof}

\begin{lemma}
\label{pione-lemma-reformulate-purity-normal}
$(A, \mathfrak m)$ をネーター局所環とする。$X = \Spec(A)$ と置き、
$U = X \setminus \{\mathfrak m\}$ とする。$A$ は次元 $\geq 2$ の正規環であると仮定する。
関手
$$
\textit{F\'Et}_U \longrightarrow
\left\{
\begin{matrix}
\text{有限正規 }A\text{-代数 }B\text{ であって} \\
\Spec(B) \to X\text{ が }U\text{ 上エタールであるもの}
\end{matrix}
\right\},
\quad
V \longmapsto \Gamma(V, \mathcal{O}_V)
$$
は同値である。さらに、$V = Y \times_X U$ となる有限エタール射 $Y \to X$ が
存在するための必要十分条件は、$B = \Gamma(V, \mathcal{O}_V)$ が $A$ 上有限エタール
であることである。
\end{lemma}

\begin{proof}
$\text{depth}(A) \geq 2$ であることに注意する。これは $A$ が正規だからである
（正規性に関する Serre の判定法、Algebra, 補題
\ref{algebra-lemma-criterion-normal}）。したがって最後の主張は補題
\ref{pione-lemma-reformulate-purity} から従う。有限エタール射 $\pi : V \to U$ が与えられたとき、
$B = \Gamma(V, \mathcal{O}_V)$ と置く。$B$ が正規で $A$ 上有限であることを示せれば、
表示した関手が得られる。明らかな擬逆関手があるので、示すべきことはこれですべてである。

\medskip\noindent
$A$ は正規なので、スキーム $V$ は正規である
（Descent, 補題 \ref{descent-lemma-normal-local-smooth}）。したがって $V$ は整スキームの
有限非交和である（Properties, 補題 \ref{properties-lemma-normal-Noetherian}）。
よって $V$ は整であると仮定してよい。この場合、関数体 $L$、すなわち $V$ の関数体
（Morphisms, 節 \ref{morphisms-section-rational-maps}）は $A$ の分数体の有限分離拡大である
（これは $V \to U$ の一般ファイバーを調べ、Morphisms, 補題
\ref{morphisms-lemma-etale-over-field} を用いることから分かる）。Algebra, 補題
\ref{algebra-lemma-Noetherian-normal-domain-finite-separable-extension} により、
$B' \subset L$ を $A$ の $L$ における整閉包とすると、これは $A$ 上有限である。More on Algebra, 補題
\ref{more-algebra-lemma-integral-closure-reflexive} により $B'$ は反射的 $A$-加群であり、
したがって More on Algebra, 補題 \ref{more-algebra-lemma-reflexive-over-normal} により
$\text{depth}_A(B') \geq 2$ である。

\medskip\noindent
$f \in \mathfrak m$ とする。このとき $B_f = \Gamma(V \times_U D(f), \mathcal{O}_V)$ である
（Properties, 補題 \ref{properties-lemma-invert-f-sections}）。したがって $B'_f = B_f$ である。
実際、上で見たように $B_f$ は正規で、$A_f$ 上有限かつ分数体 $L$ をもつ。よって
$V = \Spec(B') \times_X U$ となる。さらに、補題
\ref{pione-lemma-sections-over-punctured-spec} から $B = B'$ が従う。この補題は
$\Spec(B') \to X$ に適用する。これを適用できるのは、局所化 $B'_{\mathfrak m'}$、すなわち
$B'$ の極大イデアル $\mathfrak m' \subset B'$ で $\mathfrak m$ の上にあるものにおける
局所化が、Algebra, 補題 \ref{algebra-lemma-depth-goes-down-finite} と
前段落の深さに関する注意により、深さ $\geq 2$ をもつからである。
\end{proof}

\begin{lemma}
\label{pione-lemma-purity-and-completion}
$(A, \mathfrak m)$ をネーター局所環とする。$X = \Spec(A)$ と置き、
$U = X \setminus \{\mathfrak m\}$ とする。$V$ は $U$ 上有限エタールとする。
$A^\wedge$ を $\mathfrak m$ に関する $A$ の完備化とし、$X' = \Spec(A^\wedge)$ と置く。
$U'$ と $V'$ をそれぞれ $U$ と $V$ の $X'$ への基底変換とする。次は同値である。
\begin{enumerate}
\item $V = Y \times_X U$ となる有限エタール射 $Y \to X$ が存在する。
\item $V' = Y' \times_{X'} U'$ となる有限エタール射 $Y' \to X'$ が存在する。
\end{enumerate}
\end{lemma}

\begin{proof}
含意 (1) $\Rightarrow$ (2) は、解 $Y \to X$ の基底変換を取ることから従う。
$Y' \to X'$ を (2) のものとする。補題 \ref{pione-lemma-fill-in-missing} により、有限射
$Y \to X$ であって $V = Y \times_X U$ かつ $Y' = Y \times_X X'$ となるものを取れる。
降下により $Y \to X$ は有限エタールである
（Algebra, 補題 \ref{algebra-lemma-descend-properties-modules} および
\ref{algebra-lemma-etale}）。これで証明が終わる。
\end{proof}

\noindent
次の二つの補題の要点は、その仮定から $A$ の深さが $\geq 3$ であるとは限らないことである。
例えば、$A$ が次元 $\geq 3$ の完備正規局所整域で、$f \in \mathfrak m$ が非零ならば、
これらの仮定は満たされる。

\begin{lemma}
\label{pione-lemma-lift-simple}
状況 \ref{pione-situation-local-lefschetz} のもとで、$V$ を $U$ 上有限エタールとする。次を仮定する。
\begin{enumerate}
\item[(a)] $A$ は双対化複体をもつ。
\item[(b)] 対 $(A, (f))$ は Hensel 対である。
\item[(c)] 次のいずれかが成り立つ。
\begin{enumerate}
\item[(i)] $A_f$ は $(S_2)$ であり、$X$ のうち $X_0$ に含まれない各既約成分の
次元は $\geq 3$ である。
\item[(ii)] すべての素イデアル $\mathfrak p \subset A$ について、
$f \not \in \mathfrak p$ ならば
$\text{depth}(A_\mathfrak p) + \dim(A/\mathfrak p) > 2$ である。
\end{enumerate}
\item[(d)] $V_0 = V \times_U U_0$ は $Y_0 \times_{X_0} U_0$ に等しくなる有限エタール射
$Y_0 \to X_0$ が存在する。
\end{enumerate}
このとき $V = Y \times_X U$ となる有限エタール射 $Y \to X$ が存在する。
\end{lemma}

\begin{proof}
補題 \ref{pione-lemma-purity-and-completion} を用いて完備な場合に帰着する。
（仮定も引き継がれる。補題 \ref{pione-lemma-fully-faithful-simple} の証明を参照。）

\medskip\noindent
完備な場合、$Y_0 \to X_0$ を有限エタール射 $Y \to X$ に持ち上げられる。これは
More on Algebra, 補題 \ref{more-algebra-lemma-finite-etale-equivalence} による。また、
More on Algebra, 補題 \ref{more-algebra-lemma-complete-henselian} により
$(A, fA)$ は Hensel 対であることに注意する。そこで補題
\ref{pione-lemma-fully-faithful-simple} を用いると、$V$ は $Y \times_X U$ と同型である。
これで証明が終わる。
\end{proof}

\begin{lemma}
\label{pione-lemma-lift-purity-general}
状況 \ref{pione-situation-local-lefschetz} のもとで、$V$ を $U$ 上有限エタールとする。次を仮定する。
\begin{enumerate}
\item $H^1_\mathfrak m(A)$ と $H^2_\mathfrak m(A)$ は $f$ のある冪で零化される。
\item $V_0 = V \times_U U_0$ は $Y_0 \times_{X_0} U_0$ に等しくなる有限エタール射
$Y_0 \to X_0$ が存在する。
\end{enumerate}
このとき $V = Y \times_X U$ となる有限エタール射 $Y \to X$ が存在する。
\end{lemma}

\begin{proof}
補題 \ref{pione-lemma-purity-and-completion} を用いて完備な場合に帰着する。
（仮定も引き継がれる。Dualizing Complexes, 補題
\ref{dualizing-lemma-torsion-change-rings} を用いる。）

\medskip\noindent
完備な場合、$Y_0 \to X_0$ を有限エタール射 $Y \to X$ に持ち上げられる。これは
More on Algebra, 補題 \ref{more-algebra-lemma-finite-etale-equivalence} による。また、
More on Algebra, 補題 \ref{more-algebra-lemma-complete-henselian} により
$(A, fA)$ は Hensel 対であることに注意する。そこで補題
\ref{pione-lemma-fully-faithful-minimal} を用いると、$V$ は $Y \times_X U$ と同型である。
これで証明が終わる。
\end{proof}

\begin{lemma}
\label{pione-lemma-lift-purity}
状況 \ref{pione-situation-local-lefschetz} のもとで、$V$ を $U$ 上有限エタールとする。次を仮定する。
\begin{enumerate}
\item $A$ の深さは $\geq 3$ である。
\item $V_0 = V \times_U U_0$ は $Y_0 \times_{X_0} U_0$ に等しくなる有限エタール射
$Y_0 \to X_0$ が存在する。
\end{enumerate}
このとき $V = Y \times_X U$ となる有限エタール射 $Y \to X$ が存在する。
\end{lemma}

\begin{proof}
深さの仮定から $H^1_\mathfrak m(A) = H^2_\mathfrak m(A) = 0$ が従う。
Dualizing Complexes, 補題 \ref{dualizing-lemma-depth} を参照。したがって補題
\ref{pione-lemma-lift-purity-general} を適用できる。
\end{proof}










\section{分岐軌跡の純性}
\label{pione-section-purity}

\noindent
有限局所自由射の判別式を用いる。Discriminants, 節
\ref{discriminant-section-discriminant} を参照。

\begin{lemma}
\label{pione-lemma-find-point-codim-1}
$(A, \mathfrak m)$ を $\dim(A) \geq 1$ であるネーター局所環とする。
$f \in \mathfrak m$ とする。このとき、$\mathfrak p \in V(f)$ で
$\dim(A_\mathfrak p) = 1$ を満たすものが存在する。
\end{lemma}

\begin{proof}
$\dim(A)$ に関する帰納法を用いる。$\dim(A) = 1$ ならば
$\mathfrak p = \mathfrak m$ とすればよい。$\dim(A) > 1$ ならば、
$Z \subset \Spec(A)$ を次元 $> 1$ の既約成分とする。このとき
$V(f) \cap Z$ の次元は $> 0$ である
（Algebra, 補題 \ref{algebra-lemma-one-equation}）。素イデアル
$\mathfrak q \in V(f) \cap Z$ で $\mathfrak q \not = \mathfrak m$ となり、
$A$ の穿孔スペクトルの閉点に対応するものを選ぶ。
これは Properties, 補題 \ref{properties-lemma-complement-closed-point-Jacobson} により可能である。
このとき $\mathfrak q$ は $Z$ の一般点ではない。したがって
$0 < \dim(A_\mathfrak q) < \dim(A)$ かつ $f \in \mathfrak q A_\mathfrak q$ である。
次元に関する帰納法により、$f \in \mathfrak p \subset A_\mathfrak q$ で
$\dim((A_\mathfrak q)_\mathfrak p) = 1$ を満たすものを取れる。このとき
$\mathfrak p \cap A$ が所望のものである。
\end{proof}

\begin{lemma}
\label{pione-lemma-ramification-quasi-finite-flat}
$f : X \to Y$ を局所ネータースキームの射とし、$x \in X$ とする。次を仮定する。
\begin{enumerate}
\item $f$ は平坦である。
\item $f$ は $x$ において準有限である。
\item $x$ は $X$ のどの既約成分の一般点でもない。
\item 特殊化 $x' \leadsto x$ で $\dim(\mathcal{O}_{X, x'}) = 1$ を満たすものに対して、
$f$ は $x'$ において不分岐である。
\end{enumerate}
このとき $f$ は $x$ においてエタールである。
\end{lemma}

\begin{proof}
$f$ が不分岐である点の集合は $f$ がエタールである点の集合に等しく、しかもこの集合は開である。
Morphisms, 定義 \ref{morphisms-definition-unramified}、
\ref{morphisms-definition-etale} および補題
\ref{morphisms-lemma-flat-unramified-etale} を参照。$f$ が $x$ においてエタールであることは、
基底と終域の双方でエタール局所的に調べてよい
（Descent, 補題 \ref{descent-lemma-descending-property-etale} および
\ref{descent-lemma-etale-etale-local-source}）。したがって More on Morphisms, 補題
\ref{more-morphisms-lemma-etale-makes-quasi-finite-finite-at-point} を適用し、
$f : X \to Y$ は有限で、$x$ は $X$ において $y = f(x)$ の上にある唯一の点であると仮定できる。
すると $f$ は有限局所自由である
（Morphisms, 補題 \ref{morphisms-lemma-finite-flat}）。

\medskip\noindent
$f$ は有限局所自由で、$x$ は $X$ において $y = f(x)$ の上にある唯一の点であると仮定する。
Discriminants, 補題 \ref{discriminant-lemma-discriminant} により、局所主閉部分スキーム
$D_\pi \subset Y$ であって、$y' \in D_\pi$ であることと、$x' \in X$ で
$f(x') = y'$ かつ $f$ が $x'$ において分岐するものが存在することとが同値になるものを得る。
したがって $y \not \in D_\pi$ を示せばよい。矛盾を導くため $y \in D_\pi$ と仮定する。

\medskip\noindent
条件 (3) により $\dim(\mathcal{O}_{X, x}) \geq 1$ である。Algebra, 補題
\ref{algebra-lemma-dimension-base-fibre-equals-total} により
$\dim(\mathcal{O}_{X, x}) = \dim(\mathcal{O}_{Y, y})$ である。補題
\ref{pione-lemma-find-point-codim-1} により、$y' \in D_\pi$ で $y$ へ特殊化し、
$\dim(\mathcal{O}_{Y, y'}) = 1$ を満たすものを取れる。
$x' \in X$ で $f(x') = y'$ かつ $f$ がそこで分岐するものを選ぶ。$f$ は有限なので閉写像であり、
したがって $x' \leadsto x$ である。先と同様に
$\dim(\mathcal{O}_{X, x'}) = \dim(\mathcal{O}_{Y, y'}) = 1$ となる。
これは条件 (4) に反する。
\end{proof}

\begin{lemma}
\label{pione-lemma-local-purity}
$(A, \mathfrak m)$ を次元 $d \geq 2$ の正則局所環とする。
$X = \Spec(A)$ および $U = X \setminus \{\mathfrak m\}$ と置く。このとき、
\begin{enumerate}
\item 関手 $\textit{F\'Et}_X \to \textit{F\'Et}_U$ は本質的全射である。すなわち、
$A$ に対して純性が成り立つ。
\item $A \to B$ が有限で $B$ が正規であり、穿孔スペクトル上に有限エタール射を
誘導するならば、この射はエタールである。
\end{enumerate}
\end{lemma}

\begin{proof}
Algebra, 補題 \ref{algebra-lemma-regular-normal} により、正則局所環は正規であることを
思い出そう。したがって補題 \ref{pione-lemma-reformulate-purity-normal} により (1) と (2) は
同値である。$d$ に関する帰納法で補題を証明する。

\medskip\noindent
$d = 2$ の場合。この場合 $A \to B$ は平坦である。実際、Algebra, 命題
\ref{algebra-proposition-going-down-normal-integral} により $A \to B$ に対して下降定理が
成り立つ。すると Algebra, 補題 \ref{algebra-lemma-dimension-base-fibre-equals-total} により、
$\dim(B_{\mathfrak m'}) = 2$ が、すべての極大イデアル
$\mathfrak m' \subset B$ について成り立つ。さらに Algebra, 補題
\ref{algebra-lemma-criterion-normal} により $B_{\mathfrak m'}$ は Cohen--Macaulay である。
したがって、ここが重要な点だが、Algebra, 補題
\ref{algebra-lemma-CM-over-regular-flat} を適用して $A \to B_{\mathfrak m'}$ が平坦であると
分かる。次に Algebra, 補題 \ref{algebra-lemma-flat-localization} により $A \to B$ は平坦である。
よって補題 \ref{pione-lemma-ramification-quasi-finite-flat}（または、より容易な Discriminants, 補題
\ref{discriminant-lemma-discriminant} を用いて直接議論してもよい）を適用すると、
$A \to B$ はエタールである。

\medskip\noindent
$d \geq 3$ の場合。$V \to U$ を有限エタールとする。
$f \in \mathfrak m_A$ かつ $f \not \in \mathfrak m_A^2$ となる元を取る。このとき
$A/fA$ は次元 $d - 1 \geq 2$ の正則局所環である。Algebra, 補題
\ref{algebra-lemma-regular-ring-CM} を参照。$U_0$ を $A/fA$ の穿孔スペクトルとし、
$V_0 = V \times_U U_0$ と置く。補題 \ref{pione-lemma-lift-purity} により、$V_0$ が関手
$\textit{F\'Et}_{\Spec(A/fA)} \to \textit{F\'Et}_{U_0}$ の本質的像に属することを
示せば十分である。これは帰納法の仮定から従う。
\end{proof}

\begin{lemma}[分岐軌跡の純性]
\label{pione-lemma-purity}
\begin{reference}
\cite{Nagata-Purity} および \cite[Exp. X, Thm. 3.1]{SGA1}
\end{reference}
\begin{history}
この結果は、まず Zariski によって幾何学的な形で述べられ、証明された
\cite{Zariski-Purity}。Nagata による非完全基礎体への一般化は、同じ巻の
Proceedings of the National Academy of Sciences of the United States of America の
次の論文として \cite{Nagata-Remarks-Purity} に発表された。その翌年、Nagata は
ネーター局所環に対する結果を \cite{Nagata-Purity} で証明した。その証明は、完備局所整域に
対する Bertini の定理である Chow の結果を用いる。\cite{Chow-Bertini} を参照。
Bertini の諸定理の歴史は Kleiman の歴史的論文 \cite{Kleiman-Bertini} で論じられている。
数年後、Auslander によってまったく異なる証明が得られた。
\cite{Auslander-Purity} を参照。
\end{history}
$f : X \to Y$ を局所ネータースキームの射とする。$x \in X$ とし、
$y = f(x)$ と置く。次を仮定する。
\begin{enumerate}
\item $\mathcal{O}_{X, x}$ は正規である。
\item $\mathcal{O}_{Y, y}$ は正則である。
\item $f$ は $x$ において準有限である。
\item $\dim(\mathcal{O}_{X, x}) = \dim(\mathcal{O}_{Y, y}) \geq 1$
\item 特殊化 $x' \leadsto x$ で $\dim(\mathcal{O}_{X, x'}) = 1$ を満たすものに対して、
$f$ は $x'$ において不分岐である。
\end{enumerate}
このとき $f$ は $x$ においてエタールである。
\end{lemma}

\begin{proof}
$d = \dim(\mathcal{O}_{X, x}) = \dim(\mathcal{O}_{Y, y})$ に関する帰納法で証明する。

\medskip\noindent
$d = 1$ の場合は特に問題がない。この場合、$f$ は $x$ において不分岐であり、
$\mathcal{O}_{Y, y}$ は離散付値環であると仮定している
（Algebra, 補題 \ref{algebra-lemma-characterize-dvr}）。すると
$\mathcal{O}_{X, x}$ は $\mathcal{O}_{Y, y}$ 上平坦である
（そうでなければ、射は $x$ において準有限でない）。したがって $f$ は $x$ において平坦である。
平坦 $+$ 不分岐はエタールなので結論を得る（いくつかの細部は省略する）。

\medskip\noindent
$d \geq 2$ の場合。$d$ に関する帰納法を用いて、補題
\ref{pione-lemma-local-purity} で論じた場合に帰着する。$f$ が $x$ においてエタールであることは、
基底と終域の双方でエタール局所的に調べてよい
（Descent, 補題 \ref{descent-lemma-descending-property-etale} および
\ref{descent-lemma-etale-etale-local-source}）。したがって More on Morphisms, 補題
\ref{more-morphisms-lemma-etale-makes-quasi-finite-finite-at-point} を適用し、
$f : X \to Y$ は有限で、$x$ は $X$ において $y$ の上にある唯一の点であると仮定できる。
ここでは、局所環のエタール拡大が次元、正規性、正則性を変えないことを用いた。
More on Algebra, 節 \ref{more-algebra-section-permanence-etale} および
\'Etale Morphisms, 節 \ref{etale-section-properties-permanence} を参照。

\medskip\noindent
次に、$\Spec(\mathcal{O}_{Y, y})$ によって基底変換し、$Y$ は正則局所環のスペクトルであると
仮定できる。このとき $X = \Spec(\mathcal{O}_{X, x})$ となる。実際、$X$ のすべての点は
必ず $x$ へ特殊化する。

\medskip\noindent
環準同型 $\mathcal{O}_{Y, y} \to \mathcal{O}_{X, x}$ は有限であり、次元が等しいことから
必ず単射である。Algebra, 命題 \ref{algebra-proposition-going-down-normal-integral} により、
$\mathcal{O}_{Y, y} \to \mathcal{O}_{X, x}$ に対して下降定理が成り立つ
（Algebra, 補題 \ref{algebra-lemma-regular-normal} によれば正則環は正規環であることも用いる）。
$x' \in X$ かつ $x' \not = x$ で、像が $y' = f(x')$ となる点を選ぶ。
すると $\mathcal{O}_{X, x'}$ は正規整域の局所化なので正規である。同様に
$\mathcal{O}_{Y, y'}$ は正則である（Algebra, 補題
\ref{algebra-lemma-localization-of-regular-local-is-regular} を参照）。Algebra, 補題
\ref{algebra-lemma-dimension-base-fibre-equals-total} により
$\dim(\mathcal{O}_{X, x'}) = \dim(\mathcal{O}_{Y, y'})$ である
（上で下降定理を確認した）。これらの次元はもちろん $d$ より真に小さい。実際、
$x' \not = x$ である。
$d$ に関する帰納法により、$f$ は $x'$ においてエタールである。

\medskip\noindent
以上により、次の状況に到達した。ネーター局所環の有限局所準同型 $A \to B$ があり、
その次元は $d \geq 2$ で、$A$ は正則、$B$ は正規で、この準同型は穿孔スペクトル上に
有限エタール射 $V \to U$ を誘導する。目標は $A \to B$ がエタールであることを示すことである。
これは補題 \ref{pione-lemma-local-purity} から従い、証明が終わる。
\end{proof}

\noindent
次の補題は、有限エタール射を延長できる最大の開部分集合を求める際に有用なことがある。

\begin{lemma}
\label{pione-lemma-extend-S2}
$j : U \to X$ を局所ネータースキームの開埋め込みとし、
$\text{depth}(\mathcal{O}_{X, x}) \geq 2$ が $x \not \in U$ に対して成り立つと仮定する。
$\pi : V \to U$ を有限エタールとする。このとき、
\begin{enumerate}
\item $\mathcal{B} = j_*\pi_*\mathcal{O}_V$ は反射的連接 $\mathcal{O}_X$-代数である。
$Y = \underline{\Spec}_X(\mathcal{B})$ と置く。
\item $Y \to X$ は、$V = Y \times_X U$ かつ
$\text{depth}(\mathcal{O}_{Y, y}) \geq 2$ が $y \in Y \setminus V$ に対して成り立つ
唯一の有限射である。
\item $Y \to X$ が $y$ においてエタールであるための必要十分条件は、$Y \to X$ が
$y$ において平坦であることである。
\item $Y \to X$ がエタールであるための必要十分条件は、$\mathcal{B}$ が
$\mathcal{O}_X$-加群として有限局所自由であることである。
\end{enumerate}
さらに、(a) $\mathcal{B}$ と $Y \to X$ の構成は、局所ネータースキームの平坦射
$X' \to X$ による基底変換と可換である。(b) $V' \to U'$ が有限エタール射で、
$U \subset U' \subset X$ は開、かつ $V \to U$ と $U$ 上で一致するならば、
同型 $Y' \times_X U' = V'$ が $U'$ 上で一意に存在する。
\end{lemma}

\begin{proof}
$\pi_*\mathcal{O}_V$ は有限局所自由 $\mathcal{O}_U$-加群であり、特に反射的である。
Divisors, 補題 \ref{divisors-lemma-reflexive-S2-extend} により、
$j_*\pi_*\mathcal{O}_V$ は、$X$ 上の反射的連接加群であって
$\pi_*\mathcal{O}_V$ に $U$ 上で制限される唯一のものである。これで (1) が示された。

\medskip\noindent
構成により $Y \times_X U = V$ である。$\mathcal{B}$ は連接なので、$Y \to X$ は有限である。
Divisors, 補題 \ref{divisors-lemma-reflexive-S2} により、
$\text{depth}(\mathcal{B}_x) \geq 2$ が $x \in X \setminus U$ に対して成り立つ。
したがって Algebra, 補題 \ref{algebra-lemma-depth-goes-down-finite} により、
$\text{depth}(\mathcal{O}_{Y, y}) \geq 2$ が $y \in Y \setminus V$ に対して成り立つ。
逆に、$\pi' : Y' \to X$ を有限射で、$V = Y' \times_X U$ かつ
$\text{depth}(\mathcal{O}_{Y', y'}) \geq 2$ が $y' \in Y' \setminus V$ に対して
成り立つものとする。このとき $\pi'_*\mathcal{O}_{Y'}$ は $\pi_*\mathcal{O}_V$ に
$U$ 上で制限され、Algebra, 補題 \ref{algebra-lemma-depth-goes-down-finite} により、
$\text{depth}((\pi'_*\mathcal{O}_{Y'})_x) \geq 2$ が $x \in X \setminus U$ に対して
成り立つ。すると例えば Divisors, 補題
\ref{divisors-lemma-depth-2-hartog} により、$\pi'_*\mathcal{O}_{Y'}$ は
$j_*\pi_*\mathcal{O}_V$ と標準的に同型である。これで (2) が示された。

\medskip\noindent
$Y \to X$ が $y$ においてエタールならば、$Y \to X$ は $y$ において平坦である。
逆に、$Y \to X$ が $y$ において平坦であると仮定する。$y \in V$ ならば、
$Y \to X$ は $y$ においてエタールである。$y \not \in V$ ならば、補題
\ref{pione-lemma-ramification-quasi-finite-flat} の (1)、(2)、(3)、(4) を確認することで、
$Y \to X$ が $y$ においてエタールであると分かる。(1) と (2) は明らかであり、
$\text{depth}(\mathcal{O}_{Y, y}) \geq 2$ なので (3) も明らかである。
$y' \leadsto y$ が特殊化で $\dim(\mathcal{O}_{Y, y'}) = 1$ ならば、$y' \in V$ である。
そうでなければ Algebra, 補題 \ref{algebra-lemma-bound-depth} により、この局所環の深さが
$2$ となって矛盾するからである。したがって $Y \to X$ は $y'$ においてエタールであり、
補題 \ref{pione-lemma-ramification-quasi-finite-flat} の (4) も成り立つ。これで (3) の証明が終わる。

\medskip\noindent
(4) は (3) と次の事実から従う。$((Y \to X)_*\mathcal{O}_Y)_x$ が平坦な
$\mathcal{O}_{X, x}$-加群であるための必要十分条件は、$\mathcal{O}_{Y, y}$ が平坦な
$\mathcal{O}_{X, x}$-加群であることが、すべての $y \in Y$ で $x$ へ写るものについて
成り立つことである。Algebra, 補題
\ref{algebra-lemma-flat-localization} を参照。また、ネーター環上の有限平坦加群は有限局所自由
であることも用いる。Algebra, 補題 \ref{algebra-lemma-finite-projective} および
Algebra, 補題 \ref{algebra-lemma-Noetherian-finite-type-is-finite-presentation} を参照。

\medskip\noindent
補題の最後の主張については、(a) は平坦基底変換から従う。Cohomology of Schemes, 補題
\ref{coherent-lemma-flat-base-change-cohomology} を参照。(b) は (2) の一意性を制限
$Y \times_X U'$ に適用することから従う。
\end{proof}

\begin{lemma}
\label{pione-lemma-extend-pure}
$j : U \to X$ をネータースキームの開埋め込みとし、$\mathcal{O}_{X, x}$ の純性が
すべての $x \not \in U$ に対して成り立つと仮定する。このとき
$$
\textit{F\'Et}_X \longrightarrow \textit{F\'Et}_U
$$
は本質的全射である。
\end{lemma}

\begin{proof}
$V \to U$ を有限エタール射とする。ネーター帰納法により、$V \to U$ を $X$ の
真に大きい開部分集合上の有限エタール射へ延長できれば十分である。
$x \in X \setminus U$ を $X \setminus U$ のある既約成分の一般点とする。
このとき逆像 $U_x$、すなわち $U$ の $\Spec(\mathcal{O}_{X, x})$ における逆像は
$\mathcal{O}_{X, x}$ の穿孔スペクトルである。仮定により、$V_x = V \times_U U_x$ は
有限エタール射 $Y_x \to \Spec(\mathcal{O}_{X, x})$ の $U_x$ への制限である。
Limits, 補題 \ref{limits-lemma-glueing-near-point} により、開部分スキーム
$U \subset U' \subset X$ で $x$ を含むものと、有限表示の射 $V' \to U'$ であって、$U$ への制限が
$V \to U$ を復元し、$\Spec(\mathcal{O}_{X, x})$ への制限が $Y_x$ を復元するものを得る。
最後に、射 $V' \to U'$ は、必要なら Limits, 補題
\ref{limits-lemma-glueing-near-point-properties} により $U'$ をより小さい開部分集合へ縮小すると、
有限エタールになる。
\end{proof}
















\section{穿孔スペクトルの有限エタール被覆 II}
\label{pione-section-pi1-punctured-spec-II}

\noindent
本節では、節 \ref{pione-section-pi1-punctured-spec} で論じた内容のいくつかの変形を証明する。
ネーター局所環 $(A, \mathfrak m)$ と $f \in \mathfrak m$ が与えられているとする。
$X = \Spec(A)$ および $X_0 = \Spec(A/fA)$ と置き、$U = X \setminus \{\mathfrak m\}$ と
$U_0 = X_0 \setminus \{\mathfrak m\}$ をそれぞれ $A$ と $A/fA$ の穿孔スペクトルとする。
これはまさに状況 \ref{pione-situation-local-lefschetz} と同じである。異なるのは、制限関手
$$
\colim_{U_0 \subset U' \subset U\text{ open}} \textit{F\'Et}_{U'}
\longrightarrow
\textit{F\'Et}_{U_0}
$$
を考える点である。言い換えれば、$U_0$ の有限エタール被覆を $U$ 全体へ持ち上げようとはせず、
ある開近傍 $U'$、すなわち $U_0$ の $U$ における開近傍へだけ持ち上げる。

\begin{lemma}
\label{pione-lemma-faithful-general}
状況 \ref{pione-situation-local-lefschetz} のもとで、$U' \subset U$ を $U_0$ を含む開部分集合とする。
$\mathfrak p \subset A$ が $\mathfrak p \in U'$ かつ $\mathfrak p \not \in U_0$ を満たす
極小な素イデアルならば $\dim(A/\mathfrak p) \geq 2$ であると仮定する。このとき
$$
\textit{F\'Et}_{U'} \longrightarrow \textit{F\'Et}_{U_0},\quad
V' \longmapsto V_0 = V' \times_{U'} U_0
$$
は忠実関手である。さらに、この仮定を満たす $U'$ が存在し、任意のより小さい開部分集合
$U'' \subset U'$ で $U_0$ を含むものもこの仮定を満たす。特に、制限関手
$$
\colim_{U_0 \subset U' \subset U\text{ open}} \textit{F\'Et}_{U'}
\longrightarrow
\textit{F\'Et}_{U_0}
$$
は忠実である。
\end{lemma}

\begin{proof}
Algebra, 補題 \ref{algebra-lemma-one-equation} により、$V(\mathfrak p)$ は $U_0$ と交わる。
ここで $\mathfrak p$ は $A$ の任意の素イデアルで $\dim(A/\mathfrak p) \geq 2$ を満たすものとする。
したがって補題 \ref{pione-lemma-restriction-faithful} により、表示した関手は主張にあるような
$U$ に対して忠実である。そのような $U'$ の存在を見るため、$\mathfrak p \subset A$ が
$\mathfrak p \in U$、$\mathfrak p \not \in U_0$、$\dim(A/\mathfrak p) = 1$ を満たすならば、
$\mathfrak p$ は $U$ の閉点に対応し、したがって $V(\mathfrak p) \cap U_0 = \emptyset$ であることに
注意する。よって $U'$ として、$X$ の既約成分で $U_0$ と交わらず次元 $1$ であるものの
補集合を取れる。
\end{proof}

\begin{lemma}
\label{pione-lemma-fully-faithful-general-better}
状況 \ref{pione-situation-local-lefschetz} のもとで、次を仮定する。
\begin{enumerate}
\item $A$ は双対化複体をもち、$f$-進完備である。
\item $X$ のうち $X_0$ に含まれないすべての既約成分の次元は $\geq 3$ である。
\end{enumerate}
このとき制限関手
$$
\colim_{U_0 \subset U' \subset U\text{ open}} \textit{F\'Et}_{U'}
\longrightarrow
\textit{F\'Et}_{U_0}
$$
は充満忠実である。
\end{lemma}

\begin{proof}
基本群の位相的不変性により、証明に際して $A$ をその被約化で置き換えてよい。
補題 \ref{pione-lemma-thickening} を参照。すると結果は補題
\ref{pione-lemma-restriction-fully-faithful-general} および Algebraic and Formal Geometry, 補題
\ref{algebraization-lemma-fully-faithful-general} から従う。
\end{proof}

\begin{lemma}
\label{pione-lemma-fully-faithful-general}
状況 \ref{pione-situation-local-lefschetz} のもとで、次を仮定する。
\begin{enumerate}
\item $A$ は $f$-進完備である。
\item $f$ は非零因子である。
\item $H^1_\mathfrak m(A/fA)$ は有限 $A$-加群である。
\end{enumerate}
このとき制限関手
$$
\colim_{U_0 \subset U' \subset U\text{ open}} \textit{F\'Et}_{U'}
\longrightarrow
\textit{F\'Et}_{U_0}
$$
は充満忠実である。
\end{lemma}

\begin{proof}
補題 \ref{pione-lemma-restriction-fully-faithful-general} および
Algebraic and Formal Geometry, 補題
\ref{algebraization-lemma-fully-faithful-general-alternative} から従う。
\end{proof}








\section{穿孔スペクトルの有限エタール被覆 III}
\label{pione-section-pi1-punctured-spec-III}

\noindent
本節では、状況 \ref{pione-situation-local-lefschetz} のもとで制限関手
$$
\colim_{U_0 \subset U' \subset U\text{ open}} \textit{F\'Et}_{U'}
\longrightarrow
\textit{F\'Et}_{U_0}
$$
が圏同値となる条件を調べる。

\begin{lemma}
\label{pione-lemma-essentially-surjective-general-better}
状況 \ref{pione-situation-local-lefschetz} のもとで、次を仮定する。
\begin{enumerate}
\item $A$ は双対化複体をもち、$f$-進完備である。
\item 次のいずれかが成り立つ。
\begin{enumerate}
\item $A_f$ は $(S_2)$ であり、$X$ のうち $X_0$ に含まれないすべての既約成分の
次元は $\geq 4$ である。
\item $\mathfrak p \not \in V(f)$ かつ
$V(\mathfrak p) \cap V(f) \not = \{\mathfrak m\}$ ならば、
$\text{depth}(A_\mathfrak p) + \dim(A/\mathfrak p) > 3$ である。
\end{enumerate}
\end{enumerate}
このとき制限関手
$$
\colim_{U_0 \subset U' \subset U\text{ open}} \textit{F\'Et}_{U'}
\longrightarrow
\textit{F\'Et}_{U_0}
$$
は同値である。
\end{lemma}

\begin{proof}
補題 \ref{pione-lemma-restriction-equivalence-general} および
Algebraic and Formal Geometry, 補題
\ref{algebraization-lemma-equivalence-better} から従う。
\end{proof}

\begin{lemma}
\label{pione-lemma-essentially-surjective-general}
状況 \ref{pione-situation-local-lefschetz} のもとで、次を仮定する。
\begin{enumerate}
\item $A$ は $f$-進完備である。
\item $f$ は非零因子である。
\item $H^1_\mathfrak m(A/fA)$ および $H^2_\mathfrak m(A/fA)$
は有限 $A$-加群である。
\end{enumerate}
このとき制限関手
$$
\colim_{U_0 \subset U' \subset U\text{ open}} \textit{F\'Et}_{U'}
\longrightarrow
\textit{F\'Et}_{U_0}
$$
は同値である。
\end{lemma}

\begin{proof}
補題 \ref{pione-lemma-restriction-equivalence-general} および
Algebraic and Formal Geometry, 補題
\ref{algebraization-lemma-equivalence} から従う。
\end{proof}

\begin{remark}
\label{pione-remark-combine}
$(A, \mathfrak m)$ を完備ネーター局所環とし、$f \in \mathfrak m$ を非零とする。
$A_f$ は $(S_2)$ であり、$\Spec(A)$ のすべての既約成分の次元は $\geq 4$ であると仮定する。
すると補題 \ref{pione-lemma-essentially-surjective-general-better} により、圏
$$
\colim\nolimits_{U' \subset U\text{ open, }U_0 \subset U}
\text{有限エタール }U'\text{-スキームの圏}
$$
は有限エタール $U_0$-スキームの圏と同値である。例えば、$A$ が次元 $\geq 4$ の正規整域ならば
この条件は成り立つ。
\end{remark}




\section{穿孔スペクトルの有限エタール被覆 IV}
\label{pione-section-pi1-punctured-spec-IV}

\noindent
状況 \ref{pione-situation-local-lefschetz} と同様に $X, X_0, U, U_0$ を取る。
本節では、制限関手
$$
\textit{F\'Et}_U
\longrightarrow
\textit{F\'Et}_{U_0}
$$
が本質的全射となる条件を問う。節 \ref{pione-section-pi1-punctured-spec-III} の結果を用い、
残る隙間を純性によって埋めることでこれを行う。ネーター局所環 $(A, \mathfrak m)$ に対し、
制限関手 $\textit{F\'Et}_X \to \textit{F\'Et}_U$ が本質的全射である場合を考える。
ここで $X = \Spec(A)$ および $U = X \setminus \{\mathfrak m\}$ と置く。このとき、この環では
{\it 純性が成り立つ}というのであった。

\begin{lemma}
\label{pione-lemma-equivalence-better}
状況 \ref{pione-situation-local-lefschetz} のもとで、次を仮定する。
\begin{enumerate}
\item $A$ は双対化複体をもち、$f$-進完備である。
\item 次のいずれかが成り立つ。
\begin{enumerate}
\item $A_f$ は $(S_2)$ であり、$X$ のうち $X_0$ に含まれないすべての既約成分の
次元は $\geq 4$ である。
\item $\mathfrak p \not \in V(f)$ かつ
$V(\mathfrak p) \cap V(f) \not = \{\mathfrak m\}$ ならば、
$\text{depth}(A_\mathfrak p) + \dim(A/\mathfrak p) > 3$ である。
\end{enumerate}
\item 任意の極大イデアル $\mathfrak p \subset A_f$ に対して、$(A_f)_\mathfrak p$ では純性が成り立つ。
\end{enumerate}
このとき制限関手 $\textit{F\'Et}_U \to \textit{F\'Et}_{U_0}$ は本質的全射である。
\end{lemma}

\begin{proof}
$V_0 \to U_0$ を有限エタール射とする。補題
\ref{pione-lemma-essentially-surjective-general-better} により、開部分集合 $U' \subset U$ で
$U_0$ を含むものと有限エタール射 $V' \to U$ が存在し、その $U_0$ への基底変換は
$V_0 \to U_0$ と同型である。$U' \supset U_0$ なので、$U \setminus U'$ は (3) にある
素イデアル $\mathfrak p_1, \ldots, \mathfrak p_n$ に対応する点からなる。仮定により、有限エタール射
$V'_i \to \Spec(A_{\mathfrak p_i})$ であって、$V' \to U'$ と
$U' \times_U \Spec(A_{\mathfrak p_i})$ 上で一致するものを取れる。Limits, 補題
\ref{limits-lemma-glueing-near-closed-point} を $n$ 回適用すると、$V' \to U'$ は
有限エタール射 $V \to U$ に延長されることがわかる。
\end{proof}

\begin{lemma}
\label{pione-lemma-equivalence}
$(A, \mathfrak m)$ をネーター局所環とし、$f \in \mathfrak m$ とする。次を仮定する。
\begin{enumerate}
\item $A$ は $f$-進完備である。
\item $f$ は非零因子である。
\item $H^1_\mathfrak m(A/fA)$ および $H^2_\mathfrak m(A/fA)$ は有限 $A$-加群である。
\item 任意の極大イデアル $\mathfrak p \subset A_f$ に対して、$(A_f)_\mathfrak p$ では純性が成り立つ。
\end{enumerate}
このとき制限関手 $\textit{F\'Et}_U \to \textit{F\'Et}_{U_0}$ は本質的全射である。
\end{lemma}

\begin{proof}
証明は補題 \ref{pione-lemma-equivalence-better} の証明と同じであり、補題
\ref{pione-lemma-essentially-surjective-general} を、補題
\ref{pione-lemma-essentially-surjective-general-better} の代わりに用いればよい。
\end{proof}



\section{局所的な場合の純性 II}
\label{pione-section-local-purity-II}

\noindent
本節は節 \ref{pione-section-local-purity} の続きである。ネーター局所環 $(A, \mathfrak m)$ に対し、
制限関手 $\textit{F\'Et}_X \to \textit{F\'Et}_U$ が本質的全射である場合を考える。
ここで $X = \Spec(A)$ および $U = X \setminus \{\mathfrak m\}$ と置く。このとき、この環では
{\it 純性が成り立つ}というのであった。

\begin{lemma}
\label{pione-lemma-purity-inherited-by-hypersurface-better}
$(A, \mathfrak m)$ をネーター局所環とし、$f \in \mathfrak m$ とする。次を仮定する。
\begin{enumerate}
\item $A$ は双対化複体をもち、$f$-進完備である。
\item 次のいずれかが成り立つ。
\begin{enumerate}
\item $A_f$ は $(S_2)$ であり、$X$ のうち $X_0$ に含まれないすべての既約成分の
次元は $\geq 4$ である。
\item $\mathfrak p \not \in V(f)$ かつ
$V(\mathfrak p) \cap V(f) \not = \{\mathfrak m\}$ ならば、
$\text{depth}(A_\mathfrak p) + \dim(A/\mathfrak p) > 3$ である。
\end{enumerate}
\item 任意の極大イデアル $\mathfrak p \subset A_f$ に対して、$(A_f)_\mathfrak p$ では純性が成り立つ。
\item $A$ では純性が成り立つ。
\end{enumerate}
このとき $A/fA$ では純性が成り立つ。
\end{lemma}

\begin{proof}
$X = \Spec(A)$ と置き、$U = X \setminus \{\mathfrak m\}$ を穿孔スペクトルとする。
同様に $X_0 = \Spec(A/fA)$ および $U_0 = X_0 \setminus \{\mathfrak m\}$ を得る。
$V_0 \to U_0$ を有限エタール射とする。補題 \ref{pione-lemma-equivalence-better} により、
有限エタール射 $V \to U$ で、その $U_0$ への基底変換が $V_0 \to U_0$ と同型であるものを得る。
仮定 (5) により、$V \to U$ は有限エタール射 $Y \to X$ に延長される。すると $Y$ の
$X_0$ への制限が、求める $V_0 \to U_0$ の延長である。
\end{proof}

\begin{lemma}
\label{pione-lemma-purity-inherited-by-hypersurface}
$(A, \mathfrak m)$ をネーター局所環とし、$f \in \mathfrak m$ とする。次を仮定する。
\begin{enumerate}
\item $A$ は $f$-進完備である。
\item $f$ は非零因子である。
\item $H^1_\mathfrak m(A/fA)$ および $H^2_\mathfrak m(A/fA)$ は有限 $A$-加群である。
\item 任意の極大イデアル $\mathfrak p \subset A_f$ に対して、$(A_f)_\mathfrak p$ では純性が成り立つ。
\item $A$ では純性が成り立つ。
\end{enumerate}
このとき $A/fA$ では純性が成り立つ。
\end{lemma}

\begin{proof}
証明は補題 \ref{pione-lemma-purity-inherited-by-hypersurface-better} の証明と同じであり、補題
\ref{pione-lemma-equivalence} を、補題 \ref{pione-lemma-equivalence-better} の代わりに用いればよい。
\end{proof}

\noindent
ここまでの結果を帰納的に用いることで、次元 $\geq 3$ の完全交叉では純性が成り立つことを証明できる。
ネーター局所環の完備化が正則局所環を正則列の生成するイデアルで割った商であるとき、その環を
完全交叉というのであった。Divided Power Algebra, 節 \ref{dpa-section-lci} の議論を参照せよ。

\begin{proposition}
\label{pione-proposition-purity-complete-intersection}
$(A, \mathfrak m)$ をネーター局所環とする。$A$ が次元 $\geq 3$ の完全交叉ならば、
$A$ では純性が成り立つ。すなわち、穿孔スペクトルの任意の有限エタール被覆は延長される。
\end{proposition}

\begin{proof}
補題 \ref{pione-lemma-purity-and-completion} により、$A$ は完備局所環であると仮定してよい。仮定により、
$A = B/(f_1, \ldots, f_r)$ と書ける。ここで $B$ は完備正則局所環であり、
$f_1, \ldots, f_r$ は正則列である。$r$ に関する帰納法で証明を完了する。基底の場合 $r = 0$ は、
次元 $\geq 2$ の正則環に適用できる補題 \ref{pione-lemma-local-purity} から従う。

\medskip\noindent
$A = B/(f_1, \ldots, f_r)$ とし、この命題が $r - 1$ に対して成り立つと仮定する。
$A' = B/(f_1, \ldots, f_{r - 1})$ と置き、$f_r \in A'$ に補題
\ref{pione-lemma-purity-inherited-by-hypersurface} を適用する。これは次の理由で可能である。
条件 (1) は $f_1, \ldots, f_r$ が正則列であることから成り立つ。条件 (2) は $B$、したがって
$A'$ が完備であることから成り立つ。条件 (3) は $A = A'/f_r A'$ が次元
$\dim(A) \geq 3$ の Cohen--Macaulay 環であることから成り立つ。Dualizing Complexes, 補題
\ref{dualizing-lemma-depth} を参照せよ。条件 (4) は帰納法の仮定から成り立つ。実際、
$\dim((A'_{f_r})_\mathfrak p) \geq 3$ である。ここで $\mathfrak p$ は $A'_{f_r}$ の極大素イデアルであり、さらに
$(A'_{f_r})_\mathfrak p = B_\mathfrak q/(f_1, \ldots, f_{r - 1})$ が、ある $\mathfrak q \subset B$ に対して
成り立つ。条件 (5) も帰納法の
仮定から成り立つ。
\end{proof}




\section{局所的な場合の純性 III}
\label{pione-section-local-purity-III}

\noindent
本節では、節 \ref{pione-section-local-purity} および \ref{pione-section-local-purity-II} の議論を続ける。

\begin{lemma}
\label{pione-lemma-fully-faithful-power-series-over-depth2}
$(A, \mathfrak m)$ を深さ $\geq 2$ のネーター局所環とする。
$B = A[[x_1, \ldots, x_d]]$ とし、$d \geq 1$ とする。
$Y = \Spec(B)$ および $Y_0 = V(x_1, \ldots, x_d)$ と置く。
$V \subset Y$ を開部分スキームとし、$V_0 = V \cap Y_0$ が
$Y_0 \setminus \{\mathfrak m_B\}$ に等しいとする。このとき制限関手
$$
\textit{F\'Et}_V \longrightarrow \textit{F\'Et}_{V_0}
$$
は充満忠実である。
\end{lemma}

\begin{proof}
$I = (x_1, \ldots, x_d)$ および $X = \Spec(A)$ と置く。
写像 $Y \to X$ を用いて $Y_0$ を $X$ と同一視すると、$V_0$ は穿孔スペクトル $U$、すなわち
$A$ の穿孔スペクトルと同一視される。このアフィン射によって加群を前進すると
\begin{align*}
\lim_n \Gamma(V_0, \mathcal{O}_V/I^n\mathcal{O}_V)
& =
\lim_n \Gamma(V_0, \mathcal{O}_Y/I^n\mathcal{O}_Y) \\
& =
\lim_n \Gamma(U, \mathcal{O}_U[x_1, \ldots, x_d]/(x_1, \ldots, x_d)^n) \\
& =
\lim_n A[x_1, \ldots, x_d]/(x_1, \ldots, x_d)^n \\
& =
B
\end{align*}
実際、$A$ の深さは $\geq 2$ なので $\Gamma(U, \mathcal{O}_U) = A$ である。Local Cohomology, 補題
\ref{local-cohomology-lemma-finiteness-pushforwards-and-H1-local} を参照せよ。
したがって、補題にある任意の開部分集合 $V \subset Y$ に対し
$$
B = \Gamma(Y, \mathcal{O}_Y) \to \Gamma(V, \mathcal{O}_V) \to
\lim_n \Gamma(V_0, \mathcal{O}_Y/I^n\mathcal{O}_Y) = B
$$
を得る。これより二つの矢印はいずれも同型である（細部を省略する）。
Algebraic and Formal Geometry, 補題 \ref{algebraization-lemma-completion-fully-faithful} により、
$\textit{Coh}(\mathcal{O}_V) \to \textit{Coh}(V, I\mathcal{O}_V)$
は有限局所自由対象からなる充満部分圏上で充満忠実である。したがって補題
\ref{pione-lemma-restriction-fully-faithful} から結論を得る。
\end{proof}

\begin{lemma}
\label{pione-lemma-purity-power-series-over-depth2}
\begin{slogan}
有限エタール被覆に対する Ramanujam--Samuel
\end{slogan}
$(A, \mathfrak m)$ を深さ $\geq 2$ のネーター局所環とする。
$B = A[[x_1, \ldots, x_d]]$ とし、$d \geq 1$ とする。開部分集合
$V \subset Y = \Spec(B)$ が次を含むと仮定する。
\begin{enumerate}
\item 任意の素イデアル $\mathfrak q \subset B$ で、$\mathfrak q \cap A \not = \mathfrak m$ を満たすもの。
\item 素イデアル $\mathfrak m B$。
\end{enumerate}
このとき関手
$
\textit{F\'Et}_Y
\to
\textit{F\'Et}_V
$
は同値である。特に $B$ では純性が成り立つ。
\end{lemma}

\begin{proof}
素イデアル $\mathfrak q \subset B$ で $V$ に含まれないものは $\mathfrak m$ 上にある。この場合
$A \to B_\mathfrak q$ は平坦局所準同型であり、したがって
$\text{depth}(B_\mathfrak q) \geq 2$ である（Algebra, 補題
\ref{algebra-lemma-apply-grothendieck}）。ゆえに補題
\ref{pione-lemma-quasi-compact-dense-open-connected-at-infinity-Noetherian} と Local Cohomology, 補題
\ref{local-cohomology-lemma-depth-2-connected-punctured-spectrum} により、この関手は充満忠実である。

\medskip\noindent
$W \to V$ を有限エタール射とする。$B \to C$ を次を満たす一意的な有限環準同型とする。
すなわち、$\Spec(C) \to Y$ は補題 \ref{pione-lemma-extend-S2} で構成される、$W \to V$ を延長する有限射である。
$C = \Gamma(W, \mathcal{O}_W)$ であることに注意する。

\medskip\noindent
$Y_0 = V(x_1, \ldots, x_d)$ および $V_0 = V \cap Y_0$ と置き、$X = \Spec(A)$ と置く。
写像 $Y \to X$ を用いて $Y_0$ を $X$ と同一視すると、$V_0$ は穿孔スペクトル $U$、すなわち
$A$ の穿孔スペクトルと同一視される。したがって $W_0 = W \times_Y Y_0$ を $U$ 上の有限エタール・スキームとみなせる。このとき
$$
W_0 \times_U (U \times_X Y)
\quad\text{かつ}\quad
W \times_V (U \times_X Y)
$$
は $U \times_X Y$ 上の有限エタール・スキームであり、$V_0$ 上への制限は互いに同型である。
開部分集合 $U \times_X Y$ に補題 \ref{pione-lemma-fully-faithful-power-series-over-depth2} を適用して、同型
$$
W_0 \times_U (U \times_X Y) \longrightarrow W \times_V (U \times_X Y)
$$
を $U \times_X Y$ 上で得る。

\medskip\noindent
$C_0 = \Gamma(W_0, \mathcal{O}_{W_0})$ は有限 $A$-代数であることに注意する。これは
$W_0 \to U \subset X$ に補題 \ref{pione-lemma-extend-S2} を適用すれば従う（上で $B \to C$ に対して
行ったのとまったく同じである）。補題 \ref{pione-lemma-extend-S2} の構成は平坦基底変換および
開集合の変更と両立するので、上の同型は同型
$$
\Psi : C \longrightarrow C_0 \otimes_A B
$$
を誘導し、これは有限 $B$-代数の同型である。一方、$\Spec(C) \to Y$ は、$Y$ の少なくとも一つの点であって
$\mathfrak m \in X$ 上にあるものの上のすべての点でエタールであることがわかっている。
$\Psi$ は同型なので、$\Spec(C_0) \to X$ は
$\mathfrak m$ の上でエタールである（細部を省略する）。これはもちろん、$A \to C_0$ が有限エタールであり、
したがって $B \to C$ も有限エタールであることを意味する。
\end{proof}

\begin{lemma}
\label{pione-lemma-purity-smooth-over-depth2}
$f : X \to S$ をスキームの射とし、$U \subset X$ を開部分スキームとする。次を仮定する。
\begin{enumerate}
\item $f$ は滑らかである。
\item $S$ はネーター的である。
\item $s \in S$ が $\text{depth}(\mathcal{O}_{S, s}) \leq 1$ を満たすならば $X_s = U_s$ である。
\item $U_s \subset X_s$ は任意の $s \in S$ に対して稠密である。
\end{enumerate}
このとき $\textit{F\'Et}_X \to \textit{F\'Et}_U$ は同値である。
\end{lemma}

\begin{proof}
補題 \ref{pione-lemma-quasi-compact-dense-open-connected-at-infinity-Noetherian} と Local Cohomology, 補題
\ref{local-cohomology-lemma-depth-2-connected-punctured-spectrum} により、この関手は充満忠実である
（深さの条件を確認するため、さらに Algebra, 補題 \ref{algebra-lemma-apply-grothendieck} を適用する）。

\medskip\noindent
$\pi : V \to U$ を有限エタール射とし、$Y \to X$ を補題 \ref{pione-lemma-extend-S2} で構成した有限射とする。
$Y \to X$ が有限エタールであることを示さなければならない。すべての点 $x \in X$ について、それが
与えられた点 $s \in S$ へ写る場合にこれを示す際には、ネーター・スキームの平坦射 $S' \to S$ で
$s$ を像に含むものによって基底変換してよい。これは補題 \ref{pione-lemma-extend-S2} の構成が
平坦基底変換と両立することから従う。

\medskip\noindent
$U$ を拡大することにより、$U \subset X$ は $Y \to X$ が有限エタールとなる最大の開部分集合であると
仮定してよい。$Z \subset X$ を $U$ の補集合とする。矛盾を導くため $Z \not = \emptyset$ と仮定する。
$s \in S$ を $Z \to S$ の像に属する点で、その点 $s$ の真の生成点が一つも像に属さないものとする。
そこで $\Spec(\mathcal{O}_{S, s})$ へ基底変換すると、$S = \Spec(A)$ であり、
$(A, \mathfrak m, \kappa)$ は深さ $\geq 2$ のネーター局所環であるとしてよい。また $Z$ は閉ファイバー
$X_s$ に含まれ、$X_s$ の中で至る所稠密でない。閉点 $z \in Z$ を選ぶ。このとき
$\kappa(z)/\kappa$ は有限である（Hilbert の零点定理による。Algebra, 定理
\ref{algebra-theorem-nullstellensatz} を参照）。局所スキームの有限平坦射
$(S', s') \to (S, s)$ で剰余体拡大 $\kappa(z)/\kappa$ を実現するものを選ぶ。Algebra, 補題
\ref{algebra-lemma-finite-free-given-residue-field-extension} を参照せよ。$S' \to S$ で基底変換すれば、
$\kappa(z) = \kappa$ の場合に帰着する。

\medskip\noindent
More on Morphisms, 補題 \ref{more-morphisms-lemma-slice-smooth} により、局所閉部分スキーム
$S' \subset X$ であって $z$ を通り、$S' \to S$ が $z$ においてエタールとなるものが存在する。$S' \to S$ で
基底変換することにより、切断 $\sigma : S \to X$ で $\sigma(s) = z$ を満たすものが存在すると仮定してよい。
$\Spec(B) \subset X$ を $s$ のアフィン近傍として選ぶ。このとき $A \to B$ は滑らかな環準同型であり、切断
$\sigma : B \to A$ をもつ。$I = \Ker(\sigma)$ と置き、$B^\wedge$ を $I$-進完備化、すなわち環 $B$ の
完備化とする。このとき $B^\wedge \cong A[[x_1, \ldots, x_d]]$ が、ある $d \geq 0$ に対して成り立つ。Algebra, 補題
\ref{algebra-lemma-section-smooth} を参照せよ。$d > 0$ であることに注意する。実際、そうでなければ
$X \to S$ は $z$ においてエタールとなり、$z$ は $X_s$ の生成点、したがって仮定 (4) により
$z \in U$ となる。同様に、$d > 0$ ならば、$\mathfrak m B^\wedge$ は $U$ の中へ、射
$\Spec(B^\wedge) \to X$ によって写る。$Y \to X$ が $\Spec(B^\wedge)$ への基底変換後に
有限エタールであることを示せば十分である。$B \to B^\wedge$ は平坦なので
（Algebra, 補題 \ref{algebra-lemma-completion-flat}）、これは補題
\ref{pione-lemma-purity-power-series-over-depth2} および $Y \to X$ の構成の一意性から従う。
\end{proof}

\begin{proposition}
\label{pione-proposition-purity-smooth-over-depth2}
$A \to B$ をネーター局所環の局所準同型とする。$A$ の深さは $\geq 2$ であり、
$A \to B$ は $\mathfrak m_B$-進位相に関して形式的に滑らかであり、
$\dim(B) > \dim(A)$ であると仮定する。開部分集合 $V \subset Y = \Spec(B)$ が次を含むとする。
\begin{enumerate}
\item 任意の素イデアル $\mathfrak q \subset B$ で、$\mathfrak q \cap A \not = \mathfrak m_A$ を満たすもの。
\item 素イデアル $\mathfrak m_A B$。
\end{enumerate}
このとき関手 $\textit{F\'Et}_Y \to \textit{F\'Et}_V$ は同値である。特に $B$ では純性が成り立つ。
\end{proposition}

\begin{proof}
素イデアル $\mathfrak q \subset B$ で $V$ に含まれないものは $\mathfrak m_A$ 上にある。この場合
$A \to B_\mathfrak q$ は平坦局所準同型であり、したがって $\text{depth}(B_\mathfrak q) \geq 2$ である
（Algebra, 補題 \ref{algebra-lemma-apply-grothendieck}）。ゆえに補題
\ref{pione-lemma-quasi-compact-dense-open-connected-at-infinity-Noetherian} と Local Cohomology, 補題
\ref{local-cohomology-lemma-depth-2-connected-punctured-spectrum} により、この関手は充満忠実である。

\medskip\noindent
$A^\wedge$ および $B^\wedge$ を、それぞれ $A$ および $B$ の極大イデアルに関する完備化とする。
この命題の仮定は $A^\wedge \to B^\wedge$ に対して成り立つ。More on Algebra, 補題
\ref{more-algebra-lemma-completion-dimension},
\ref{more-algebra-lemma-completion-depth}, および
\ref{more-algebra-lemma-formally-smooth-completion} を参照せよ。補題 \ref{pione-lemma-extend-S2} の構成の一意性と
平坦基底変換との両立性により、$A^\wedge \to B^\wedge$ および $V$ の逆像について本質的全射性を
証明すれば十分である（細部を省略する。$V$ が穿孔スペクトルである場合については補題
\ref{pione-lemma-purity-and-completion} と比較せよ）。More on Algebra, 命題
\ref{more-algebra-proposition-fs-regular} により、これは $A \to B$ が正則であると仮定してよいことを意味する。

\medskip\noindent
$W \to V$ を有限エタール射とする。Popescu の定理（Smoothing Ring Maps, 定理
\ref{smoothing-theorem-popescu}）により、$B = \colim B_i$ を滑らかな $A$-代数のフィルター付き余極限として
書ける。ある $i$ と開部分集合 $V_i \subset \Spec(B_i)$ で、その逆像が $V$ となるものを選べる
（Limits, 補題 \ref{limits-lemma-descend-opens}）。$i$ を大きくした後、有限エタール射
$W_i \to V_i$ で、その $V$ への基底変換が $W \to V$ となるものが存在すると仮定してよい。
Limits, 補題 \ref{limits-lemma-descend-finite-presentation},
\ref{limits-lemma-descend-finite-finite-presentation}, および
\ref{limits-lemma-descend-etale} を参照せよ。$V_i$ の補集合は $\Spec(B_i) \to \Spec(A)$ の閉ファイバーに
含まれると仮定してよい。これは $V$ について成り立つからである（$V_i$ をこのように選ぶか、上の補題を
用いて十分大きい $i$ について成り立つことを示せばよい）。$\eta$ を $\Spec(B) \to \Spec(A)$ の
閉ファイバーの生成点とする。$\eta \in V$ なので、$\eta$ の像は $V_i$ に属する。したがって $V_i$ を
$\Spec(B)$ の閉点の像のアフィン開近傍で置き換えた後、$\Spec(B_i) \to \Spec(A)$ の閉ファイバーは既約で、
その生成点は $V_i$ に含まれると仮定してよい（細部を省略する。体上滑らかなスキームが既約スキームの
非交和であることを用いる）。ここで補題 \ref{pione-lemma-purity-smooth-over-depth2} を適用すると、
$W_i \to V_i$ は有限エタール射 $\Spec(C_i) \to \Spec(B_i)$ に延長される。これを $\Spec(B)$ へ
引き戻すと、$W$ は所望のとおり関手 $\textit{F\'Et}_Y \to \textit{F\'Et}_V$ の本質的像に属する。
\end{proof}





\section{基本群に対する Lefschetz}
\label{pione-section-global-lefschetz}

\noindent
もちろん、この種の結果は局所的な場合についてすでにいくつも証明した。本節では射影的な場合を論じる。

\begin{proposition}
\label{pione-proposition-lefschetz-fully-faithful}
$k$ を体とし、$X$ を $k$ 上固有なスキームとする。$\mathcal{L}$ を豊富な可逆
$\mathcal{O}_X$-加群とする。$s \in \Gamma(X, \mathcal{L})$ とし、$Y = Z(s)$ を $s$ の零点スキームとする。
任意の $x \in X \setminus Y$ に対して
$$
\text{depth}(\mathcal{O}_{X, x}) + \dim(\overline{\{x\}}) > 1
$$
が成り立つと仮定する。このとき制限関手 $\textit{F\'Et}_X \to \textit{F\'Et}_Y$ は充満忠実である。
実際、$V \subset X$ を $Y$ を含む任意の開部分スキームとすると、制限関手
$\textit{F\'Et}_V \to \textit{F\'Et}_Y$ は充満忠実である。
\end{proposition}

\begin{proof}
第一の主張は補題 \ref{pione-lemma-restriction-fully-faithful-special} および
Algebraic and Formal Geometry, 命題 \ref{algebraization-proposition-lefschetz} の形式的帰結である。
第二の主張は補題 \ref{pione-lemma-restriction-fully-faithful-special} および
Algebraic and Formal Geometry, 補題 \ref{algebraization-lemma-lefschetz-addendum} から従う。
\end{proof}

\begin{proposition}
\label{pione-proposition-lefschetz-equivalence-general}
$k$ を体とし、$X$ を $k$ 上固有なスキームとする。$\mathcal{L}$ を豊富な可逆
$\mathcal{O}_X$-加群とする。$s \in \Gamma(X, \mathcal{L})$ とし、$Y = Z(s)$ を $s$ の零点スキームとする。
$\mathcal{V}$ を、$X$ の開部分スキームで $Y$ を含むもの全体の集合に包含関係の逆順を入れたものとする。
任意の $x \in X \setminus Y$ に対して
$$
\text{depth}(\mathcal{O}_{X, x}) + \dim(\overline{\{x\}}) > 2
$$
が成り立つと仮定する。このとき制限関手
$$
\colim_\mathcal{V} \textit{F\'Et}_V \to \textit{F\'Et}_Y
$$
は同値である。
\end{proposition}

\begin{proof}
これは補題 \ref{pione-lemma-restriction-equivalence-general} および
Algebraic and Formal Geometry, 命題 \ref{algebraization-proposition-lefschetz-equivalence} の形式的帰結である。
\end{proof}

\begin{proposition}
\label{pione-proposition-lefschetz-equivalence}
$k$ を体とし、$X$ を $k$ 上固有なスキームとする。$\mathcal{L}$ を豊富な可逆
$\mathcal{O}_X$-加群とする。$s \in \Gamma(X, \mathcal{L})$ とし、$Y = Z(s)$ を $s$ の零点スキームとする。
任意の $x \in X \setminus Y$ に対して
$$
\text{depth}(\mathcal{O}_{X, x}) + \dim(\overline{\{x\}}) > 2
$$
が成り立ち、かつ閉点 $x \in X \setminus Y$ に対して $\mathcal{O}_{X, x}$ では純性が成り立つと仮定する。
このとき制限関手
$\textit{F\'Et}_X \to \textit{F\'Et}_Y$
は同値である。$X$、同値なことだが $Y$、が連結ならば
$$
\pi_1(Y, \overline{y}) \to \pi_1(X, \overline{y})
$$
は任意の幾何点 $\overline{y}$ が $Y$ 上にあるとき同型である。
\end{proposition}

\begin{proof}
命題 \ref{pione-proposition-lefschetz-fully-faithful} により充満忠実性が成り立つ。命題
\ref{pione-proposition-lefschetz-equivalence-general} により、$\textit{F\'Et}_Y$ の任意の対象はファイバー積
$U \times_V Y$ と同型である。ここで、ある有限エタール射 $U \to V$ があり、$V \subset X$ は $Y$ を含む
開部分スキームである。補集合 $T = X \setminus V$ は\footnote{実際、$T$ は $k$ 上固有であり
（$X$ で閉だからである）、かつアフィンである（アフィン・スキーム $X \setminus Y$ で閉だからである。
Morphisms, 補題 \ref{morphisms-lemma-proper-ample-delete-affine} を参照）。したがって $k$ 上有限である
（Morphisms, 補題 \ref{morphisms-lemma-finite-proper}）。ゆえに $T$ は閉点の有限集合である。}
$X \setminus Y$ の閉点の有限集合である。$T = \{x_1, \ldots, x_n\}$ とする。仮定により、有限エタール射
$V'_i \to \Spec(\mathcal{O}_{X, x_i})$ であって、$U \to V$ と
$V \times_X \Spec(\mathcal{O}_{X, x_i})$ 上で一致するものを取れる。Limits, 補題
\ref{limits-lemma-glueing-near-closed-point} を $n$ 回適用すると、$U \to V$ は所望の有限エタール射
$U' \to X$ に延長される。最後の主張については補題 \ref{pione-lemma-what-equivalence-gives} を参照せよ。
\end{proof}










\section{分岐軌跡の純性}
\label{pione-section-purity-ramification}

\noindent
本節では、生成的有限射に対する分岐軌跡の純性の類似を論じる。この結果は Gabber によるものらしい。
その特別な場合が van der Waerden の純性定理であり、正規多様体から滑らかな多様体への双有理射が
同型でない軌跡について述べる。

\begin{lemma}
\label{pione-lemma-characterize-rational-singularity}
$A$ を次元 $2$ のネーター正規局所整域とする。$A$ は Nagata であり、双対化加群 $\omega_A$ をもち、
特異点解消 $f : X \to \Spec(A)$ をもつと仮定する。$\omega_X$ を Resolution of Surfaces, 注意
\ref{resolve-remark-dualizing-setup} におけるものとする。ある有効 Cartier 因子 E について
$\omega_X \cong \mathcal{O}_X(E)$ が成り立ち、この $E \subset X$ の台が例外ファイバーに含まれるならば、
$A$ は有理特異点を定める。
$f$ が極小解消ならば $E = 0$ である。
\end{lemma}

\begin{proof}
跡写像 $Rf_*\omega_X \to \omega_A$ が存在する。Duality for Schemes, 節
\ref{duality-section-trace} を参照せよ。Grauert--Riemenschneider により
（Resolution of Surfaces, 命題 \ref{resolve-proposition-Grauert-Riemenschneider}）、
$R^1f_*\omega_X = 0$ である。したがって跡写像は写像 $f_*\omega_X \to \omega_A$ である。そこで
$$
\mathcal{O}_{\Spec(A)} = f_*\mathcal{O}_X \to f_*\omega_X \to \omega_A
$$
を考えられる。最初の写像は、この補題の主張で存在を仮定した写像
$\mathcal{O}_X \to \mathcal{O}_X(E) = \omega_X$ から来る。合成は Divisors, 補題
\ref{divisors-lemma-check-isomorphism-via-depth-and-ass} により同型である。実際、$A$ の穿孔スペクトル上で
同型であり（補題の仮定と、$f$ が穿孔スペクトル上で同型であることによる）、$A$ と $\omega_A$ は
$A$-加群として深さ $2$ をもつ（Algebra, 補題 \ref{algebra-lemma-criterion-normal} および
Dualizing Complexes, 補題 \ref{dualizing-lemma-depth-dualizing-module} による）。したがって
$f_*\omega_X \to \omega_A$ は全射であり、それゆえ同型である。ゆえに $Rf_*\omega_X = \omega_A$ であり、
双対性により $Rf_*\mathcal{O}_X = \mathcal{O}_{\Spec(A)}$ を得る。従って
$H^1(X, \mathcal{O}_X) = 0$ であり、これは $A$ が有理特異点を定めることを意味する
（Resolution of Surfaces, 節 \ref{resolve-section-bounded} の議論、特に補題
\ref{resolve-lemma-regular-rational} および \ref{resolve-lemma-exact-sequence} を参照）。$f$ が極小ならば
$E = 0$ である。なぜなら、写像 $f^*\omega_A \to \omega_X$ は Resolution of Surfaces, 補題
\ref{resolve-lemma-dualizing-blow-up-rational} の繰り返し適用により全射であり、さらに上で見たように
$\omega_A \cong A$ だからである。
\end{proof}

\begin{lemma}
\label{pione-lemma-key-purity-ramification}
$f : X \to \Spec(A)$ を有限型射とする。点 $x \in X$ をとる。次を仮定する。
\begin{enumerate}
\item $A$ は優秀正則局所環である。
\item $\mathcal{O}_{X, x}$ は次元 $2$ の正規環である。
\item $f$ は $\overline{\{x\}}$ の外でエタールである。
\end{enumerate}
このとき $f$ は $x$ でエタールである。
\end{lemma}

\begin{proof}
まず $X$ を $x$ のアフィン開近傍で置き換える。$\mathcal{O}_{X, x}$ は優秀局所環である
（More on Algebra, 補題 \ref{more-algebra-lemma-finite-type-over-excellent}）。
したがって、特異点の極小解消 $W \to \Spec(\mathcal{O}_{X, x})$ を選べる。
Resolution of Surfaces, 定理 \ref{resolve-theorem-resolve} を参照せよ。必要なら $X$ を $x$ の
アフィン開近傍で置き換えることにより、固有射 $b : X' \to X$ で
$X' \times_X \Spec(\mathcal{O}_{X, x}) = W$ を満たすものを見いだせる。
Limits, 補題 \ref{limits-lemma-glueing-near-closed-point} を参照せよ。さらに $X$ を縮小して、
$X'$ は正則であると仮定してよい。実際、$W$ は正則で $X'$ は優秀であり、優秀環のスペクトルの
正則軌跡は開である。また $W \to \Spec(\mathcal{O}_{X, x})$ は射影的である
（正規化された吹き上げの列だから）ので、$X$ を縮小した後 $b$ は射影的であると仮定してよい
（詳細は省略する）。$U = X \setminus \overline{\{x\}}$ とおく。
$W \to \Spec(\mathcal{O}_{X, x})$ は穿孔スペクトル上で同型であるから、
$b : X' \to X$ は $U$ 上で同型であると仮定してよい。したがって $U$ を $X'$ の開部分スキームとも
みなすことにする。$f' = f \circ b : X' \to \Spec(A)$ とおく。

\medskip\noindent
$A$ は正則なので、$\mathcal{O}_Y$ は $Y$ の双対化複体である。したがって
$f^!\mathcal{O}_Y$ は $X$ 上の双対化複体である
（Duality for Schemes, 補題 \ref{duality-lemma-shriek-dualizing}）。
Duality for Schemes, 補題 \ref{duality-lemma-dualizing-module-CM-scheme} により
$X$ の Cohen--Macaulay 軌跡は開である（これは優秀性を用いて証明することもできる）。
$\mathcal{O}_{X, x}$ は Cohen--Macaulay なので、$X$ を縮小した後、$X$ は Cohen--Macaulay
であると仮定してよい。エタール射は局所完全交叉であることに注意する。したがって
Duality for Schemes, 補題 \ref{duality-lemma-fundamental-class-almost-lci} を $r = 0$ として
適用でき、写像
$$
\mathcal{O}_X \longrightarrow \omega_{X/Y} = H^0(f^!\mathcal{O}_Y)
$$
を得る。これは $X \setminus \overline{\{x\}}$ 上で同型である。
Duality for Schemes, 補題 \ref{duality-lemma-shriek-over-CM} により $\omega_{X/Y}$ は $(S_2)$
であるから、Divisors, 補題 \ref{divisors-lemma-check-isomorphism-via-depth-and-ass} により
この写像は同型である。これで既に $X$、特に $\mathcal{O}_{X, x}$ が Gorenstein であることが分かる。

\medskip\noindent
$\omega_{X'/Y} = H^0((f')^!\mathcal{O}_Y)$ とおく。上とまったく同様に論じると、
$(f')^!\mathcal{O}_Y = \omega_{X'/Y}[0]$ は $X'$ の双対化複体である。$X'$ は正則なので、
射 $X' \to Y$ は局所完全交叉射である。More on Morphisms, 補題
\ref{more-morphisms-lemma-morphism-regular-schemes-is-lci} を参照せよ。
Duality for Schemes, 補題 \ref{duality-lemma-fundamental-class-lci} により写像
$$
\mathcal{O}_{X'} \longrightarrow \omega_{X'/Y}
$$
が存在し、これは $U$ 上で同型である。従って、ある有効 Cartier 因子について
$\omega_{X'/Y} = \mathcal{O}_{X'}(E)$ であり、この $E \subset X'$ は $U$ と交わらない。

\medskip\noindent
$\omega_{X/Y} = \mathcal{O}_Y$ なので、
$\omega_{X'/Y} = b^! f^!\mathcal{O}_Y = b^!\mathcal{O}_X$ である。
$W \to \Spec(\mathcal{O}_{X, x})$ に戻ると、$\omega_W = \mathcal{O}_W(E|_W)$ が分かる。
補題 \ref{pione-lemma-characterize-rational-singularity} により $E|_W = 0$ を得る。従って、既に用いた
Duality for Schemes, 補題 \ref{duality-lemma-fundamental-class-lci} により
$f' : X' \to Y$ はエタールである。$f'$ がエタールならば $b$ は同型でなければならないので、
これで直ちに証明が完了する。
\end{proof}

\begin{lemma}[分岐軌跡の純性]
\label{pione-lemma-purity-ramification}
\begin{reference}
複素空間に対するこの結果は \cite{Fischer} の 170 ページにある。一般の場合は
Gabber に帰せられている \cite[Theorem 2.4]{Zong} である。
\end{reference}
$f : X \to Y$ を局所ネーター・スキームの射とする。$x \in X$ とし、$y = f(x)$ とおく。次を仮定する。
\begin{enumerate}
\item $\mathcal{O}_{X, x}$ は次元 $\geq 1$ の正規環である。
\item $\mathcal{O}_{Y, y}$ は正則である。
\item $f$ は局所有限型である。
\item 特殊化 $x' \leadsto x$ が $\dim(\mathcal{O}_{X, x'}) = 1$ を満たすならば、
$f$ は $x'$ でエタールである。
\end{enumerate}
このとき $f$ は $x$ でエタールである。
\end{lemma}

\begin{proof}
この補題を $d = \dim(\mathcal{O}_{X, x})$ に関する帰納法で証明する。

\medskip\noindent
$d = 1$ の場合は、仮定により射 $f$ が $x$ でエタールなので自明である。$d \geq 2$ と仮定する。

\medskip\noindent
この補題の結論に影響を与えずに $\Spec(\mathcal{O}_{Y, y}) \to Y$ による基底変換を行える。
Morphisms, 補題 \ref{morphisms-lemma-set-points-where-fibres-etale} を参照せよ。従って
$Y = \Spec(A)$ と仮定してよい。ここで $A$ は正則局所環であり、$y$ は極大イデアル
$\mathfrak m$ に対応し、これは $A$ の極大イデアルである。

\medskip\noindent
特殊化 $x' \leadsto x$ で $x' \not = x$ を満たすものをとる。このとき $\mathcal{O}_{X, x'}$ は
$\mathcal{O}_{X, x}$ の局所化として正規である。$x'$ が $X$ の生成点でなければ
$1 \leq \dim(\mathcal{O}_{X, x'}) < d$ であり、帰納法の仮定により $f$ は $x'$ でエタールである。
従って、$f$ は $x$ に特殊化するすべての点でエタールであると仮定してよい。$f$ がエタールである
点の集合は（定義により）$X$ で開なので、$X$ を $x$ の開近傍で置き換えた後、
$f$ は $\overline{\{x\}}$ の外でエタールであると仮定してよい。特に、$f$ は閉点
$y \in Y = \Spec(A)$ 上にある点を除いてエタールである。

\medskip\noindent
$X' = X \times_{\Spec(A)} \Spec(A^\wedge)$ とおく。$x' \in X'$ を $x$ 上にある唯一の点とする。
上の議論により、$X'$ は閉ファイバーの外で $\Spec(A^\wedge)$ 上エタールであり、従って $X'$ は
閉ファイバーの外で正規である。$X$ は正規なので、Resolution of Surfaces, 補題
\ref{resolve-lemma-normalization-completion} により $X'$ は正規である。そこで
$X' \to \Spec(A^\wedge)$ が $x'$ でエタールであることを示せば、前述の Morphisms, 補題
\ref{morphisms-lemma-set-points-where-fibres-etale} により $f$ は $x$ でエタールとなる。
従って $A$ は正則完備局所環であると仮定してよく、実際そう仮定する。

\medskip\noindent
$d = 2$ の場合は補題 \ref{pione-lemma-key-purity-ramification} から従う。

\medskip\noindent
$d > 2$ と仮定する。$t \in \mathfrak m$、$t \not \in \mathfrak m^2$ をとる。
$Y_0 = \Spec(A/tA)$ および $X_0 = X \times_Y Y_0$ とおく。このとき $X_0 \to Y_0$ は
閉点上のファイバーの外でエタールである。$d > 2$ なので
$\dim(\mathcal{O}_{X_0, x}) = d - 1$ は $\geq 2$ である。正規化 $X_0' \to X_0$ は全射かつ有限である
（完備局所環上で作業しており、そのような環は Nagata だからである）。点
$x' \in X_0'$ で $x$ に写るものをとる。帰納法の仮定により、射 $X'_0 \to Y$ は $x'$ でエタールである。包含
$\kappa(y) \subset \kappa(x) \subset \kappa(x')$ から、$\kappa(x)$ は $\kappa(y)$ 上有限である。
従って $x$ は $X \to Y$ の $y$ 上のファイバーの閉点である。しかし $x$ はこのファイバーの
生成点でもあるので、$f$ は $x$ で準有限であると結論でき、分岐軌跡の純性の場合に帰着する。
補題 \ref{pione-lemma-purity} を参照せよ。
\end{proof}




\section{分岐軌跡の補集合のアフィン性}
\label{pione-section-stronger-purity}

\noindent
$f : X \to Y$ をネーター・スキームの有限型射とし、$X$ は正規、$Y$ は正則であるとする。
$V \subset X$ を $f$ がエタールである最大の開部分スキームとする。
\cite[Chapter IV, Section 21.12]{EGA} の議論は、$V \to X$ がアフィン射であろうことを示唆している。
$V \to X$ がアフィンならば、分岐軌跡の純性（補題 \ref{pione-lemma-purity-ramification}）が
Divisors, 補題 \ref{divisors-lemma-complement-affine-open-immersion} を用いて従うことに注意せよ。
従って $V \to X$ のアフィン性は、分岐軌跡に対する純性の「強い」形である。
本節では、$V \to X$ が、$X$ と $Y$ が等標数かつ優秀であるときアフィンであることを証明する。
定理 \ref{pione-theorem-global} を参照せよ。$X$ と $Y$ が混標数の場合にも（依然として優秀ならば）
この結果は真であろうと予想するのが妥当に思われる。

\begin{lemma}
\label{pione-lemma-structure-cohomology}
$(A, \mathfrak m)$ を体を含む正則局所環とする。
$f : V \to \Spec(A)$ をエタールかつ準コンパクトとする。
$\mathfrak m \not \in f(V)$ であり、$g : V \to \Spec(A) \setminus \{\mathfrak m\}$ は
アフィンであると仮定する。このとき $H^i(V, \mathcal{O}_V)$ は、$i > 0$ に対して
$A$ の剰余体の入射包のコピーの直和に同型である。
\end{lemma}

\begin{proof}
穿孔スペクトルを $U = \Spec(A) \setminus \{\mathfrak m\}$ と書く。このとき $g : V \to U$ は
アフィンである。Cohomology of Schemes, 補題
\ref{coherent-lemma-relative-affine-cohomology} により
$H^i(V, \mathcal{O}_V) = H^i(U, g_*\mathcal{O}_V)$ である。
Schemes, 補題 \ref{schemes-lemma-push-forward-quasi-coherent} により、
$\mathcal{O}_U$-加群 $g_*\mathcal{O}_V$ は準連接である。任意の準連接
$\mathcal{O}_U$-加群 $\mathcal{F}$ に対して、コホモロジー $H^i(U, \mathcal{F})$ は $i > 0$ ならば
$\mathfrak m$-冪ねじれである。例えば Local Cohomology, 補題
\ref{local-cohomology-lemma-local-cohomology} を参照せよ。特に $A$-加群
$H^i(V, \mathcal{O}_V)$ は $i > 0$ ならば $\mathfrak m$-冪ねじれである。
任意の平坦環準同型 $A \to A'$ に対し、平坦基底変換 Cohomology of Schemes, 補題
\ref{coherent-lemma-flat-base-change-cohomology} により
$H^i(V, \mathcal{O}_V) \otimes_A A' = H^i(V', \mathcal{O}_{V'})$ である。ここで
$V' = V \times_{\Spec(A)} \Spec(A')$ である。$A'$ を $A$ の完備化とすると
（More on Algebra, 節 \ref{more-algebra-section-permanence-completion} により平坦である）、次を得る。
$$
H^i(V, \mathcal{O}_V) = H^i(V, \mathcal{O}_V) \otimes_A A' =
H^i(V', \mathcal{O}_{V'}),\quad\text{ただし } i > 0
$$
最初の等号は、いま証明したねじれ性と More on Algebra, 補題
\ref{more-algebra-lemma-neighbourhood-equivalence} による。さらに、剰余体 $k$ の入射包は
$A$ と $A'$ に対して同じである。Dualizing Complexes, 補題
\ref{dualizing-lemma-compare} を参照せよ。このようにして
$A = k[[x_1, \ldots, x_d]]$ の場合に帰着する。Algebra, 節
\ref{algebra-section-cohen-structure-theorem} を参照せよ。

\medskip\noindent
$k$ の標数が $p > 0$ であると仮定する。$F : A \to A$、$a \mapsto a^p$ は平坦であり
（Local Cohomology, 補題 \ref{local-cohomology-lemma-frobenius-flat-regular}）、
また \'Etale Morphisms, 補題 \ref{etale-lemma-relative-frobenius-etale} により
$V \times_{\Spec(A), \Spec(F)} \Spec(A) \cong V$ は $\Spec(A)$ 上のスキームとして同型なので、
上の議論から
$H^i(V, \mathcal{O}_V) \otimes_{A, F} A \cong H^i(V, \mathcal{O}_V)$ を得る。従って
Local Cohomology, 補題
\ref{local-cohomology-lemma-structure-torsion-Frobenius-regular} により結果が従う。

\medskip\noindent
$k$ の標数が $0$ であると仮定する。Local Cohomology, 補題
\ref{local-cohomology-lemma-etale-derivation} により、加法的作用素 $D_j$、
$j = 1, \ldots, d$ が $H^i(V, \mathcal{O}_V)$ 上にあり、
$\partial_j = \partial/\partial x_j$ に関する Leibniz 則を満たす。従って Local Cohomology, 補題
\ref{local-cohomology-lemma-structure-torsion-D-module-regular} により結果を得る。
\end{proof}

\begin{lemma}
\label{pione-lemma-conclude}
補題 \ref{pione-lemma-structure-cohomology} の状況で、
$H^i(V, \mathcal{O}_V) = 0$ が $i \geq \dim(A) - 1$ に対して成り立つと仮定する。
このとき $V$ はアフィンである。
\end{lemma}

\begin{proof}
$k = A/\mathfrak m$ とおく。$V \times_{\Spec(A)} \Spec(k) = \emptyset$ なので、
コホモロジーと基底変換により
$$
R\Gamma(V, \mathcal{O}_V) \otimes_A^\mathbf{L} k = 0
$$
である。Derived Categories of Schemes, 補題 \ref{perfect-lemma-compare-base-change} を参照せよ。
従ってスペクトル系列（More on Algebra, 例 \ref{more-algebra-example-tor}）
$$
E_2^{p, q} = \text{Tor}_{-p}(k, H^q(V, \mathcal{O}_V)),\quad
d_2^{p, q} : E_2^{p, q} \to E_2^{p + 2, q - 1}
$$
と、零に収束する $d_r^{p, q} : E_r^{p, q} \to E_r^{p + r, q - r + 1}$ がある。
補題 \ref{pione-lemma-structure-cohomology}、Dualizing Complexes, 補題
\ref{dualizing-lemma-tor-injective-hull}、および
$H^i(V, \mathcal{O}_V) = 0$ が $i \geq \dim(A) - 1$ に対して成り立つという仮定から、
位置 $(p, q) = (0, 0)$ に入る、またはそこから出る非零の微分はないと結論できる。従って
$H^0(V, \mathcal{O}_V) \otimes_A k = 0$ である。これは、もし
$\mathfrak m = (x_1, \ldots, x_d)$ ならば、開被覆
$V = \bigcup V \times_{\Spec(A)} \Spec(A_{x_i})$ をアフィン開部分スキーム
$V \times_{\Spec(A)} \Spec(A_{x_i})$ によってもち
（$V$ は $A$ の穿孔スペクトル上アフィンだから）、
$x_1, \ldots, x_d$ が $\Gamma(V, \mathcal{O}_V)$ の単位イデアルを生成することを意味する。
Properties, 補題 \ref{properties-lemma-characterize-affine} により、これは $V$ がアフィンであることを
意味する。
\end{proof}

\begin{theorem}
\label{pione-theorem-global}
$Y$ を体上の優秀正則スキームとする。$f : X \to Y$ をスキームの有限型射とし、$X$ は正規であるとする。
$V \subset X$ を $f$ がエタールである最大の開部分スキームとする。このとき包含射 $V \to X$ は
アフィンである。
\end{theorem}

\begin{proof}
$x \in X$ を、その像が $y \in Y$ である点とする。$V \cap W$ がアフィンとなるような
アフィン開近傍 $W$ を $x$ に対して見いだせば十分である。$\Spec(\mathcal{O}_{X, x})$ はスキーム $W$ の
極限なので、これは
$$
V_x = V \times_X \Spec(\mathcal{O}_{X, x})
$$
がアフィンであることと同値である（Limits, 補題 \ref{limits-lemma-limit-affine}）。
従って、射
$X \times_Y \Spec(\mathcal{O}_{Y, y}) \to \Spec(\mathcal{O}_{Y, y})$ に対して定理が成り立てば、
元の定理も成り立つ。特に $Y$ は有限次元の正則スキームであると仮定でき、次元
$d = \dim(Y)$ に関する帰納法を行える。同じ議論をもう一度組み合わせると、
$Y$ は閉点 $y$ をもつ局所スキームであり、
$V \cap (X \setminus f^{-1}(\{y\}) \to X \setminus f^{-1}(\{y\})$ は
アフィンであると仮定してよい。

\medskip\noindent
$x \in X$ を $y$ 上の点とする。$x \in V$ ならば証明すべきことはない。
$f^{-1}(\{y\}) \cap V$ は閉点の有限集合であることに注意せよ
（エタール射のファイバーは離散的だからである）。従って $X$ を $x$ のアフィン開近傍で置き換えた後、
$y \not \in f(V)$ と仮定してよい。$V$ がアフィンであることを証明しなければならない。

\medskip\noindent
$e(V)$ を、$i$ のうち $H^i(V, \mathcal{O}_V) \not = 0$ を満たすものの最大値とする。$X$ はアフィンなので、
整数 $e(V)$ は数 $e(V_x)$（ここで $x \in X \setminus V$）の最大値である。Local Cohomology, 補題
\ref{local-cohomology-lemma-cd-local} および Local Cohomology, 補題
\ref{local-cohomology-lemma-cd} におけるコホモロジー次元の特徴づけを参照せよ。
Local Cohomology, 補題 \ref{local-cohomology-lemma-cd-dimension} により
$e(V_x) \leq \dim(\mathcal{O}_{X, x}) - 1$ である。
$\dim(\mathcal{O}_{X, x}) \geq 2$ ならば、分岐軌跡の純性
（補題 \ref{pione-lemma-purity-ramification}）により、$V_x$ は $\mathcal{O}_{X, x}$ の穿孔スペクトルより
真に小さい。$\mathcal{O}_{X, x}$ は正規かつ優秀なので、Hartshorne--Lichtenbaum 消滅
（Local Cohomology, 補題 \ref{local-cohomology-lemma-affine-complement}）により
$e(V_x) \leq \dim(\mathcal{O}_{X, x}) - 2$ が従う。他方、$X \to Y$ は有限型であり、
$V \subset X$ は稠密なので（必要なら $X$ を $V$ の閉包で置き換える）、
次元公式（Morphisms, 補題 \ref{morphisms-lemma-dimension-formula}）により
$\dim(\mathcal{O}_{X, x}) \leq d$ である。従って $e(V) \leq \max(0, d -  2)$ である。
$V$ は、$d \geq 2$ ならば補題 \ref{pione-lemma-conclude} によりアフィンである。
$d = 1$ または $d = 0$ ならば、$\mathcal{O}_{Y, y}$ の穿孔スペクトルはアフィンであり、
従って $V$ はアフィンである。
\end{proof}




\section{滑らかかつ固有の場合の特殊化写像}
\label{pione-section-specialization-smooth-proper}

\noindent
本節では次の結果を論じる。$f : X \to S$ をスキームの固有滑らかな射とする。
$s \leadsto s'$ を $S$ の点の特殊化とする。このとき、節 \ref{pione-section-specialization-map} の特殊化写像
$$
sp : \pi_1(X_{\overline{s}}) \longrightarrow \pi_1(X_{\overline{s}'})
$$
は全射であり、次が成り立つ。
\begin{enumerate}
\item $\kappa(s')$ の標数が零ならば、これは同型である。
\item $\kappa(s')$ の標数が $p > 0$ ならば、これは $p$ と素な最大商の間の同型を誘導する。
\end{enumerate}

\begin{lemma}
\label{pione-lemma-specialization-map-surjective}
$f : X \to S$ を幾何学的に連結なファイバーをもつ平坦固有射とする。
$s' \leadsto s$ を特殊化とする。$X_s$ が幾何学的に被約ならば、特殊化写像
$sp : \pi_1(X_{\overline{s}'}) \to \pi_1(X_{\overline{s}})$ は全射である。
\end{lemma}

\begin{proof}
$X_s$ は幾何学的に被約なので、必要なら $S$ を縮小して、すべてのファイバーが幾何学的に被約であると
仮定してよい。More on Morphisms, 補題
\ref{more-morphisms-lemma-geometrically-reduced-open} を参照せよ。
$\mathcal{O}_{S, s} \to A \to \kappa(\overline{s}')$ を、特殊化写像の構成に現れるものとする。
節 \ref{pione-section-specialization-map} を参照せよ。従って、
$$
\pi_1(X_{\overline{s}'}) \to \pi_1(X_A)
$$
が全射であることを示せば十分である。これは命題
\ref{pione-proposition-first-homotopy-sequence} と $\pi_1(\Spec(A)) = \{1\}$ から従う。
\end{proof}

\begin{proposition}
\label{pione-proposition-specialization-map-isomorphism}
$f : X \to S$ を幾何学的に連結なファイバーをもつ滑らかかつ固有な射とする。
$s' \leadsto s$ を特殊化とする。$\kappa(s)$ の標数が零ならば、特殊化写像
$$
sp : \pi_1(X_{\overline{s}'}) \to \pi_1(X_{\overline{s}})
$$
は同型である。
\end{proposition}

\begin{proof}
この写像は補題 \ref{pione-lemma-specialization-map-surjective} により全射である。従って単射性を示せばよい。

\medskip\noindent
$S$ はアフィンであると仮定してよい。このとき $S$ は $\mathbf{Z}$ 上有限型のアフィン・スキームの
余フィルター付き極限である。従って、$X \to S$ は $X_0 \to S_0$ の基底変換であると仮定できる。
ここで $S_0$ は有限型 $\mathbf{Z}$-代数のスペクトルであり、$X_0 \to S_0$ は滑らかかつ固有である。
Limits, 補題 \ref{limits-lemma-descend-finite-presentation}、
\ref{limits-lemma-descend-smooth}、および \ref{limits-lemma-eventually-proper} を参照せよ。補題
\ref{pione-lemma-specialization-map-base-change} により、基底がネーターである場合に帰着する。

\medskip\noindent
補題 \ref{pione-lemma-specialization-map-discrete-valuation-ring} を適用すると、基底 $S$ が強 Hensel 離散付値環
$A$ のスペクトルであり、$A$ 上の特殊化写像を考える場合に帰着する。$K$ を $A$ の分数体とする。
$\overline{K}$ を、幾何学的生成点 $\overline{\eta}$（$\Spec(A)$ の点）に対応する代数閉包として選ぶ。
$\overline{K}/L/K$ が有限分離であるとき、$B \subset L$ を $A$ の整閉包とし、この整閉包は $L$ 内でとる。
これは More on Algebra, 注意 \ref{more-algebra-remark-finite-separable-extension} により離散付値環である。

\medskip\noindent
$X \to \Spec(A)$ を前段落のものとする。特殊化写像の単射性を示すには、任意の有限エタール被覆
$V$ で $X_{\overline{\eta}}$ の被覆であるものが、有限エタール被覆 $Y \to X$ の基底変換であることを
証明すれば十分である。実際、このとき補題 \ref{pione-lemma-functoriality-galois-injective} により
$\pi_1(X_{\overline{\eta}}) \to \pi_1(X) = \pi_1(X_s)$ は単射である。

\medskip\noindent
$V$ が与えられたとき、まず補題 \ref{pione-lemma-perfection} により $V$ を
$V' \to X_{K^{sep}}$ に降下し、次に補題 \ref{pione-lemma-limit} により $V'' \to X_L$ に降下できる。
$Z \to X_B$ を $X_B$ の正規化とし、これは $V''$ の中でとる。$Z$ は正規であり、
$Z_L = V''$ が $X_L$ 上のスキームとして成り立つことに注意せよ。従って $Z \to X_B$ は
生成ファイバー上で有限エタールである。問題は $Z \to X_B$ が至る所エタールかどうか分からないことである。
$X \to \Spec(A)$ は幾何学的に連結な滑らかなファイバーをもつので、特殊ファイバー $X_s$ は
幾何学的既約である。従って $X_B \to \Spec(B)$ の特殊ファイバーは既約である。その生成点を
$\xi_B$ とする。$\xi_1, \ldots, \xi_r$ を $Z$ の点で $\xi_B$ に写るものとする。いま最初の問題
（そして唯一の問題であることが分かる）は、離散付値環の拡大
$$
\mathcal{O}_{X_B, \xi_B} \subset \mathcal{O}_{Z, \xi_i}
$$
が分岐している可能性である。この拡大の分岐指数を $e_i$ とする。$\kappa(s)$ の標数は零なので、
この分岐は馴分岐であることに注意せよ。

\medskip\noindent
分岐を取り除くため、さらなる有限分離拡大 $K^{sep}/L'/L/K$ で、その分岐指数 $e$、すなわち
誘導される拡大 $B'/B$ の分岐指数が $e_i$ で割り切れるものを選ぶ。$Z'$ を $Z$ の
$\Spec(B') \to \Spec(B)$ に関する正規化基底変換とする。More on Morphisms, 節
\ref{more-morphisms-section-reduced-fibre-theorem} の議論を参照せよ。
$\xi_{i, j}$ を $Z'$ の点で、$\xi_{B'}$ および $\xi_i$（$Z$ の点）に写るものとする。このとき局所環
$$
\mathcal{O}_{Z', \xi_{i, j}}
$$
は、$\mathcal{O}_{Z, \xi_i}$ の $L' \otimes_L F_i$ における整閉包の局所化である。ここで $F_i$ は
$\mathcal{O}_{Z, \xi_i}$ の分数体である。詳細は省略する。従って Abhyankar の補題
（More on Algebra, 補題 \ref{more-algebra-lemma-abhyankar}）により、
$$
\mathcal{O}_{X_{B'}, \xi_{B'}} \subset \mathcal{O}_{Z', \xi_{i, j}}
$$
は不分岐である。従って射 $Z' \to X_{B'}$ は余次元 $1$ の外でエタールである。
ゆえに分岐軌跡の純性（補題 \ref{pione-lemma-purity}）により、$Z' \to X_{B'}$ は有限エタールである。

\medskip\noindent
しかし、$A \to B'$ が誘導する剰余体拡大は自明である（$A$ の剰余体は標数零の分離閉体なので
代数閉だからである）。従って補題 \ref{pione-lemma-finite-etale-on-proper-over-henselian} を二度適用すると、
$Z'$ は有限エタール被覆 $Y \to X$ の基底変換であると分かる
（まず $Y$ を $A$ 上で得て、次に $B$ への引き戻しが $Z'$ と同型であることを示す）。
これで証明が完了する。
\end{proof}

\noindent
$G$ を副有限群とし、$p$ を素数とする。{\it $p$ と素な最大商}は、定義により
$$
G' = \lim_{U \subset G\text{ は開かつ正規で、その指数は }p\text{ と素}} G/U
$$
である。$X$ が連結スキームで $p$ が与えられたとき、$p$ と素な最大商を
$\pi_1(X)$ に対して $\pi'_1(X)$ と書く。

\begin{theorem}
\label{pione-theorem-specialization-map-isomorphism-prime-to-p}
$f : X \to S$ を幾何学的に連結なファイバーをもつ滑らかかつ固有な射とする。
$s' \leadsto s$ を特殊化とする。$\kappa(s)$ の標数が $p$ ならば、特殊化写像
$$
sp : \pi_1(X_{\overline{s}'}) \to \pi_1(X_{\overline{s}})
$$
は全射であり、同型
$$
\pi'_1(X_{\overline{s}'}) \cong \pi'_1(X_{\overline{s}})
$$
を p と素な最大商の間に誘導する。
\end{theorem}

\begin{proof}
これは命題 \ref{pione-proposition-specialization-map-isomorphism} とまったく同様に証明されるが、
次の相違がある。
\begin{enumerate}
\item $X/A$ が与えられたとき、関手
$\textit{F\'Et}_X \to \textit{F\'Et}_{X_{\overline{\eta}}}$ が本質的全射であることまでは示さない。
Galois 群の位数が $p$ と素である Galois 対象が本質像に属することだけを示す。同じ議論により
$\pi'_1(X_{\overline{s}'}) \to \pi'_1(X_{\overline{s}})$ の単射性を結論するには、これで十分である。
\item 拡大
$\mathcal{O}_{X_B, \xi_B} \subset \mathcal{O}_{Z, \xi_i}$
は馴分岐である。なぜなら、対応する分数体の拡大は Galois で、その群の位数は $p$ と素だからである。
More on Algebra, 補題 \ref{more-algebra-lemma-galois-conclusion} を参照せよ。
\item 拡大 $\kappa_B/\kappa_A$ はもはや必ずしも自明ではないが、純非分離である。従って射
$X_{\kappa_B} \to X_{\kappa_A}$ は普遍的同相であり、命題
\ref{pione-proposition-universal-homeomorphism} により基本群の同型を誘導する。
\end{enumerate}
\end{proof}













\section{馴分岐}
\label{pione-section-tame}

\noindent
$X \to Y$ を $\mathbf{Z}$ 上有限型なスキームの有限エタール射とする。$f$ が $\infty$ で
馴分岐することの定義には多くの方法がある。論文 \cite{Kerz-Schmidt} では、これらの概念が
どの程度一致するかが論じられている。

\medskip\noindent
本節では、無限遠での馴分岐性という概念に先立つ、別のより初等的な問題を論じる。
\cite{Grothendieck-Murre} の（やや異なる）議論と比較されたい。次が与えられていると仮定する。
\begin{enumerate}
\item 局所ネーター・スキーム $X$、
\item 稠密開集合 $U \subset X$、
\item 有限エタール射 $f : Y \to U$。
\end{enumerate}
さらに、すべての素因子 $Z \subset X$ で $Z \cap U = \emptyset$ を満たすものに対して、
局所環 $\mathcal{O}_{X, \xi}$、すなわち $X$ の、生成点 $\xi$（$Z$ のもの）における局所環は
離散付値環であるとする。
$K_\xi$ を $\mathcal{O}_{X, \xi}$ の分数体とおけば、スキームのデカルト図式
$$
\xymatrix{
\Spec(K_\xi) \ar[r] \ar[d] & U \ar[d] \\
\Spec(\mathcal{O}_{X, \xi}) \ar[r] & X
}
$$
を得る。特に、$Y \times_U \Spec(K_\xi)$ は有限分離代数 $L_\xi/K_\xi$ のスペクトルである。
このとき、{\it $Y$ は $X$ 上余次元 $1$ で不分岐である}、それぞれ
{\it $Y$ は $X$ 上余次元 $1$ で馴分岐する}という。ただし、$L_\xi/K_\xi$ が
$\mathcal{O}_{X, \xi}$ に関して、上のようなすべての $(Z, \xi)$ についてそれぞれ不分岐、馴分岐である
ことを意味する。More on Algebra, 定義
\ref{more-algebra-definition-types-of-extensions} を参照せよ。より正確には、$L_\xi$ を
$K_\xi$ の有限分離体拡大の積に分解し、その各々が $\mathcal{O}_{X, \xi}$ に関してそれぞれ不分岐、
馴分岐であることを要求する。

\begin{lemma}
\label{pione-lemma-pullback-tame-codim1}
$X' \to X$ を局所ネーター・スキームの射とし、$U \subset X$ を稠密開集合とする。次を仮定する。
\begin{enumerate}
\item $U' = f^{-1}(U)$ は $X'$ の稠密開集合である。
\item すべての素因子 $Z \subset X$ で $Z \cap U = \emptyset$ を満たすものに対して、
局所環 $\mathcal{O}_{X, \xi}$、すなわち $X$ の、生成点 $\xi$（$Z$ のもの）における局所環は
離散付値環である。
\item すべての素因子 $Z' \subset X'$ で $Z' \cap U' = \emptyset$ を満たすものに対して、
局所環 $\mathcal{O}_{X', \xi'}$、すなわち $X'$ の、生成点 $\xi'$（$Z'$ のもの）における局所環は
離散付値環である。
\item $\xi' \in X'$ が (3) のような点ならば、$\xi = f(\xi')$ は (2) のような点である。
\end{enumerate}
このとき $f : Y \to U$ が有限エタールで、$Y$ が $X$ 上余次元 $1$ でそれぞれ不分岐、馴分岐ならば、
$Y' = Y \times_X X' \to U'$ は有限エタールであり、$Y'$ は $X'$ 上余次元 $1$ でそれぞれ不分岐、
馴分岐である。
\end{lemma}

\begin{proof}
この補題で唯一興味深い事実は、More on Algebra, 補題
\ref{more-algebra-lemma-tame-goes-up} に与えられている可換代数の結果である。
\end{proof}

\noindent
上で導入した用語を用いると、先に得た純性の結果を次の簡潔な形で言い換えられる。

\begin{lemma}
\label{pione-lemma-purity-one-divisor-general}
$X$ を局所ネーター・スキームとする。$U \subset X$ を稠密開集合とする。
$Y \to U$ を有限エタール射とする。次を仮定する。
\begin{enumerate}
\item $Y$ は $X$ 上余次元 $1$ で不分岐である。
\item $\mathcal{O}_{X, x}$ はすべての $x \in X \setminus U$ に対して正則である。
\end{enumerate}
このとき、有限エタール射 $Y' \to X$ で、その $U$ への制限が $Y$ であるものが存在する。
\end{lemma}

\begin{proof}
$\xi \in X \setminus U$ を、$X \setminus U$ の余次元 $1$ の既約成分の生成点とする。このとき
$\mathcal{O}_{X, \xi}$ は離散付値環である。上の議論と同様に
$Y \times_U \Spec(K_\xi) = \Spec(L_\xi)$ と書く。$B_\xi$ を
$\mathcal{O}_{X, \xi}$ の $L_\xi$ における整閉包とする。$Y$ が $X$ 上余次元 $1$ で不分岐である
という仮定は、$\mathcal{O}_{X, \xi} \to B_\xi$ が有限エタールであることを意味する。
従って有限エタール射 $Y_\xi \to \Spec(\mathcal{O}_{X, \xi})$ と同型
$$
Y \times_U \Spec(K_\xi) \cong
Y_\xi \times_{\Spec(\mathcal{O}_{X, \xi})} \Spec(K_\xi)
$$
を得る。これは $\Spec(K_\xi)$ 上の同型である。Limits, 補題
\ref{limits-lemma-glueing-near-point} により、開部分スキーム
$U \subset U' \subset X$ で $\xi$ を含むものと、有限表示射 $Y' \to U'$ で、その $U$ への制限が $Y$ を復元し、
$\Spec(\mathcal{O}_{X, \xi})$ への制限が $Y_\xi$ を復元するものを得る。最後に、Limits, 補題
\ref{limits-lemma-glueing-near-point-properties} により、射 $Y' \to U'$ は、必要なら $U'$ をより小さい
開集合に縮小した後、有限エタールである。$X \setminus U$ の余次元 $1$ の他の生成点についてこの議論を
繰り返すと、有限エタール射 $Y' \to U'$ で $Y \to U$ を延長するものをもつと仮定してよい。
ここで開部分スキーム $U' \subset X$ は $U$ と、すべての余次元 $1$ の点
（$X \setminus U$ のもの）を含む。
最後に補題 \ref{pione-lemma-extend-pure} を $Y' \to U'$ に適用する。実際、すべての局所環
$\mathcal{O}_{X, x}$（ここで $x \in X \setminus U'$）は正則であり、
$\dim(\mathcal{O}_{X, x}) \geq 2$ を満たす。従って補題 \ref{pione-lemma-local-purity} により
$\mathcal{O}_{X, x}$ に対して純性が成り立つ。
\end{proof}

\begin{lemma}
\label{pione-lemma-purity-one-divisor}
$X$ を局所ネーター・スキームとする。$D \subset X$ を有効 Cartier 因子とし、$D$ は正則スキームであるとする。
$Y \to X \setminus D$ を有限エタール射とする。$Y$ が $X$ 上余次元 $1$ で不分岐ならば、
有限エタール射 $Y' \to X$ で、その $X \setminus D$ への制限が $Y$ であるものが存在する。
\end{lemma}

\begin{proof}
これは補題 \ref{pione-lemma-purity-one-divisor-general} の特別な場合である。まず $D$ は $X$ の
至る所稠密でない部分集合であり（Divisors, 節
\ref{divisors-section-effective-Cartier-divisors} の議論を参照）、従って
$X \setminus D$ は $X$ で稠密である。次に、環 $\mathcal{O}_{X, x}$ はすべての $x \in D$ に対して
正則局所環である。これは Algebra, 補題 \ref{algebra-lemma-regular-mod-x} と
$\mathcal{O}_{D, x}$ が正則であるという仮定による。
\end{proof}

\begin{example}[標準的な馴分岐射]
\label{pione-example-tamely-ramified}
$A$ をネーター環とする。$f \in A$ を非零因子とし、$A/fA$ は被約であるとする。すると
$A_\mathfrak p$ は素元 $f$ をもつ離散付値環である。これは、任意の極小素イデアル
$\mathfrak p$ で $f$ 上のものに対して成り立つ。$e \geq 1$ を $A$ で可逆な整数とする。
$$
C = A[x]/(x^e - f)
$$
とおく。このとき $\Spec(C) \to \Spec(A)$ は有限局所自由射であり、$A_f$ のスペクトル上でエタールである。
有限エタール射
$$
\Spec(C_f) \longrightarrow \Spec(A_f)
$$
は $\Spec(A)$ 上余次元 $1$ で馴分岐する。この馴分岐性は More on Algebra, 補題
\ref{more-algebra-lemma-characterize-tame} における馴分岐拡大の特徴づけから直ちに従う。
\end{example}

\noindent
正則因子に対する Abhyankar の補題の一つの形を述べる。

\begin{lemma}[正則因子に対する Abhyankar の補題]
\label{pione-lemma-abhyankar-one-divisor}
$X$ を局所ネーター・スキームとする。$D \subset X$ を有効 Cartier 因子とし、$D$ は正則スキームであるとする。
$Y \to X \setminus D$ を有限エタール射とする。$Y$ が $X$ 上余次元 $1$ で馴分岐するならば、
$X$ 上エタール局所的に、射 $Y \to X$ は例
\ref{pione-example-tamely-ramified} に記述した標準的な馴分岐射の有限非交和として与えられる。
\end{lemma}

\begin{proof}
すべての $x \in X$ に対し、エタール近傍 $(U, u) \to (X, x)$ で
$U = \Spec(A)$ を満たし、基底変換 $Y \times_X U \to U$ が例
\ref{pione-example-tamely-ramified} のような標準的馴分岐射の有限非交和となるものを見いだす。
$x \in D$ と仮定する。$x \not \in D$ の場合は \'Etale Morphisms, 補題
\ref{etale-lemma-finite-etale-etale-local} と、例 \ref{pione-example-tamely-ramified} で
$e = 1$ および $f = 1$ ととることから従う。

\medskip\noindent
この段落では、$X$ が強 Hensel 局所環のスペクトルで、$x$ が閉点である場合に帰着する。実際、$X$ を
縮小して、$X = \Spec(A)$ であり、$D \subset X$ が非零因子 $f \in A$ により与えられると仮定してよい。
このとき $Y$ は $\Spec(A_f)$ の有限エタール被覆としてアフィンである。$Y = \Spec(B)$ と書く。
$A^{sh}$ を $\mathcal{O}_{X, x}$ の強 Hensel 化とする。$A \to A^{sh}$ は平坦で $x \in D$ なので、
$f$ は $A^{sh}$ の極大イデアルの非零因子に写る。また
$A^{sh}/fA^{sh} = (A/fA)^{sh}$ は（Algebra, 補題
\ref{algebra-lemma-quotient-strict-henselization}）、$\mathcal{O}_{D, x}$ の強 Hensel 化として
正則局所環である。More on Algebra, 補題
\ref{more-algebra-lemma-henselization-regular} を参照せよ。補題
\ref{pione-lemma-pullback-tame-codim1} により、$Y$ の $\Spec(A^{sh})$ への基底変換は余次元 $1$ で
馴分岐する。この $Y$ の $\Spec(A^{sh})$ への基底変換について主張を証明したと仮定する。
$A^{sh}$ は強 Hensel なので、任意のエタール近傍は切断をもち、同型
$$
A^{sh} \otimes_A B \cong
\prod\nolimits_{i \in I} A^{sh}_f[x]/(x^{e_i} - f)
$$
を得る。ここで $I$ は有限集合で、各 $e_i$ は $A^{sh}$ で可逆な整数である。
$A^{sh} = \colim A_\lambda$ をエタール $A$-代数 $A_\lambda$ のフィルター付き余極限として書く。
上の同型は同型
$$
A_\lambda \otimes_A B \cong \prod\nolimits_{i \in I}
(A_\lambda)_f[x]/(x^{e_i} - f)
$$
へ降下する。ただし $\lambda$ は十分大きくとる。例えば Algebra, 補題
\ref{algebra-lemma-colimit-category-fp-algebras} を参照せよ。$\lambda$ をもう少し大きくすれば、
$e_i$ は $A_\lambda$ で可逆であると仮定してよい。このとき
$\Spec(A_\lambda) \to \Spec(A)$ は求める $x$ のエタール近傍である。

\medskip\noindent
$X = \Spec(A)$ と仮定する。ここで $A$ は強 Hensel 局所環であり、$x \in X$ は極大イデアル
$\mathfrak m$（$A$ のもの）に対応する。$f \in \mathfrak m$ を $D$ を切り出す非零因子とする。
このとき $A/fA$ は正則であり、Algebra, 補題 \ref{algebra-lemma-regular-mod-x} により $A$ も正則である。
正則局所環のいくつかの性質、例えばそれが正規整域であることを用いる。Algebra, 補題
\ref{algebra-lemma-regular-domain} および \ref{algebra-lemma-regular-normal} を参照せよ。特に $A$ は整域である。
上と同様に $Y = \Spec(B)$ はアフィンであり、$A_f \to B$ は有限エタールである。従って $B$ も正規である
（Algebra, 補題 \ref{algebra-lemma-normal-goes-up}）。よって Algebra, 補題
\ref{algebra-lemma-characterize-reduced-ring-normal} により、$Y$ は正規整域のスペクトルの有限非交和である。
従って $B$ も整域であると仮定してよい。$A$ と $A/fA$ は整域なので、$f$ は $A$ の素元であり、
高さ $1$ の素イデアル $\mathfrak p = fA$ を生成する。$A_\mathfrak p$ は $f$ を素元とする
離散付値環であることに注意せよ。

\medskip\noindent
$K$ を $A$ の分数体、$L$ を $B$ の分数体とする。馴分岐の仮定は、$L$ が
$A_\mathfrak p$ に関して馴分岐することを意味する。$e \geq 1$ を、More on Algebra, 注意
\ref{more-algebra-remark-finite-separable-extension} にあるように、$L$ の $A_\mathfrak p$ 上の
分岐指数の積とする。このとき $e$ は $\kappa(\mathfrak p)$ で可逆であるが、この時点では $e$ が
$A/fA$ または $A$ で可逆かどうかは分からない。

\medskip\noindent
有限自由 $A$-代数
$$
A' = A[x]/(x^e - f)
$$
を考える。これは強 Hensel 局所環の局所有限拡大であり、従って強 Hensel である。
$f' = x$ は $A'$ の非零因子であり、$A'/f'A' \cong A/fA$ は正則環であることに注意せよ。
従って前と同様に $A'$ は正則であり、まして整域である。
$B' = B \otimes_A A' = B \otimes_{A_f} A'_{f'}$.
とおく。Abhyankar の補題（More on Algebra, 補題 \ref{more-algebra-lemma-abhyankar}）により、
$\Spec(B')$ は $\Spec(A')$ 上余次元 $1$ で不分岐である。実際、補題
\ref{pione-lemma-pullback-tame-codim1} により $\Spec(B')$ は少なくともなお $\Spec(A')$ 上余次元 $1$ で
馴分岐する。しかし Abhyankar の補題により、分岐指数はすべて $1$ になっている。補題
\ref{pione-lemma-purity-one-divisor} により、$\Spec(B') \to \Spec(A'_{f'})$ は有限エタール射
$\Spec(C) \to \Spec(A')$ に延長される。しかし $A'$ は強 Hensel なので、$\Spec(C)$ は
$\Spec(A')$ のコピーの有限非交和である。結論として、少なくとも一つの射
$\Spec(A'_{f'}) \to \Spec(B')$ が存在する。従って包含
$B \subset A'_{f'}$ が $A_f$-代数として存在する。

\medskip\noindent
従って包含 $K \subset L \subset K[f^{1/e}]$ をもつ。More on Algebra, 補題
\ref{more-algebra-lemma-pull-root-uniformizer} により、$L$ は $K[f^{1/n}]$ に等しい。ここで $n$ は
$e$ のある約数である。すると $B$ の正規性から $B \cong A_f[y]/(y^n - f)$ が従う。詳細は省略する。
しかし環準同型 $A_f \to A_f[y]/(y^n - f)$ の分岐軌跡は $ny^{n - 1}$ で切り出されるので、
$n$ は $A_f$ の単元でなければならない。$f$ は素元なので、これは $n = u f^s$ であることを意味する。
ここで $u$ は $A$ のある単元で、$s \in \mathbf{Z}$ である。$n$ は
$\kappa(\mathfrak p)$ の単元に写るので $s = 0$ となり、証明が完了する。
\end{proof}

\begin{lemma}
\label{pione-lemma-extend-tame-covering-normal}
補題 \ref{pione-lemma-abhyankar-one-divisor} の状況で、$X$ の $Y$ における正規化は有限局所自由射
$\pi : Y' \to X$ であり、次を満たす。
\begin{enumerate}
\item $Y'$ の $X \setminus D$ への制限は $Y$ と同型である。
\item $D' = \pi^{-1}(D)_{red}$ は $Y'$ 上の有効 Cartier 因子である。
\item $D'$ は正則スキームである。
\end{enumerate}
さらに、$X$ 上エタール局所的に、射 $Y' \to X$ は射
$$
\Spec(A[x]/(x^e - f)) \to \Spec(A)
$$
の有限非交和である。ここで $A$ はネーター環、$f \in A$ は非零因子で $A/fA$ は正則、
$e \geq 1$ は $A$ で可逆である。
\end{lemma}

\begin{proof}
これは補題 \ref{pione-lemma-abhyankar-one-divisor} への補足にすぎず、実際、その補題の証明を読めば
この補題が真であることはほとんど直ちに従う。しかし補題
\ref{pione-lemma-abhyankar-one-divisor} の結果から直接導くこともできる。実際、$X$ の $Y$ における
正規化をとる操作はエタール基底変換と可換である。More on Morphisms, 補題
\ref{more-morphisms-lemma-normalization-smooth-localization} を参照せよ。従って、$Y' \to X$ の局所構造に
関する主張は $X$ 上エタール局所的に証明してよい。ゆえに補題
\ref{pione-lemma-abhyankar-one-divisor} により、$X = \Spec(A)$ で $A$ はネーター環であり、非零因子
$f\in A$ をもち、$A/fA$ は正則で、$Y$ は環 $A_f[x]/(x^e - f)$ のスペクトルの有限非交和で、
$e$ は $A$ で可逆であると仮定してよい。$A$ の整閉包（$A_f[x]/(x^e - f)$ におけるもの）が
$A' = A[x]/(x^e - f)$ に等しいことの確認は省略する
（これを見るには、$A'$ の局所化で、$(f)$ 上にある素イデアルにおけるものが正則であることを示せばよい）。
詳細は省略する。
\end{proof}

\begin{lemma}
\label{pione-lemma-tame-covering-split}
補題 \ref{pione-lemma-abhyankar-one-divisor} の状況で、$Y' \to X$ を補題
\ref{pione-lemma-extend-tame-covering-normal} のものとする。$R$ を分数体 $K$ をもつ離散付値環とする。
$$
t : \Spec(R) \to X
$$
を射とし、スキーム論的逆像 $t^{-1}D$ は $\Spec(R)$ の被約閉点であるとする。
\begin{enumerate}
\item $t|_{\Spec(K)}$ が $Y$ の点へ持ち上がるならば、持ち上げ
$t' : \Spec(R) \to Y'$ で、$Y' \to X$ が $t'(\Spec(R))$ に沿ってエタールとなるものを得る。
\item $\Spec(K) \times_X Y$ が $\Spec(K)$ のコピーの非交和と同型ならば、$Y' \to X$ は
$t(\Spec(R))$ のある開近傍上で有限エタールである。
\end{enumerate}
\end{lemma}

\begin{proof}
有限射 $Y' \to X$ に固有性の付値判定法を適用すると、$\Spec(K)$-値点で $Y$ のもののうち、
$t|_{\Spec(K)}$ と $X$ への写像として一致するものは、射 $t' : \Spec(R) \to Y'$ に一意に持ち上がる。
従って主張 (1) は意味をもつ。

\medskip\noindent
エタール近傍 $(U, u) \to (X, t(\mathfrak m_R))$ で、$U = \Spec(A)$ かつ
$Y' \times_X U \to U$ が補題 \ref{pione-lemma-extend-tame-covering-normal} のような記述を、ある
$f \in A$ に対してもつものを選ぶ。このとき $\Spec(R) \times_X U \to \Spec(R)$ はエタールかつ全射である。
$R'$ を、$\Spec(R) \times_X U$ の局所環で $\Spec(R)$ の閉点上にあるものとすると、$R'$ は
離散付値環で、$R \subset R'$ は離散付値環の不分岐拡大である
（More on Algebra, 補題 \ref{more-algebra-lemma-Dedekind-etale-extension}）。$t$ に関する仮定は、
写像 $A \to R'$ で
$$
\Spec(R') \to \Spec(R) \times_X U \to U
$$
に対応するものが $f$ を素元 $\pi \in R'$ に写すことを意味する。いま
$$
Y' \times_X U =
\coprod\nolimits_{i \in I} \Spec(A[x]/(x^{e_i} - f))
$$
がある $e_i \geq 1$ に対して成り立つと仮定する。このとき
$$
\Spec(R') \times_U (Y' \times_X U) =
\coprod\nolimits_{i \in I} \Spec(R'[x]/(x^{e_i} - \pi))
$$
を得る。環 $R'[x]/(x^{e_i} - f)$ は離散付値環である
（More on Algebra, 補題 \ref{more-algebra-lemma-pull-root-uniformizer}）。従って $R'$ の分数体への
写像をもつには $e_i = 1$ でなければならない。

\medskip\noindent
(1) の証明。この場合、写像 $t' : \Spec(R) \to Y'$ は基底変換により対応する写像
$t'' : \Spec(R') \to Y' \times_X U$ を定め、上の議論により、これは $i \in I$ かつ $e_i = 1$ に
対応する成分へ写らなければならない。従って明らかに $Y' \times_X U \to U$ は $t''$ の像に沿って
エタールである。エタール性はエタール基底変換後に判定できる性質なので、これで (1) が証明された。

\medskip\noindent
(2) の証明。この場合、仮定は、$e_i = 1$ がすべての $i \in I$ に対して成り立つことを意味する。
従って $Y' \times_X U \to U$ は有限エタールであり、前と同様に結論を得る。
\end{proof}

\begin{lemma}
\label{pione-lemma-extend-covering}
$S$ を生成点 $\eta$ をもつ整正規ネーター・スキームとする。
$f : X \to S$ を幾何学的に連結なファイバーをもつ滑らかな射とする。
$\sigma : S \to X$ を $f$ の切断とする。$Z \to X_\eta$ を有限エタール Galois 被覆
（節 \ref{pione-section-finite-etale-under-galois}）とし、その群 $G$ の位数は $S$ 上可逆で、
$Z$ は $\kappa(\eta)$-有理点で $\sigma(\eta)$ に写るものをもつとする。このとき
有限エタール Galois 被覆 $Y \to X$ で群 $G$ をもち、その $X_\eta$ への制限が $Z$ であるものが存在する。
\end{lemma}

\begin{proof}
まず $S = \Spec(R)$ は離散付値環 $R$ のスペクトルで、閉点を $s \in S$ とすると仮定する。
このとき $X_s$ は $X$ の有効 Cartier 因子であり、体上滑らかなスキームとして $X_s$ は正則である。
さらに生成ファイバー $X_\eta$ は開部分スキーム $X \setminus X_s$ である。
More on Algebra, 補題 \ref{more-algebra-lemma-galois-conclusion} と $G$ に関する仮定から、$Z$ は
$X$ 上余次元 $1$ で馴分岐する。$Z' \to X$ を補題
\ref{pione-lemma-extend-tame-covering-normal} のものとする。$G$ の $Z$ への作用は $G$ の $Z'$ への作用に
延長されることに注意せよ。補題 \ref{pione-lemma-tame-covering-split} により、$Z' \to X$ は
$\sigma(y)$ のある開近傍上で有限エタールである。$X_s$ は既約なので、これは $Z \to X_\eta$ が
$X$ 上余次元 $1$ で不分岐であることを意味する。すると補題
\ref{pione-lemma-purity-one-divisor} により有限エタール射 $Y \to X$ で、その $X_\eta$ への制限が
$Z$ であるものを得る。もちろん $Y \cong Z'$ であり（詳細は省略する。ヒント：エタール局所的に計算せよ）、
従って $Y$ は群 $G$ をもつ Galois 被覆である。

\medskip\noindent
一般の場合。$U \subset S$ を次の性質をもつ最大の開部分スキームとする。すなわち、
有限エタール Galois 被覆 $Y \to X \times_S U$ で群 $G$ をもち、その $X_\eta$ への制限が $Z$ と同型であるものが
存在する。矛盾を導くため $U \not = S$ と仮定する。$s \in S \setminus U$ を
$S \setminus U$ の既約成分の生成点とする。このとき $U_s$ を $U$ の
$\Spec(\mathcal{O}_{S, s})$ における逆像とすると、これは $\mathcal{O}_{S, s}$ の穿孔スペクトルである。写像
$Y \times_S U_s \to X \times_S U_s$ は、有限エタール Galois 被覆
$Y'_s \to X \times_S \Spec(\mathcal{O}_{S, s})$ で群 $G$ をもつものの制限であると主張する。

\medskip\noindent
まず、この主張から所望の矛盾が生じることを示す。Limits, 補題
\ref{limits-lemma-glueing-near-point} により、開部分スキーム
$U \subset U' \subset S$ で $s$ を含むものと有限表示射 $Y'' \to U'$ で、その $U$ への制限が $Y' \to U$ を復元し、
$\Spec(\mathcal{O}_{S, s})$ への制限が $Y'_s$ を復元するものを得る。さらに、この補題が与える圏の同値により、
$U'$ を縮小した後、射 $Y'' \to U' \times_S X$ と、$G$ の $Y''$ への作用で $U' \times_S X$ 上のものが
存在し、これらは $U$ および $\Spec(\mathcal{O}_{S, s})$ への基底変換後に与えられた射と作用に整合すると
仮定してよい。必要なら $U'$ をさらに縮小して、$Y'' \to U \times_S X$ は有限エタールであると
仮定してよい。Limits, 補題 \ref{limits-lemma-glueing-near-point-properties} を参照せよ。
これは、$S$ の真に大きい開集合で、その上へ $Y$ が群 $G$ をもつ有限エタール Galois 被覆として
延長されるものを見いだしたことを意味し、求めていた矛盾を与える。

\medskip\noindent
主張の証明。$S$ を $\Spec(\mathcal{O}_{S, s})$ で置き換えてよく、実際そうする。このとき
$S = \Spec(A)$ であり、$(A, \mathfrak m)$ は正規局所整域である。また $U \subset S$ は穿孔スペクトルで、
有限エタール Galois 被覆 $Y \to X \times_S U$ で群 $G$ をもつものがある。$\dim(A) = 1$ ならば、
証明の最初の段落により $Y$ の $X$ の Galois 被覆への延長を構成できる。従って
$\dim(A) \geq 2$ と仮定してよく、$\text{depth}(A) \geq 2$ である。これは $S$ が正規だからである。
Algebra, 補題 \ref{algebra-lemma-criterion-normal} を参照せよ。$X \to S$ は平坦なので、
$\text{depth}(\mathcal{O}_{X, x}) \geq 2$ と、すべての点 $x \in X$ で $s$ に写るものに対して結論できる。Algebra, 補題
\ref{algebra-lemma-apply-grothendieck} を参照せよ。
$$
Y' \longrightarrow X
$$
を、補題 \ref{pione-lemma-extend-S2} で $Y \to X \times_S U$ を用いて構成される有限射とする。
$G$ の標準的な作用を $Y$ 上に得ることに注意せよ。従って残るのは、$Y'$ が $X$ 上エタールであることを
示すことだけである。実際、例えば補題 \ref{pione-lemma-purity-smooth-over-depth2} により、
$Y' \to X$ がファイバー $X_s$ の（唯一の）生成点上でエタールであることを示せば十分である。
これは $\sigma(s)$ の（形式的）近傍における局所計算によって示す。

\medskip\noindent
アフィン開集合 $\Spec(B) \subset X$ で $\sigma(s)$ を含むものを選ぶ。このとき $A \to B$ は滑らかな環準同型で、
切断 $\sigma : B \to A$ をもつ。$I = \Ker(\sigma)$ と書き、$B^\wedge$ を $I$-進完備化（$B$ のもの）とする。
すると $B^\wedge \cong A[[x_1, \ldots, x_d]]$ である。ここである $d \geq 0$ が存在する。Algebra, 補題
\ref{algebra-lemma-section-smooth} を参照せよ。もちろん $B \to B^\wedge$ は平坦であり
（Algebra, 補題 \ref{algebra-lemma-completion-flat}）、$\Spec(B^\wedge) \to X$ の像は $X_s$ の生成点を含む。
$V \subset \Spec(B^\wedge)$ を $U$ の逆像とする。有限エタール射
$$
W = Y \times_{(X \times_S U)} V \longrightarrow V
$$
を考える。補題 \ref{pione-lemma-extend-S2} における $Y'$ の構成と平坦基底変換との整合性により、基底変換
$Y' \times_X \Spec(B^\wedge) \to \Spec(B^\wedge)$
は、$W \to V$ から $\Spec(B^\wedge)$ 上で補題 \ref{pione-lemma-extend-S2} の手続きによって構成される。
$V_0 = V \cap V(x_1, \ldots, x_d) \subset V$ および $W_0 = W \times_V V_0$ とおく。
これは正規整スキームであり、$\sigma(S)$ の中へ、射 $\Spec(B^\wedge) \to X$ により写り、実際
$\sigma(U)$ と同一視される。従って $W_0 \to V_0 = U$ は完全に分解する。実際、その生成ファイバーに
ついてこれは、$Z \to X_\eta$ が $\kappa(\eta)$-有理点で $\sigma(\eta)$ 上にあるものをもつという仮定から
成り立つ（そしてもちろん $G$-作用から、ファイバー全体 $Z_{\sigma(\eta)}$ はスキーム
$\eta = \Spec(\kappa(\eta))$ のコピーの非交和である）。最後に補題
\ref{pione-lemma-fully-faithful-power-series-over-depth2} により
$$
W_0 \times_U V \cong W
$$
である。これは $W$ が $V$ のコピーの非交和であることを示す。従って
$Y' \times_X \Spec(B^\wedge)$ は $\Spec(B^\wedge)$ のコピーの非交和であり、証明が完了する。
\end{proof}

\begin{lemma}
\label{pione-lemma-extend-covering-general}
$S$ を生成点 $\eta$ をもつ準コンパクトかつ準分離な整正規スキームとする。
$f : X \to S$ を準コンパクトかつ準分離で滑らかな射で、幾何学的に連結なファイバーをもつものとする。
$\sigma : S \to X$ を $f$ の切断とする。$Z \to X_\eta$ を有限エタール Galois 被覆
（節 \ref{pione-section-finite-etale-under-galois}）とし、その群 $G$ の位数は $S$ 上可逆で、
$Z$ は $\kappa(\eta)$-有理点で $\sigma(\eta)$ に写るものをもつとする。このとき有限エタール Galois 被覆
$Y \to X$ で群 $G$ をもち、その $X_\eta$ への制限が $Z$ であるものが存在する。
\end{lemma}

\begin{proof}
$S$ がネーターならば、これは補題 \ref{pione-lemma-extend-covering} の結果である。
一般の場合は標準的な極限論法によりこれから従う。
読者には証明を読み飛ばすことを強く勧める。

\medskip\noindent
$S = \lim S_i$ と、アフィン推移射をもつスキーム系の有向極限として書け、
$S_i$ は $\mathbf{Z}$ 上有限型である。Limits, 命題
\ref{limits-proposition-approximate} を参照せよ。
各 $i$ に対し、$S \to S'_i \to S_i$ を、$S_i$ を $S$ の中で正規化して得られる射とする。
Morphisms, 節 \ref{morphisms-section-normalization-X-in-Y} を参照せよ。
Algebra, 命題 \ref{algebra-proposition-ubiquity-nagata}、Morphisms, 補題
\ref{morphisms-lemma-nagata-normalization-finite} および
\ref{morphisms-lemma-normal-normalization} を合わせると、$S'_i$ は $\mathbf{Z}$ 上有限型で、
$S_i$ 上有限であり、さらに $S'_i$ は整正規スキームで $S \to S'_i$ が支配的であると結論できる。
Morphisms, 補題 \ref{morphisms-lemma-functoriality-normalization} により、推移射
$S'_{i'} \to S'_i$ で、推移射 $S_{i'} \to S_i$ および始域が $S$ である射と整合するものを得る。
$S = \lim S'_i$ であると主張する。この主張の証明は省略する（ヒント：$S_i$ の選んだアフィン開集合上の
アフィン開集合を、ある $i$ に対して考えれば、これは素朴な代数の問題に翻訳される）。
従って $S = \lim S_i$ と、正規整スキーム $S_i$ の系の有向極限として書け、推移射はアフィンで、
$S_i$ は $\mathbf{Z}$ 上有限型であるとしてよい。

\medskip\noindent
ある $i$ に対し有限表示の滑らかな射 $X_i \to S_i$ で、$S$ への基底変換が $X \to S$ であるものを
見いだせる。Limits, 補題 \ref{limits-lemma-descend-finite-presentation} および
\ref{limits-lemma-descend-smooth} を参照せよ。
$i$ を大きくした後、切断 $\sigma$ は切断 $\sigma_i : S_i \to X_i$ に持ち上がると仮定してよい
（Limits, 補題 \ref{limits-lemma-descend-finite-presentation} の圏同値による）。
$X_i$ を、More on Morphisms, 節 \ref{more-morphisms-section-connected-components} で調べた
その開部分スキーム $X_i^0$ で置き換えてよい。実際、$X \to X_i$ の像は明らかにそこへ写る
（開性については More on Morphisms, 補題
\ref{more-morphisms-lemma-connected-along-section-open} による）。
従って $X_i \to S_i$ のファイバーは幾何学的に連結であると仮定してよい。
$i$ を大きくした後、$|G|$ は $S_i$ 上可逆であると仮定してよい。
$\eta_i \in S_i$ を生成点とする。$X_\eta$ はスキーム $X_{i, \eta_i}$ の極限なので、
まったく同じ論法により、$Z \to X_\eta$ をある有限エタール Galois 被覆
$Z_i \to X_{i, \eta_i}$ へ降下できるよう、必要なら $i$ を大きくできる。補題 \ref{pione-lemma-limit} を参照せよ。
必要ならさらに一度 $i$ を大きくした後、$Z_i$ は $\kappa(\eta_i)$-有理点で
$\sigma_i(\eta_i)$ に写るものをもつと仮定してよい。
そこでネーターの場合の補題を適用し、$X$ へ引き戻せば結論を得る。
\end{proof}











\section{正標数における技巧}
\label{pione-section-achinger}

\noindent
Piotr Achinger の論文 \cite{Achinger} では、正標数のアフィン・スキームが常に
$K(\pi, 1)$ であることが示されている。この節では \cite{Achinger} のうち、より初等的な部分を説明する。
すなわち、正標数の体 $k$ に対し、$\mathbf{A}^n_k$ 上エタールなアフィン・スキームは、
実際には（別の射によって）$\mathbf{A}^n_k$ 上有限エタールであることを示す。また、正標数の
連結アフィン・スキームの閉埋め込みが、エタール基本群の間の単射を誘導することも示す。

\medskip\noindent
$k$ を標数 $p > 0$ の体とする。
全射
$$
k[x_1, \ldots, x_n] \longrightarrow A
$$
を有限型 $k$-代数の全射とし、その始域は $x_1, \ldots, x_n$ 上の多項式環とする。
核を $I \subset k[x_1, \ldots, x_n]$ と書けば、
$A = k[x_1, \ldots, x_n]/I$ である。$A$ が非零であるとは仮定しない
（言い換えれば、$A$ が零環で $I = k[x_1, \ldots, x_n]$ となる場合も許す）。最後に、
有限エタール環準同型 $\pi : A \to B$ が与えられていると仮定する。

\medskip\noindent
$k, n, k[x_1, \ldots, x_n] \to A, I, \pi : A \to B$ が与えられているとする。
$C$ を $k$-代数とする。可換図式
$$
\xymatrix{
& B \\
C \ar[r] & C/\varphi(I)C \ar[u]^\tau \\
k[x_1, \ldots, x_n] \ar[u]^\varphi \ar[r] &
A \ar[u] \ar@/_3em/[uu]_\pi
}
$$
を考える。ここで $\varphi$ はエタール $k$-代数準同型、$\tau$ は全射 $k$-代数準同型である。
$C, \varphi, \tau$ が与えられているとする。任意の $r \geq 0$ と $y_1, \ldots, y_r \in C$ で、
$C$ を $\Im(\varphi)$ 上の代数として生成するものに対し、
$s = s(r, y_1, \ldots, y_r) \in \{0, \ldots, r\}$ を、$y_i$ が $\Im(\varphi)$ 上整となる
$1 \leq i \leq s$ の最大の元とする。$NF(C, \varphi, \tau)$ を、
$r - s = r - s(r, y_1, \ldots, y_r)$ が、上のような $r$ および $y_1, \ldots, y_r$ の
すべての選択にわたって取る最小値として定める。
$NF(C, \varphi, \tau)$ が $0$ であることと $\varphi$ が有限であることは同値である。

\begin{lemma}
\label{pione-lemma-lower-invariant}
上の状況で $NF(C, \varphi, \tau) > 0$ ならば、エタール $k$-代数準同型 $\varphi'$ と全射
$k$-代数準同型 $\tau'$ で、可換図式
$$
\xymatrix{
& B \\
C \ar[r] & C/\varphi'(I)C \ar[u]_{\tau'} \\
k[x_1, \ldots, x_n] \ar[u]^{\varphi'} \ar[r] &
A \ar[u] \ar@/_3em/[uu]_\pi
}
$$
に入り、$NF(C, \varphi', \tau') < NF(C, \varphi, \tau)$ を満たすものが存在する。
\end{lemma}

\begin{proof}
$r \geq 0$ および $y_1, \ldots, y_r \in C$ で $C$ を $\Im(\varphi)$ 上生成するものを選び、
$0 \leq s \leq r$ を、$y_1, \ldots, y_s$ が $\Im(\varphi)$ 上整で、
$r - s = NF(C, \varphi, \tau) > 0$ となるように取る。
$B$ は $A$ 上有限なので、$y_{s + 1}$ の $B$ における像は $A$ 上のモニック多項式を満たす。
従って $d \geq 1$ および $f_1, \ldots, f_d \in k[x_1, \ldots, x_n]$ で、
$$
z = y_{s + 1}^d + \varphi(f_1) y_{s + 1}^{d - 1} + \ldots + \varphi(f_d) \in
J = \Ker(C \to C/\varphi(I)C \xrightarrow{\tau} B)
$$
となるものを見いだせる。$\varphi : k[x_1, \ldots, x_n] \to C$ はエタールなので、
非零かつ非定数の多項式 $g \in k[T_1, \ldots, T_{n + 1}]$ で、
$$
g(\varphi(x_1), \ldots, \varphi(x_n), z) = 0
\quad\text{（以下の環において）}\quad
C
$$
となるものを見いだせる。これを見るには、例えば
$C \otimes_{\varphi, k[x_1, \ldots, x_n]} k(x_1, \ldots, x_n)$ が
$k(x_1, \ldots, x_n)$ の有限次分離体拡大の有限直積であること
（Algebra, 補題 \ref{algebra-lemma-etale-over-field} を参照）を用いればよい。
従って $z$ は $k(x_1, \ldots, x_n)$ 上のモニック多項式を満たす。分母を払えば $g$ を得る。

\medskip\noindent
$g$ の存在と Algebra, 補題 \ref{algebra-lemma-helper-polynomial} により、整数
$e_1, e_2, \ldots, e_n \geq 1$ で、$z$ が部分環 $C'$（$C$ のもの）、すなわち
$t_1 = \varphi(x_1) + z^{pe_1}, \ldots, t_n = \varphi(x_n) + z^{pe_n}$ で生成されるものの上で
整となるようなものを得る。もちろん元 $\varphi(x_1), \ldots, \varphi(x_n)$ も $C'$ 上整であり、
元 $y_1, \ldots, y_s$ も同様である。最後に、$z$ の選び方により元 $y_{s + 1}$ も $C'$ 上整である。

\medskip\noindent
環準同型
$$
\varphi' : k[x_1, \ldots, x_n] \longrightarrow C, \quad
x_i \longmapsto t_i
$$
を考える。その像は $C'$ である。
$\text{d}(\varphi(x_i)) = \text{d}(t_i) = \text{d}(\varphi'(x_i))$
が $\Omega_{C/k}$ で成り立つので（ここで $k$ の標数が $p > 0$ であることを使う）、
$\varphi'$ がエタールであると結論できる。これは $\varphi$ がエタールだからである。Algebra, 補題
\ref{algebra-lemma-characterize-etale-over-polynomial-ring} を参照せよ。
$\varphi'(x_i) - \varphi(x_i) = t_i - \varphi(x_i) = z^{pe_i}$ は、核 $J$、すなわち写像
$C \to C/\varphi(I)C \to B$ の核に属することに注意せよ。これは $z$ を $J$ の元として選んだからである。
従って $f \in I$ に対して元
$$
\varphi'(f) =
f(t_1, \ldots, t_n) =
f(\varphi(x_1) + z^{pe_1}, \ldots, \varphi(x_n) + z^{pe_n}) =
\varphi(f) + \text{（次のイデアルの元）}(z)
$$
も $J$ に属する。言い換えれば $\varphi'(I)C \subset J$ であり、全射
$$
\tau' : C/\varphi'(I)C \longrightarrow C/J \cong B
$$
を $A$ 上エタールな代数の間に得る。最後に、代数 $C$ は元
$\varphi(x_1), \ldots, \varphi(x_n), y_1, \ldots, y_r$ によって
$C' = \Im(\varphi')$ 上生成され、
$\varphi(x_1), \ldots, \varphi(x_n), y_1, \ldots, y_{s + 1}$ は
$C' = \Im(\varphi')$ 上整である。従って
$NF(C, \varphi', \tau') < r - s = NF(C, \varphi, \tau)$ である。これで証明が完了する。
\end{proof}

\begin{lemma}
\label{pione-lemma-affine-etale-over-affine-space}
$k$ を標数 $p > 0$ の体とする。$X \to \mathbf{A}^n_k$ をエタール射とし、$X$ はアフィンとする。
このとき有限エタール射 $X \to \mathbf{A}^n_k$ が存在する。
\end{lemma}

\begin{proof}
$X = \Spec(C)$ と書く。$A = 0$ とおき、$I = k[x_1, \ldots, x_n]$ と書く。
仮定により、あるエタール $k$-代数準同型
$\varphi : k[x_1, \ldots, x_n] \to C$ が存在する。$\tau : C/\varphi(I)C \to 0$ を一意な全射と書く。
$\varphi$ および $\tau$ は $N(C, \varphi, \tau)$ が最小となるように選べる。補題
\ref{pione-lemma-lower-invariant} により $N(C, \varphi, \tau) = 0$ を得る。従って $\varphi$ は有限エタールである。
\end{proof}

\begin{lemma}
\label{pione-lemma-dominate-affine-space}
$k$ を標数 $p > 0$ の体とする。$Z \subset \mathbf{A}^n_k$ を閉部分スキームとする。
$Y \to Z$ を有限エタールとする。このとき有限エタール射 $f : U \to \mathbf{A}^n_k$ で、
開かつ閉な埋め込み $Y \to f^{-1}(Z)$ が $Z$ 上存在するようなものが存在する。
\end{lemma}

\begin{proof}
この問題を代数に翻訳しよう。
$\mathbf{A}^n_k = \Spec(k[x_1, \ldots, x_n])$.
と書く。このとき $Z = \Spec(A)$ であり、$A = k[x_1, \ldots, x_n]/I$、ここで
$I \subset k[x_1, \ldots, x_n]$ はあるイデアルである。
$Y = \Spec(B)$ と書けば、$Y \to Z$ は有限エタール $k$-代数準同型 $A \to B$ に対応する。

\medskip\noindent
Algebra, 補題 \ref{algebra-lemma-lift-etale} により、エタール環準同型
$$
\varphi : k[x_1, \ldots, x_n] \to C
$$
と全射 $A$-代数準同型 $\tau : C/\varphi(I)C \to B$ が存在する。
（$C, \varphi, \tau$ は $\tau$ が同型となるようにさえ選べるが、これは用いない。）
$\varphi$ および $\tau$ は $N(C, \varphi, \tau)$ が最小となるように選べる。補題
\ref{pione-lemma-lower-invariant} により $N(C, \varphi, \tau) = 0$ を得る。従って $\varphi$ は有限エタールである。

\medskip\noindent
$f : U = \Spec(C) \to \mathbf{A}^n_k$ を $\varphi$ に対応する有限エタール射とする。
射 $Y \to f^{-1}(Z) = \Spec(C/\varphi(I)C)$ は $\tau$ によって誘導され、$\tau$ が全射なので閉埋め込みであり、
Morphisms, 補題 \ref{morphisms-lemma-etale-permanence} によりエタール射なので開でもある。
これで証明が完了する。
\end{proof}

\noindent
以下が主要結果である。

\begin{proposition}
\label{pione-proposition-injective}
$p$ を素数とする。$i : Z \to X$ を $\mathbf{F}_p$ 上の連結アフィン・スキームの閉埋め込みとする。
任意の幾何点 $\overline{z}$（$Z$ のもの）に対し、写像
$$
\pi_1(Z, \overline{z}) \to \pi_1(X, \overline{z})
$$
は単射である。
\end{proposition}

\begin{proof}
$Y \to Z$ を有限エタール射とする。有限エタール射 $f : U \to X$ で、$Y$ が
$f^{-1}(Z)$ の開かつ閉な部分スキームと同型になるようなものを構成すれば十分である。
補題 \ref{pione-lemma-functoriality-galois-injective} を参照せよ。
$Y = \Spec(A)$ および $X = \Spec(R)$ と書けば、閉埋め込み $Y \to X$ は全射 $R \to A$ によって与えられる。
$A = \colim A_i$ と、その有限型 $\mathbf{F}_p$-部分代数のフィルター付き余極限として書ける。
補題 \ref{pione-lemma-limit} により、ある $i$ と有限エタール射 $Y_i \to Z_i = \Spec(A_i)$ で、
$Y = Z \times_{Z_i} Y_i$ となるものを見いだせる。

\medskip\noindent
全射 $\mathbf{F}_p[x_1, \ldots, x_n] \to A_i$ を選ぶ。これは閉埋め込み
$$
Z_i = \Spec(A_i)
\longrightarrow
X_i = \mathbf{A}^n_{\mathbf{F}_p} = \Spec(\mathbf{F}_p[x_1, \ldots, x_n])
$$
を定める。多項式環の普遍性と $R \to A$ が全射であることにより、可換図式
$$
\xymatrix{
\mathbf{F}_p[x_1, \ldots, x_n] \ar[r] \ar[d] &
A_i \ar[d] \\
R \ar[r] &
A
}
$$
を $\mathbf{F}_p$-代数の間に見いだせる。従って可換図式
$$
\xymatrix{
Y_i \ar[r] &
Z_i \ar[r] &
X_i \\
Y \ar[u] \ar[r] &
Z \ar[u] \ar[r] &
X \ar[u]
}
$$
を得、その右側の正方形はデカルトである。有限エタールな $f_i : U_i \to X_i$ で、$Y_i$ が
$f_i^{-1}(Z_i)$ の開かつ閉な部分スキームと同型となるものを見いだせれば、
$f : U \to X$、すなわち $f_i$ の $X \to X_i$ による基底変換が問題の解になることは明らかである。
従って補題 \ref{pione-lemma-dominate-affine-space} を
$Y_i \to Z_i \to X_i = \mathbf{A}^n_{\mathbf{F}_p}$ に適用すれば結論を得る。
\end{proof}
